\chapter{Zhu's modularity theorem}\label{ch:L15}
\begingroup\setlength{\emergencystretch}{3em}

\providecommand{\lfifteenH}{\mathbb{H}}
\providecommand{\lfifteenB}{\mathfrak{B}}
\providecommand{\lfifteenvac}{\mathbf{1}}
\providecommand{\lfifteenSL}{SL_2(\Z)}
\providecommand{\lfifteenZ}{\mathcal{Z}}
\providecommand{\lfifteenq}{\mathsf{q}}

This chapter documents the release of the LeanVOA library named in the front matter; it was first written against an earlier state of the library (``the snapshot'' below), and checks that were
run at an earlier state of the library and not re-run for this edition are marked as such. Every Lean quotation below is taken from the files of the release and is machine-checked by \texttt{scripts/textbook/sync.py} to
occur verbatim (up to whitespace and comments) in the source. The chapter covers the files assigned to it in
\texttt{Notes/Textbook/assignment.txt}: essentially all of \leanpath{VOA/\allowbreak Modularity/} except the modular-forms
toolkit (Eisenstein series, the ring $\C[E_4,E_6]$, the Serre derivative, the slash operator, evaluation of $q$-series;
Chapter~\ref{ch:L14}), the Arakawa--Kawasetsu files (Chapter~\ref{ch:L16}), the super files (Chapter~\ref{ch:L17}) and
the logarithmic files (Chapter~\ref{ch:L18}); plus four example files about the trivial vertex operator algebra
$\C\lfifteenvac$. That was $122$ files at the snapshot (two more were added for the release), roughly one and a half megabytes of Lean. We do not tour them. We state exactly
what the headline theorem says, present every definition on which its statement depends, explain Zhu's proof on paper
in the form in which the library actually carries it out, and describe the remaining files by role at the end
(Section~\ref{l15:sec:files}).

\emph{Release revision.} The chapter was revised for the first public release, applying the release audit \texttt{Notes/\allowbreak Textbook/\allowbreak referee/\allowbreak
RELEASE\_AUDIT\_ZHU.md} (2026-10-27). The quotations added in this revision --- the canonical-family headline
(Section~\ref{l15:sec:canonical}), the closed trivial instance (Section~\ref{l15:sec:nonvac}) and the first nontrivial
instances (Section~\ref{l15:sec:instances}) --- were copied from the files of that release; the older quotations were
re-verified there by \texttt{sync.py}. The two new source files are \leanpath{VOA/\allowbreak Modularity/\allowbreak
ZhuModularityCanonical.lean} (this chapter) and \leanpath{VOA/\allowbreak Classification/\allowbreak COne/\allowbreak
ZhuModularityV2mPlus.lean} (also assigned to this chapter, although it lives in the $c=1$ classification directory of
Chapter~\ref{ch:L22}, because its theorems are instances of this chapter's headline).

\emph{Second release revision (2026-10-28).} The chapter was revised again, applying finding F2 of \texttt{Notes/\allowbreak Textbook/\allowbreak referee/\allowbreak
RELEASE\_AUDIT\_NEW\_DECLS.md} and the entry ``Zhu for rational VOAs'' of \texttt{Notes/\allowbreak RELEASE\_PLAN.md}. The
first release revision said that the implication ``rational $\Rightarrow A(V)$ semisimple'' was not formalised and that the
lattice VOAs $V_{\sqrt{2m}\Z}$ were not covered. Both statements were false: the library already proved the implication
(\lean{VOA.Huang.ZhuSemisimple.}\allowbreak\lean{isSemisimpleRing_}\allowbreak\lean{A_of_}\allowbreak\lean{admissiblyRational}, in \leanpath{VOA/\allowbreak Huang/\allowbreak
ZhuSemisimple.lean}, Chapter~\ref{ch:L21}). It now also states Zhu's theorem in Zhu's own form, ``$C_2$-cofinite and
rational'' (\lean{zhuModularity_}\allowbreak\lean{canonical_of_}\allowbreak\lean{C2Cofinite_}\allowbreak\lean{admissiblyRational} and
\lean{zhuModularity_}\allowbreak\lean{canonical_of_}\allowbreak\lean{stronglyRational}, Section~\ref{l15:sec:rational}), and it proves the instances
\lean{zhuModularity_V2m} (every $V_{\sqrt{2m}\Z}$, $m\ge1$) and \lean{zhuModularity_of_nice} (every nice VOA;
Section~\ref{l15:sec:instances}). The quotations added in this revision were copied from the files at that revision. The
one open point left at that revision was the relation between the library's rationality predicate \lean{AdmissiblyRational}
and Dong--Li--Mason rationality (Remark~\ref{l15:rem:hyp}); since then the direction DLM $\Rightarrow$
\lean{AdmissiblyRational} is proved under $C_2$ (\lean{admissiblyRational_of_dlmRational}), and only the converse is open.

\emph{Third release revision (2026-10-28).} At this revision the library proves, in
\leanpath{VOA/\allowbreak Classification/\allowbreak COne/\allowbreak LiteratureInputsProved.lean} (Chapter~\ref{ch:L12},
Theorems~\ref{l12:thm:litproved} and~\ref{l12:thm:srlit}), the Dong--Li--Mason finiteness predicate
\lean{DLMFinitelyManyIrredAdmissible} for every $C_2$-cofinite $V$ (\lean{dlmFinitelyManyIrredAdmissible_of_c2Cofinite}), hence
for every strongly rational $V$ (\lean{StronglyRational.dlmFinitelyManyIrredAdmissible}), and restates Zhu's semisimplicity
predicate as the hypothesis-free \lean{zhuAdmissibleSemisimple_holds}. So both literature hypotheses \lean{hFin} and \lean{hZhu}
of the older packaging \lean{zhuModularity_stronglyRational_voaPBW} (Theorem~\ref{l15:thm:sr}) are now theorems for every
strongly rational $V$; that section was corrected (it previously called \lean{hFin} ``unproved''). No quotation changed.

Throughout we use, without repeating them, the following material from earlier chapters: vertex operator algebras,
modes $a_{(n)}$ and the library's mode convention (Chapter~\ref{ch:01}); ordinary modules \lean{OrdVModule}, their weight
spaces, the zero mode $o(v)$, graded traces and formal $q$-series, the predicates \lean{HasConformalWeight},
\lean{IrrModule}, $\Phi(V)$, Zhu's algebra $A(V)$, $C_2$-cofiniteness, and the rationality predicates
(Chapter~\ref{ch:L12}); the datum \lean{BottomProjection}, the PBW theorem \lean{VOA.pbw}, the predicates
\lean{IsSimpleAMod} and \lean{SimpleAModsRealized}, and Zhu's bijection (Chapter~\ref{ch:L13}); Eisenstein series,
slash operators and the analytic evaluation of $q$-series (Chapter~\ref{ch:L14}).

\paragraph{The findings of this chapter, in brief.} The reader who only wants the verdict should read this paragraph,
Section~\ref{l15:sec:headline}, and the warning boxes.
\begin{itemize}
\item \textbf{The headline, in Zhu's form.} The two theorems \lean{zhuModularity_}\allowbreak\lean{canonical_of_}\allowbreak\lean{C2Cofinite_}\allowbreak\lean{admissiblyRational}
  and \lean{zhuModularity_}\allowbreak\lean{canonical_of_}\allowbreak\lean{stronglyRational} (Theorem~\ref{l15:thm:rational})
  say: let $V$ be a \emph{rational, $C_2$-cofinite} vertex operator algebra (of CFT type, which every vertex operator algebra
  of the library is), where ``rational'' is the library's \lean{AdmissiblyRational} (every admissible module completely
  reducible); or let $V$ be strongly rational, which contains these two conditions. Then the conclusion displayed in the
  next item holds. These are \texttt{\#print axioms}-clean (the standard three axioms) and use the theorem \lean{VOA.pbw}.
  The step ``rational $\Rightarrow A(V)$ semisimple'' is the library theorem
  \lean{isSemisimpleRing_}\allowbreak\lean{A_of_}\allowbreak\lean{admissiblyRational} (Zhu's Theorem~2.2.3; Dong--Li--Mason).
\item \textbf{The engine: the semisimple form.} \lean{zhuModularity_}\allowbreak\lean{canonical_of_}\allowbreak\lean{C2Cofinite_semisimple}
  (Theorem~\ref{l15:thm:canonical}) is the more general form behind the headline. It says: let $V$ be a $C_2$-cofinite vertex
  operator algebra whose Zhu algebra $A(V)$ is a \emph{semisimple ring}. Then $\Phi(V)$, the set
  of isomorphism classes of irreducible ordinary $V$-modules, is finite, say with $r$ elements, and for the canonical
  representatives $M_1,\dots,M_r$ of these classes, with their conformal weights $h_1,\dots,h_r$, there is a group
  homomorphism $\rho:\lfifteenSL\to\mathrm{Mat}_r(\C)$ with $\rho(T)$ the diagonal matrix of $e^{2\pi i(h_j-c/24)}$ such
  that, for \emph{every} $\gamma=\left(\begin{smallmatrix}a&b\\ c&d\end{smallmatrix}\right)\in\lfifteenSL$, every $n\in\N$,
  every $v$ of square-bracket weight $n$ and every $\tau$ in the upper half-plane,
  \begin{equation*}
    \lfifteenZ_{M_i}\!\left(v,\tfrac{a\tau+b}{c\tau+d}\right)=(c\tau+d)^n\sum_{j=1}^r\rho(\gamma)_{ij}\,\lfifteenZ_{M_j}(v,\tau),
    \qquad \lfifteenZ_{M}(v,\tau)=\tr_{M}\,o(v)\,q^{L_0-c/24}.
  \end{equation*}
  The only hypotheses are $C_2$-cofiniteness and semisimplicity of $A(V)$. Finiteness of $\Phi(V)$ is a theorem
  (\lean{finite_Φ_of_C2Cofinite_pbw}), and the family is the canonical one, so no family data are assumed. The theorem is
  \texttt{\#print axioms}-clean (the standard three axioms) and uses the PBW theorem \lean{VOA.pbw}, which is a theorem.
  It covers \emph{one-point functions} only, not the $n$-point functions on the torus. This is Zhu's modular-invariance
  theorem for one-point functions (Zhu 1996, Theorem~5.3.3; Gannon's Theorem~5.3.8), with $\rho$ attached to the actual
  modules (``pinned'').
\item \textbf{The one remaining caveat: what ``rational'' means.} Zhu assumes $V$ rational and $C_2$-cofinite, and the
  headline assumes exactly that, with ``rational'' read as \lean{AdmissiblyRational}. That predicate asks every \emph{weak}
  submodule of an admissible module to have a weak complement. Dong--Li--Mason's rationality speaks of graded submodules.
  For $C_2$-cofinite $V$, DLM-rational implies \lean{AdmissiblyRational} \emph{in Lean}
  (\lean{VOA.Classification.COne.DLMRat.admissiblyRational_of_dlmRational}, Chapter~\ref{ch:L12},
  Theorem~\ref{l12:thm:dlmrat}; \texttt{\#print axioms}-clean, recorded by WS-RELEASE2; aggregate build at the release
  commit pending). Since the headline assumes $C_2$ anyway, it applies to every $C_2$-cofinite $V$ that is rational in
  Dong--Li--Mason's sense, with no paper step. The converse (\lean{AdmissiblyRational} $\Rightarrow$ DLM-rational) is open.
\item \textbf{Nontrivial instances.} \lean{zhuModularity_V2m} (Theorem~\ref{l15:thm:instancesV2m}) proves Zhu's theorem for
  every rank-one lattice VOA $V_{\sqrt{2m}\Z}$, $m\ge1$, from the library's \lean{V2m_stronglyRational};
  \lean{zhuModularity_of_nice} proves it for every \emph{nice} VOA (strongly rational and pseudo-unitary), so for every VOA
  that the library proves nice. \lean{zhuModularity_V2mPlus} and \lean{zhuModularity_cubicModel}
  (Theorem~\ref{l15:thm:instances}) discharge the hypotheses of the engine for the $c=1$ orbifolds $V_{\sqrt{2m}\Z}^+$,
  $m\ge2$, and for the tetrahedral orbifold, through a direct proof that $A(V)$ is semisimple. All four instances except
  \lean{zhuModularity_of_nice} have no hypotheses beyond the range of $m$, and all four use only the standard three axioms. These are vertex
  operator algebras with nonconstant characters at which the library's Zhu theorem says something. The trivial VOA
  $\C\lfifteenvac$ is covered with no hypotheses (\lean{trivialVOA_zhuModularity_closed}). \emph{Correction:} the first
  release revision of this chapter said that $V_{\sqrt{2m}\Z}$ was not covered and that ``rational $\Rightarrow A(V)$
  semisimple'' was unproved. Both claims were wrong, and so were the matching module docstrings
  (\texttt{RELEASE\_AUDIT\_NEW\_DECLS.md}, F2). The docstrings have since been corrected.
\item \textbf{The form for a given family.} \lean{zhuModularity_}\allowbreak\lean{pinned_of_}\allowbreak\lean{C2Cofinite}
  (Theorem~\ref{l15:thm:headline}) is the same law for any list $M_1,\dots,M_r$ of irreducible modules of length
  $r=\#\Phi(V)$, with conformal weights, whose bottoms are pairwise non-isomorphic when the weights agree; the headline is
  its specialisation to the canonical family. It is \texttt{\#print axioms}-clean, \emph{conditional on} the datum
  \lean{Pr : BottomProjection (parabolicRepOfAMod W (regularAMod W))}; the variant
  \lean{zhuModularity_}\allowbreak\lean{pinned_of_}\allowbreak\lean{C2Cofinite_voaPBW} supplies \lean{Pr} from the theorem
  \lean{VOA.pbw}.
\item \textbf{Convergence and holomorphy are proved, not assumed, but by a separate theorem.} The growth theorem
  \lean{onePointSubexp_of_c2} (from $C_2$-cofiniteness alone) says the coefficients of every one-point function grow
  subexponentially; the corollary \lean{qSummableAt_of_subexponential} turns that into convergence on $\lfifteenH$. Neither
  is part of the conclusion of the headline: the predicate \lean{ZhuModularityAnalytic} evaluates $q$-series with a
  \lean{tsum}, which returns $0$ on a divergent series. A separate
  Lean check against the library
  (\lean{headline_}\allowbreak\lean{traces_are_genuine}) combines them and proves, from the headline's hypotheses, that the
  functions in the conclusion are genuinely convergent, holomorphic and not identically zero. That check is not part of
  the library or of this release.
\item \textbf{The block count is an upper bound.} Zhu proves that the space of genus-one blocks has dimension exactly
  $\#\Phi(V)$. The library proves finite-dimensionality and $\dim\le\#\Phi(V)$
  (\lean{finiteDimensional_}\allowbreak\lean{and_finrank_}\allowbreak\lean{analyticBlocks_}\allowbreak\lean{le_card_Φ}); the equality is a
  consequence of the chapter's results on paper but is not a Lean theorem.
\item \textbf{Two redundant hypotheses, since removed from the headline} (checked in Lean; the check is not part of this release): \lean{hzero : VacuumGradeZero W}
  holds for every vertex operator algebra of the library (it is a weakening of an axiom), and \lean{hsimp} follows from
  \lean{hss} and \lean{Pr}. After this finding the library dropped both from \lean{zhuModularity_pinned_of_C2Cofinite} and its
  \lean{_voaPBW} wrapper (the quoted statements are the current ones); some auxiliary theorems still carry \lean{hzero}.
\item \textbf{What is not proved:} anything about the $S$-matrix beyond its existence (no unitarity, symmetry,
  $S^2=C$, Verlinde formula) for this chapter's $\rho$ (Chapter~\ref{ch:L21} proves Huang's Verlinde formula and the
  modular tensor category for nice VOAs); a modular linear differential equation with modular-form coefficients for $C_2$-cofinite
  $V$; Zhu's theorem for $n$-point functions; and the orbifold (twisted) version, which appears only as the
  \emph{statement} \lean{Conjecture_5_3_10} --- and that statement is mis-normalised (Section~\ref{l15:sec:twisted}).
\end{itemize}

\section{The theorem on paper, and the plan of the proof}\label{l15:sec:plan}

\subsection{Why characters should be modular}

Let $V$ be a vertex operator algebra of central charge $c$ and $M$ a module with conformal weight $h$, meaning that
$M=\bigoplus_{n\ge0}M_{h+n}$ with $L_0=h+n$ on $M_{h+n}$, $M_h\neq0$ and every $M_{h+n}$ finite-dimensional. Its
\emph{character}, or graded dimension, is the formal series
\begin{equation*}
  \chi_M(q)=\tr_M q^{L_0-c/24}=\sum_{n\ge0}\dim M_{h+n}\,q^{h+n-c/24}.
\end{equation*}
For the Moonshine module $\Vnat$ one has $\chi_{\Vnat}(q)=J(q)=q^{-1}+196884q+\cdots$, the normalised modular
invariant; for the lattice VOA $V_L$ of an even unimodular lattice one has $\chi_{V_L}=\theta_L/\eta^{\mathrm{rk}L}$; for
the Virasoro minimal models the characters are Rocha-Caridi's theta quotients. In every known rational example the
characters, viewed as functions of $\tau$ in the upper half-plane $\lfifteenH$ through $q=e^{2\pi i\tau}$, span a space
that is mapped to itself by $\tau\mapsto\gamma\tau=(a\tau+b)/(c\tau+d)$, $\gamma\in\lfifteenSL$. Zhu's theorem explains
this. More generally it treats the \emph{one-point functions}
\begin{equation*}
  \lfifteenZ_M(v,\tau)=\tr_M o(v)\,q^{L_0-c/24}=\sum_{n\ge0}\tr\bigl(o(v)|_{M_{h+n}}\bigr)\,q^{h+n-c/24},
\end{equation*}
where $o(v)$ is the zero mode of $v$ (for $v\in V_k$, $o(v)=v_{(k-1)}$, the mode that preserves each weight space;
Chapter~\ref{ch:L12}). For $v=\lfifteenvac$, $o(\lfifteenvac)=\mathrm{id}$ and $\lfifteenZ_M(\lfifteenvac,\tau)=\chi_M$.

\begin{theorem}[Zhu 1996, Theorem~5.3.3; Dong--Li--Mason 2000]\label{l15:thm:lit}
Let $V$ be rational and $C_2$-cofinite, with irreducible modules $M_1,\dots,M_r$. Then:
\begin{enumerate}
\item each $\lfifteenZ_{M_i}(v,\tau)$ converges to a holomorphic function on $\lfifteenH$;
\item there is a representation $\rho:\lfifteenSL\to GL_r(\C)$ such that for every $v$ in the square-bracket weight space
  $V_{[n]}$ (Section~\ref{l15:sec:square})
  \begin{equation*}
  \lfifteenZ_{M_i}(v,\gamma\tau)=(c\tau+d)^n\sum_{j}\rho(\gamma)_{ij}\lfifteenZ_{M_j}(v,\tau);
  \end{equation*}
\item the functions $\lfifteenZ_{M_1}(\cdot,\tau),\dots,\lfifteenZ_{M_r}(\cdot,\tau)$ are linearly independent, so $\rho$ is
  unique.
\end{enumerate}
\end{theorem}

\begin{physics}
This is the mathematical form of \emph{modular invariance of the torus} in two-dimensional conformal field theory.
A chiral CFT placed on a torus $\C/(\Z+\Z\tau)$ has a space of conformal blocks; the torus partition function of the
full (left-right) theory is $Z(\tau,\bar\tau)=\sum_{i,j}N_{ij}\chi_i(\tau)\overline{\chi_j(\tau)}$, and it must not depend
on how the torus was presented, i.e.\ it must be invariant under the mapping class group $\lfifteenSL$ of the torus.
That is only possible if the $\chi_i$ span a representation of $\lfifteenSL$ --- which is what Zhu proves. The one-point
function $\lfifteenZ_M(v,\tau)$ is the torus expectation value $\langle v\rangle_\tau$ of the field $v$ inserted at one point
of the torus with the $M$-sector propagating around the cycle; inserting a field of conformal weight $n$ (in the
torus coordinate) turns the partition function into a modular form of weight $n$, which is the factor $(c\tau+d)^n$.
The matrix $\rho(S)$, $S=\left(\begin{smallmatrix}0&-1\\1&0\end{smallmatrix}\right)$, is the \emph{modular $S$-matrix}: it
exchanges the space and time cycles of the torus, and through the Verlinde formula it determines the fusion rules.
\end{physics}

\begin{remark}[How the library uses rationality]\label{l15:rem:hyp}
Theorem~\ref{l15:thm:lit} is the statement in the literature, with the literature's hypothesis ``rational''. Zhu uses
rationality in exactly two places: to know that $\Phi(V)$ is finite and that $A(V)$ is a semisimple algebra (his
Theorem~2.2.3; Dong--Li--Mason for the version with their definition of rationality). The library proves the first from
$C_2$-cofiniteness alone (\lean{finite_Φ_of_C2Cofinite_pbw}, Section~\ref{l15:sec:canonical}). It proves the second from
its rationality predicate (\lean{isSemisimpleRing_}\allowbreak\lean{A_of_}\allowbreak\lean{admissiblyRational}, Theorem~\ref{l15:thm:zhuss}). The theorem is
stated in two layers. The \emph{engine} (Theorem~\ref{l15:thm:canonical}) has the two hypotheses
\begin{equation*}
  \text{(a) } V \text{ is } C_2\text{-cofinite},\qquad \text{(b) } A(V) \text{ is a semisimple ring}.
\end{equation*}
Rationality implies (b), on paper and (for \lean{AdmissiblyRational}) in Lean, so the engine implies Zhu's theorem. The
converse is not claimed:
Dong and Jiang characterise rationality by semisimplicity of \emph{all} the higher Zhu algebras $A_n(V)$, so (b) alone is
not known to imply rationality, and the library's form may be more general. The \emph{headline}
(Theorem~\ref{l15:thm:rational}) has Zhu's hypotheses,
\begin{equation*}
  \text{(a) } V \text{ is } C_2\text{-cofinite},\qquad \text{(r) } V \text{ is rational (\texttt{AdmissiblyRational})},
\end{equation*}
and the step (r) $\Rightarrow$ (b) is the Lean theorem \lean{isSemisimpleRing_}\allowbreak\lean{A_of_}\allowbreak\lean{admissiblyRational}. (An earlier
revision of this remark said that this step was missing and existed only as the named hypothesis
\lean{ZhuAdmissibleSemisimple} of Chapter~\ref{ch:L12}. That was wrong: the library proves the step, and
\lean{zhuAdmissibleSemisimple_of_Pr} proves the named predicate itself.)

\emph{What remains open} is the meaning of ``rational''. The library's \lean{AdmissiblyRational} asks for complements of
weak submodules; Dong--Li--Mason ask for complements of graded submodules. The library's own notes record that the relation
is not settled (\leanpath{VOA/\allowbreak Modules/\allowbreak AdmissibleRationality.lean}; Chapter~\ref{ch:L12},
Section~\ref{l12:warn:admcr}); one direction is now a theorem: for $C_2$-cofinite $V$, DLM-rationality implies
\lean{AdmissiblyRational} (\lean{admissiblyRational_of_dlmRational}, Theorem~\ref{l12:thm:dlmrat}). So for a $C_2$-cofinite
$V$ that is rational in the paper's sense the headline applies, through that Lean theorem; only the converse is open. For a $V$ that the
library proves \lean{AdmissiblyRational}, such as the lattice VOAs $V_{\sqrt{2m}\Z}$ (Section~\ref{l15:sec:instances}), it
applies with nothing on paper. Throughout this chapter ``the library proves Zhu's theorem for rational $C_2$-cofinite $V$''
therefore means rational in the sense of \lean{AdmissiblyRational}.
\end{remark}

\subsection{Zhu's proof in seven steps}\label{l15:sec:steps}

We describe the proof as the library carries it out. Each step is a section of this chapter.
\begin{enumerate}
\item \textbf{The second VOA structure.} Zhu changes coordinates from the plane to the cylinder, $z\mapsto e^{z}-1$,
  obtaining new vertex operators $Y[v,z]=Y(e^{zL_0}v,e^z-1)$, new modes $v[n]$ and a new Virasoro algebra $L[n]$ with the
  same central charge (Section~\ref{l15:sec:square}).
\item \textbf{The recursion.} By a direct computation with traces and the Jacobi identity, the one-point functions of
  every module satisfy three identities, Zhu's \emph{genus-one block conditions}: $\lfifteenZ(a[0]v)=0$; a recursion
  expressing $\lfifteenZ(a[-2]v)$ through Eisenstein series times $\lfifteenZ(a[2k-2]v)$; and an identity expressing
  $\frac{1}{2\pi i}\frac{d}{d\tau}\lfifteenZ(v)$ through $\lfifteenZ(L[-2]v)$ and Eisenstein series
  (Section~\ref{l15:sec:recursion}). No finiteness hypothesis is needed for this.
\item \textbf{Convergence.} Under $C_2$-cofiniteness the formal $q$-series converge on all of $\lfifteenH$
  (Section~\ref{l15:sec:convergence}).
\item \textbf{The block space.} Call a \emph{genus-one block} any family $F(v,\tau)$, linear in $v$ and holomorphic in
  $\tau$, satisfying the three block conditions. They form a vector space $\lfifteenB$; the trace functions are in it
  (Section~\ref{l15:sec:blocks}).
\item \textbf{Finiteness and the count.} $\dim\lfifteenB<\infty$ (Zhu's 5.3.2) and $\dim\lfifteenB\le\#\Phi(V)$ (Zhu's
  5.3.3); the latter uses semisimplicity of $A(V)$ and the existence half of Zhu's correspondence, which itself rests on
  the Gaberdiel--Neitzke spanning theorem (Sections~\ref{l15:sec:finite}--\ref{l15:sec:gn}).
\item \textbf{The action.} The block conditions are covariant under $\lfifteenSL$, because Eisenstein series $G_{2k}$,
  $k\ge2$, are modular and the defect of $G_2$ is compensated by the derivative term (Section~\ref{l15:sec:slash}).
\item \textbf{Reading off $\rho$.} The $r=\#\Phi(V)$ trace functions are linearly independent
  (Section~\ref{l15:sec:independence}); they lie in a space of dimension $\le r$, so they are a basis; the slash of a
  trace function is again a block, hence a combination of traces; this defines $\rho$, and uniqueness makes it a
  homomorphism (Section~\ref{l15:sec:rho}). The headline is assembled in Section~\ref{l15:sec:headline}.
\end{enumerate}

\begin{remark}[What is different from Zhu]\label{l15:rem:diff}
Two steps are carried out differently from Zhu's paper, and the reader comparing with the literature should know where.
(a) Convergence (step 3): Zhu derives a linear differential equation in $q\frac{d}{dq}$ with modular-form coefficients
and uses the theory of regular singular points. The library does not; it proves that the span of the values of a
formal block over a ring of ``subexponential'' power series is finitely generated and stable under $q\frac{d}{dq}$, and
deduces subexponential growth of the coefficients from the resulting \emph{first-order system}
(Section~\ref{l15:sec:convergence}). (b) Finite-dimensionality (step 5): Zhu proves $\dim\lfifteenB<\infty$ by the same
differential equation. The library proves it with the identity theorem: condition (iii) expresses the derivative of a
block through values of the block, so a block vanishing at one point on finitely many vectors vanishes identically
(Section~\ref{l15:sec:finite}). Both are legitimate proofs; neither uses a modular linear differential equation.
\end{remark}

\section{One-point functions on the torus}\label{l15:sec:traces}

\subsection{Formal one-point functions}

Recall from Chapter~\ref{ch:L12} (Definition~\ref{l12:def:char}) that a \lean{FormalQSeries} is a function
$\chi:\C\to\C$, read as the formal sum $\sum_{e}\chi(e)q^e$ over complex exponents; that \lean{oMod N v} is the zero
mode $o(v)$ on the module \lean{N}; and that \lean{onePoint N v} is the formal series whose coefficient at $q^e$ is the
trace of $o(v)$ on the weight space $M_{e+c/24}$. The modularity files re-export it under a torus name.

\paragraph{A note on two gradings.} Two gradings of $V$ occur in this chapter. $V_n$ is the usual $L_0$-grading.
$V_{[n]}$ (Lean: \lean{sqWeightSpace W n}) is the eigenspace with eigenvalue $n$ of $L[0]$, the zero mode of a
\emph{different} Virasoro algebra $L[m]$, the ``square-bracket'' or torus Virasoro algebra constructed in
Section~\ref{l15:sec:square} (Definition~\ref{l15:def:sqmode}, Theorem~\ref{l15:thm:sqvir}). $L[0]=L_0+(\text{operators lowering the}
\ L_0\text{-degree})$, so $V_{[n]}\neq V_n$ in general, although they have the same dimensions. The modular weight in
Zhu's law is the square-bracket weight $n$ of $v\in V_{[n]}$; the zero mode $o(v)$ is always computed in the usual grading.

\begin{definition}[Torus one-point function]\label{l15:def:torusonepoint}
For an ordinary module $M$ and $v\in V$, the \emph{torus one-point function} is the formal series
$\lfifteenZ_M(v;q)=\sum_{e\in\C}\tr\bigl(o(v)|_{M_{e+c/24}}\bigr)q^{e}$, i.e.\ $\tr_M o(v)q^{L_0-c/24}$; the
\emph{partition function} is $\lfifteenZ_M(\lfifteenvac;q)=\chi_M(q)$. The map $v\mapsto\lfifteenZ_M(v;q)$ is linear.
\end{definition}
\noindent File \leanpath{VOA/\allowbreak Modularity/\allowbreak OnePointFunctions.lean}.
\begin{leancode}
noncomputable def torusOnePoint (N : OrdVModule W M) (v : 𝒱) : FormalQSeries := onePoint N v

noncomputable def torusPartition (N : OrdVModule W M) : FormalQSeries := gradedDim N

theorem torusOnePoint_vacuum (N : OrdVModule W M) : torusOnePoint N W.vacuum = torusPartition N :=
\end{leancode}
\noindent File \leanpath{VOA/\allowbreak Modularity/\allowbreak ZhuBlockConditions.lean}.
\begin{leancode}
noncomputable def torusOnePointHom (N : OrdVModule W M) : 𝒱 →ₗ[ℂ] FormalQSeries where
  toFun := torusOnePoint N
  map_add' := torusOnePoint_add N
  map_smul' := fun c v => torusOnePoint_smul N c v
\end{leancode}

\begin{dictionary}
\textbf{Decoding.} \lean{𝒱} is the underlying space of $V$ and \lean{W : VertexOperatorAlgebra 𝒱} the VOA structure.
\lean{OrdVModule W M} is an ordinary module structure on the space \lean{M} (graded by $L_0$-eigenvalues
$\alpha\in\C$ with finite-dimensional weight spaces, Definition~\ref{l12:def:ordmodule}). \lean{𝒱 →ₗ[ℂ] FormalQSeries}
is the type of $\C$-linear maps; \lean{torusOnePointHom N} is $v\mapsto\lfifteenZ_M(v;q)$ packaged as such a map, the
linearity being the two proved lemmas named in the fields.

\textbf{Exponent convention.} The coefficient of \lean{torusOnePoint N v} at the complex number $e$ is the trace on
the weight space of $L_0$-eigenvalue $e+c/24$. So for a module of conformal weight $h$ the series is supported on
$\{h-c/24+n:n\in\N\}$, as in $\tr_M o(v)q^{L_0-c/24}$. We shall check this convention at every place where a $q$-series
is evaluated.

\textbf{Fidelity.} Same as the textbook definition. No grading hypothesis on $v$ is needed: $o(v)$ is defined for every
$v$ by decomposing $v$ into homogeneous components (Chapter~\ref{ch:L12}).
\end{dictionary}

\begin{physics}
$\lfifteenZ_M(v;q)$ is the torus one-point function $\langle v\rangle_{M,\tau}$: the path integral on a torus of modulus
$\tau$ with one insertion of the field $v$ and the states of the sector $M$ running around the time cycle, weighted
by $e^{2\pi i\tau(L_0-c/24)}$. The shift $-c/24$ is the Casimir energy of the cylinder: mapping the plane to the cylinder
shifts $L_0$ by $-c/24$ (Section~\ref{l15:sec:square}). For $v=\lfifteenvac$ one gets the chiral partition function of
the sector $M$.
\end{physics}

\begin{status}
Definitions. \texttt{\#print axioms VOA.torusOnePointHom\_isFormalBlock} (which uses them): \texttt{propext,
Classical.choice, Quot.sound}.
\end{status}

\subsection{From formal series to functions on the upper half-plane}

\begin{definition}[Evaluation at $\tau$; the trace functions]\label{l15:def:traceFun}
For $\tau\in\lfifteenH$ put $q=e^{2\pi i\tau}$, so $|q|<1$. For a formal series $\chi$ and an exponent offset
$h_0\in\C$ define
\begin{equation*}
  \mathrm{ev}_{h_0}(\chi)(\tau)=e^{2\pi i\tau h_0}\sum_{n\ge0}\chi(h_0+n)\,q^n,
\end{equation*}
the sum being the (unconditional) sum of the series when it converges, and $0$ by convention otherwise. For a family
$(M_i)_{i\in\iota}$ of irreducible modules with conformal weights $h_i$, the \emph{trace function} is
\begin{equation*}
  S_i(v,\tau)=\mathrm{ev}_{h_i-c/24}\bigl(\lfifteenZ_{M_i}(v;q)\bigr)(\tau)=\sum_{n\ge0}\tr\bigl(o(v)|_{(M_i)_{h_i+n}}\bigr)e^{2\pi i\tau(h_i+n-c/24)}.
\end{equation*}
\end{definition}
\noindent File \leanpath{VOA/\allowbreak Modularity/\allowbreak ConvergenceBridge.lean}.
\begin{leancode}
noncomputable def qParam (τ : UpperHalfPlane) : ℂ :=
  Complex.exp (2 * (Real.pi : ℂ) * Complex.I * (τ : ℂ))

def QSummableAt (χ : FormalQSeries) (h : ℂ) (τ : UpperHalfPlane) : Prop :=
  Summable (fun n : ℕ => χ (h + n) * (qParam τ) ^ n)

noncomputable def qEvalAt (χ : FormalQSeries) (h : ℂ) (τ : UpperHalfPlane) : ℂ :=
  Complex.exp (2 * (Real.pi : ℂ) * Complex.I * (τ : ℂ) * h) *
    ∑' n : ℕ, χ (h + n) * (qParam τ) ^ n
\end{leancode}
\noindent File \leanpath{VOA/\allowbreak Modularity/\allowbreak BlockSpanning.lean}.
\begin{leancode}
noncomputable def traceFun (Mfam : ι → IrrModule W) (h : ι → ℂ) (i : ι) (v : 𝒱)
    (τ : UpperHalfPlane) : ℂ :=
  qEvalAt (torusOnePoint (Mfam i).isMod v) (h i - W.centralCharge / 24) τ
\end{leancode}

\begin{dictionary}
\textbf{Decoding.} \lean{UpperHalfPlane} is Mathlib's $\lfifteenH=\{\tau:\mathrm{Im}\,\tau>0\}$; \lean{(τ : ℂ)} its
coercion to $\C$. \lean{∑' n : ℕ, f n} is Mathlib's \lean{tsum}: the sum of $f$ if $f$ is summable, and $0$
otherwise. \lean{Summable} is unconditional convergence (for $\C$-valued sequences: absolute convergence of the
series). \lean{Mfam : ι → IrrModule W} is a family of bundled irreducible modules (Definition~\ref{l12:def:phi}) indexed
by an arbitrary type \lean{ι}; \lean{(Mfam i).isMod} is the module structure. \lean{h : ι → ℂ} is a separately supplied
list of numbers, which the theorems below tie to the modules by the hypothesis
\lean{hcw : ∀ i, HasConformalWeight (Mfam i).isMod (h i)}.

\textbf{Exponents.} \lean{traceFun} evaluates along the offset $h_i-c/24$, so the $n$-th term reads the coefficient
$\chi(h_i-c/24+n)$, which by the convention of Definition~\ref{l15:def:torusonepoint} is the trace on
$(M_i)_{h_i+n}$, and multiplies by $e^{2\pi i\tau(h_i-c/24)}q^n$. That is exactly
$\tr_{M_i}o(v)q^{L_0-c/24}$ with the standard branch $q^{s}:=e^{2\pi i\tau s}$ for the complex power. The convention is
consistent; compare the twisted case in Section~\ref{l15:sec:twisted}, where it is not.

\textbf{The junk value.} Because \lean{tsum} returns $0$ on a divergent series, \lean{traceFun} is a total function
whether or not the trace series converges. Every theorem below that is \emph{about} convergence says so with
\lean{QSummableAt} or with the growth predicate \lean{OnePointSubexp} (Definition~\ref{l15:def:subexp}); the headline's
conclusion does not, and we return to this in the fidelity check of the headline.
\end{dictionary}

\begin{physics}
$S_i(v,\tau)$ is the torus one-point function as an honest function of the modulus $\tau$. Physically the convergence
of the sum is the statement that the sector $M_i$ has a finite partition function at every positive temperature
$\mathrm{Im}\,\tau>0$, i.e.\ that the number of states grows slower than exponentially in the energy. For a rational CFT
the growth is $e^{2\pi\sqrt{c_{\mathrm{eff}}n/6}}$ (Cardy's formula), far below any exponential.
\end{physics}

\begin{status}
Definitions; no axioms beyond Mathlib's. The analytic facts about \lean{qEvalAt} (holomorphy, termwise derivative,
behaviour under $\tau\mapsto\tau+1$) are in Chapter~\ref{ch:L14}.
\end{status}

\subsection{The \texorpdfstring{$T$}{T}-transformation: proved for every module}

The first and easiest part of modularity needs no hypothesis on $V$ at all.

\begin{theorem}[Zhu's law at $T$]\label{l15:thm:T}
Let $M$ have conformal weight $h$ and let $v\in V$ be arbitrary. Then
$\lfifteenZ_M(v,\tau+1)=e^{2\pi i(h-c/24)}\lfifteenZ_M(v,\tau)$, as formal $q$-series (where ``$\tau\mapsto\tau+1$'' acts on
$q^e$ by $e^{2\pi i e}$). Consequently, for a finite family $(M_i,h_i)$ the law of Theorem~\ref{l15:thm:lit}(2) at
$\gamma=T=\left(\begin{smallmatrix}1&1\\0&1\end{smallmatrix}\right)$ holds with $\rho(T)=\mathrm{diag}\bigl(e^{2\pi i(h_i-c/24)}\bigr)$.
\end{theorem}
\noindent File \leanpath{VOA/\allowbreak Modularity/\allowbreak ZhuTheorem.lean}.
\begin{leancode}
noncomputable def modularTEntry (W : VertexOperatorAlgebra 𝒱) (h : ℂ) : ℂ :=
  Complex.exp (2 * (Real.pi : ℂ) * Complex.I * (h - W.centralCharge / 24))

noncomputable def modularRepT {ι : Type*} [DecidableEq ι] (W : VertexOperatorAlgebra 𝒱)
    (h : ι → ℂ) : ι → ι → ℂ :=
  fun i j => if i = j then modularTEntry W (h i) else 0

theorem onePoint_T_transform (N : OrdVModule W M) (v : 𝒱) (h : ℂ) (hcw : HasConformalWeight N h) :
    qShift (onePoint N v)
      = (Complex.exp (2 * (Real.pi : ℂ) * Complex.I * (h - W.centralCharge / 24)))
          • onePoint N v :=

theorem zhu_modularity_T_onePoint {ι : Type*} [Fintype ι] [DecidableEq ι]
    {W : VertexOperatorAlgebra 𝒱} (D : ModularData W ι) (v : 𝒱) (i : ι) :
    qShift (torusOnePoint (D.Mfam i).isMod v)
      = ∑ j : ι, D.rep ModularGroup.T i j • torusOnePoint (D.Mfam j).isMod v := by
\end{leancode}

\begin{dictionary}
\lean{qShift χ} is the formal series $e\mapsto e^{2\pi ie}\chi(e)$ (Chapter~\ref{ch:L12}, Theorem~\ref{l12:thm:Tperiod}),
i.e.\ the effect of $\tau\mapsto\tau+1$ on $q^e$. \lean{•} is scalar multiplication. \lean{ModularData} is defined in
Section~\ref{l15:sec:rho}; its field \lean{rep_T} forces \lean{D.rep ModularGroup.T} to be \lean{modularRepT}, so the second
theorem is the first one read through that field. The analytic version (functions on $\lfifteenH$ rather than formal
series) is \lean{zhu_modularity_}\allowbreak\lean{T_analytic} in \leanpath{ConvergenceBridge.lean}, also unconditional. Same as the
textbook statement; no hypothesis on $V$.
\end{dictionary}

\begin{proof}[Proof on paper]
If $M_{e+c/24}\neq0$ then $e+c/24=h+n$ with $n\in\N$ (this is the second half of \lean{HasConformalWeight}), so
$e^{2\pi ie}=e^{2\pi i(h-c/24)}e^{2\pi in}=e^{2\pi i(h-c/24)}$; if $M_{e+c/24}=0$ both sides have coefficient $0$.
\end{proof}

\begin{physics}
$T$ is the Dehn twist of the torus. Its eigenvalue on the sector $M$ is $e^{2\pi i(h-c/24)}$: the twist rotates every
state of energy $h+n-c/24$ by a phase, and the integer part $n$ drops out. The $T$-matrix of a rational CFT is therefore
diagonal with entries $e^{2\pi i(h_i-c/24)}$, the ``topological spins'' (up to the central-charge phase).
\end{physics}

\begin{status}
\lean{VOA.char_T_modularity}, \lean{VOA.zhu_}\allowbreak\lean{modularity_}\allowbreak\lean{T_onePoint}: \texttt{propext, Classical.choice, Quot.sound}.
Unconditional apart from \lean{HasConformalWeight}.
\end{status}

\section{The square-bracket vertex operator algebra}\label{l15:sec:square}

\subsection{The change of coordinates}

The torus is the quotient of the cylinder $\C/\Z$ (coordinate $z$) by $z\sim z+\tau$; the cylinder is the punctured plane
(coordinate $w$) through $w=e^{2\pi iz}$. Zhu works with $w=e^{z}-1$ (so that $z=0$ goes to $w=0$; the $2\pi i$ is put
back at the end). Vertex operators transform under a change of coordinates, and for the map $z\mapsto e^z-1$ Zhu proves
that
\begin{equation*}
  Y[v,z]:=Y\bigl(e^{zL_0}v,\,e^z-1\bigr)=\sum_{n\in\Z}v[n]z^{-n-1}
\end{equation*}
defines a \emph{second} vertex operator algebra structure on the same space $V$, with the same vacuum and the conformal
vector $\tilde\omega=\omega-\frac{c}{24}\lfifteenvac$. For $v\in V_k$ the change of variable $w=e^z-1$ in a residue gives
\begin{equation}\label{l15:eq:sqmode}
  v[m]=\Res_z Y[v,z]z^m=\Res_w Y(v,w)(1+w)^{k-1}\bigl(\log(1+w)\bigr)^m=\sum_{i\ge m}c(k,m,i)\,v_{(i)},
\end{equation}
where $c(k,m,i)$ is the coefficient of $w^i$ in $(1+w)^{k-1}(\log(1+w))^m$. For $m<0$, $(\log(1+w))^m$ is understood as
$w^m u(w)^m$ with $u(w)=\log(1+w)/w=1-w/2+\cdots$, a unit of $\C[[w]]$; for $k=0$, $(1+w)^{-1}$ is the inverse unit. The
sum in \eqref{l15:eq:sqmode} is finite on each vector, because $v_{(i)}x=0$ for $i$ large.

\begin{definition}[Square-bracket modes]\label{l15:def:sqmode}
For $k\in\N$, $m,i\in\Z$ let $c(k,m,i)=[w^{i-m}]\bigl((1+w)^{k-1}u(w)^m\bigr)$ if $i\ge m$ and $0$ otherwise. For
$v\in V$ and a declared weight $k$, the operator $v[m]_k=\sum_i c(k,m,i)v_{(i)}$. The torus Virasoro modes are
$L[n]=\omega[n+1]_2-\frac{c}{24}\lfifteenvac[n+1]_0$.
\end{definition}
\noindent File \leanpath{VOA/\allowbreak Modularity/\allowbreak OnePointFunctions.lean}.
\begin{leancode}
noncomputable def torusConformalVector : 𝒱 :=
  W.conformalVector - (W.centralCharge / 24) • W.vacuum

noncomputable def sqBracketCoeffZ (k : ℕ) (m i : ℤ) : ℂ :=
  if 0 ≤ i - m then
    PowerSeries.coeff (R := ℂ) (i - m).toNat
      ((↑(onePlusXUnit ^ ((k : ℤ) - 1)) : PowerSeries ℂ) * (↑(uUnit ^ m) : PowerSeries ℂ))
  else 0

noncomputable def vSqZ (v : 𝒱) (k : ℕ) (m : ℤ) : 𝒱 →ₗ[ℂ] 𝒱 where
  toFun x := ∑ᶠ i : ℤ, sqBracketCoeffZ k m i • mode (W.Y v) i x

noncomputable def Lsq (n : ℤ) : Module.End ℂ 𝒱 :=
  vSqZ W W.conformalVector 2 (n + 1)
    - (W.centralCharge / 24) • vSqZ W W.vacuum 0 (n + 1)

theorem Lsq_neg_one : Lsq W (-1) = W.L (-1) + W.L 0 := by
\end{leancode}

\begin{dictionary}
\textbf{Decoding.} \lean{onePlusXUnit} is $1+w$ as a unit of $\C[[w]]$ and \lean{uUnit} is $u(w)=\log(1+w)/w$
(\lean{uSeries}, coefficients $(-1)^i/(i+1)$) as a unit, so \lean{↑(onePlusXUnit ^ ((k : ℤ) - 1))} is $(1+w)^{k-1}$
for all $k\in\N$, including $k=0$, and \lean{↑(uUnit ^ m)} is $u^m$ for all $m\in\Z$. \lean{PowerSeries.coeff j}
is the coefficient of $w^j$. So \lean{sqBracketCoeffZ k m i} is $c(k,m,i)$. \lean{mode (W.Y v) i} is $v_{(i)}$
(Chapter~\ref{ch:01}). \lean{∑ᶠ} is a sum over a set of finite support (\lean{finsum}); the support is finite because
$c(k,m,i)=0$ for $i<m$ and $v_{(i)}x=0$ for large $i$ (\lean{sqModeZ_support_finite}). \lean{vSqZ W v k m} is the
operator $v[m]_k$.

\textbf{The declared weight $k$.} The square-bracket mode depends on a natural number $k$ supplied separately; it is the
genuine $v[m]$ of \eqref{l15:eq:sqmode} only when $v\in V_k$. Every theorem below that uses \lean{vSqZ W a ka m}
carries the hypothesis \lean{a ∈ W.grading (ka : ℤ)}, so this is harmless: the library defines $v[m]$ on homogeneous
vectors and extends linearly where needed (\lean{sqMode} in \leanpath{SqGradedBlock.lean}).

\textbf{$L[n]$.} \lean{Lsq W n} is $\omega[n+1]-\frac c{24}\lfifteenvac[n+1]$, i.e.\ $\tilde\omega[n+1]$, Zhu's $L[n]$.
The vacuum term matters: \lean{𝟙[n+1]} is the identity for $n=-2$ and $0$ otherwise, so $L[-2]=\omega[-1]-\frac c{24}$.
The file's docstring explains that without it the $\omega[n+1]$ satisfy a Virasoro algebra with a spurious extra
cocycle. \lean{Lsq_neg_one} is Zhu's $L[-1]=L_{-1}+L_0$ (Gannon writes $2\pi i(L_{-1}+L_0)$ because he uses
$w=e^{2\pi iz}-1$).

\textbf{Fidelity.} These are Zhu's definitions verbatim in the $w=e^z-1$ normalisation (Zhu 1996, \S4.2). The
Bernoulli numbers in Zhu's recursion are, by \lean{subst_expMinusOne_uSeries} (\leanpath{SquareBracketSubst.lean}), the
coefficients of $u$ read in the variable $z$: $u(e^z-1)=z/(e^z-1)$. This is a proof that the coefficients are the right
ones, not a numerical check.
\end{dictionary}

\begin{physics}
Mapping the plane to the cylinder is radial quantisation in reverse: the cylinder $\mathbb{R}\times S^1$ is the natural
space-time of a closed string, and $L[0]$ is the Hamiltonian that generates translations along the cylinder. The
Schwarzian derivative of $e^{z}-1$ produces the constant $-\frac{c}{24}$ in the stress tensor: that is the Casimir energy of
the cylinder, and it is the origin of $q^{-c/24}$ in every character. $L[-1]=L_{-1}+L_0$ is the cylinder translation
operator.
\end{physics}

\begin{status}
\lean{VOA.Lsq_neg_one}: \texttt{propext, Classical.choice, Quot.sound}. The coefficients were also checked by the
project against Mathematica tables; that check is evidence,
the Lean definitions and theorems are the proof.
\end{status}

\subsection{The torus Virasoro algebra and the square-bracket grading}

\begin{theorem}[Zhu; Gannon 5.3.16]\label{l15:thm:sqvir}
The operators $L[n]$ satisfy the Virasoro relations with the same central charge:
$[L[m],L[n]]=(m-n)L[m+n]+\frac{m^3-m}{12}\delta_{m+n,0}\,c$. The operator $L[0]$ is diagonalisable on $V$: $V$ is the
internal direct sum of the eigenspaces $V_{[k]}=\ker(L[0]-k)$, and only $k\in\N$ occur.
\end{theorem}
\noindent Files {\raggedright\leanpath{VOA/\allowbreak Modularity/\allowbreak GelfandFuks.lean} (a file of Chapter~\ref{ch:L14}),
\leanpath{VOA/\allowbreak Modularity/\allowbreak SquareBracketGrading.lean},
\leanpath{VOA/\allowbreak Modularity/\allowbreak SqWeightSpanning.lean}, \leanpath{VOA/\allowbreak Modularity/\allowbreak SqWeightNat.lean}.\par}
\begin{leancode}
theorem Lsq_virasoro_apply' (W : VertexOperatorAlgebra 𝒱) (m n : ℤ) (x : 𝒱) :
    Lsq W m (Lsq W n x) - Lsq W n (Lsq W m x)
      = ((m - n : ℤ) : ℂ) • Lsq W (m + n) x
        + (if m + n = 0 then (((m : ℂ) ^ 3 - (m : ℂ)) / 12 * W.centralCharge) else 0) • x :=

noncomputable def sqWeightSpace (W : VertexOperatorAlgebra 𝒱) (k : ℂ) : Submodule ℂ 𝒱 :=
  LinearMap.ker (Lsq W 0 - k • (LinearMap.id : 𝒱 →ₗ[ℂ] 𝒱))

theorem sqWeightSpace_isInternal (W : VertexOperatorAlgebra 𝒱) :
    DirectSum.IsInternal (sqWeightSpace W) :=

theorem sqWeightSpace_nat_spans_top (W : VertexOperatorAlgebra 𝒱) :
    (⨆ m : ℕ, sqWeightSpace W ((m : ℕ) : ℂ)) = ⊤ := by
\end{leancode}

\begin{dictionary}
\lean{LinearMap.ker (Lsq W 0 - k • id)} is $\ker(L[0]-k)$. \lean{DirectSum.IsInternal (sqWeightSpace W)} says that the
natural map $\bigoplus_{k\in\C}V_{[k]}\to V$ is a linear isomorphism, i.e.\ $V$ is the direct sum of the $L[0]$-eigenspaces.
\lean{⨆ m : ℕ, …} is the sum of submodules, so the last statement says that the eigenspaces with $k\in\N$ already span.
\lean{Lsq_virasoro_apply'} is the Virasoro relation applied to a vector \lean{x}; the central term is
$\frac{m^3-m}{12}c\,\delta_{m+n,0}$. All as in Zhu.
\end{dictionary}

\begin{proof}[Proof on paper]
Write $L[n]=\sum_a d(n,a)L_a$ (finite on each vector) with $d(n,a)=c(2,n+1,a+1)$, plus the vacuum term. The commutator
$[L[m],L[n]]$ is a double sum of Virasoro commutators $[L_a,L_b]=(a-b)L_{a+b}+\frac{a^3-a}{12}\delta_{a+b,0}c$, so the
relation reduces to two identities among the coefficients $d$: a ``Witt'' identity
$\sum_a d(m,a)d(n,j-a)(2a-j)=(m-n)d(m+n,j)$ and a ``central'' identity for $\sum_a d(m,a)d(n,-a)(a^3-a)$. Both are
identities of formal power series: they express that $z\mapsto e^z-1$ is a coordinate change, under which vector fields
transform as vector fields and the central term transforms with the Schwarzian derivative (Gelfand--Fuks cocycle). The
second identity produces an extra term $-\frac{m-n}{2}\delta_{m+n,-2}$ which is exactly cancelled by the vacuum term
$-\frac c{24}\lfifteenvac[n+1]$. For the grading: $L[0]=L_0+\sum_{a\ge1}d(0,a)L_a$ is $L_0$ plus operators that strictly
lower the $L_0$-degree. On the finite-dimensional space $V_{\le N}=\bigoplus_{k\le N}V_k$, $L[0]$ is therefore triangular
with diagonal $L_0$, hence diagonalisable with eigenvalues in $\{0,\dots,N\}$ (triangular with distinct diagonal blocks
that are scalar). Eigenvectors for distinct eigenvalues are independent, and every vector lies in some $V_{\le N}$.
\end{proof}

\begin{physics}
The torus has its own conformal symmetry algebra, generated by $L[n]$; $L[0]$ is the energy measured on the cylinder,
and its eigenvalue $n$ on $v$ is the modular weight with which $v$ enters the torus one-point function. That the central
charge is unchanged is the statement that $c$ is an intrinsic property of the theory, independent of the surface.
\end{physics}

\begin{status}
\lean{VOA.Lsq_virasoro_apply'} and \lean{VOA.sqWeightSpace_}\allowbreak\lean{isInternal}: \texttt{propext, Classical.choice, Quot.sound};
unconditional, for every VOA of the library. (\lean{Lsq_virasoro_apply} takes a hypothesis
\lean{CentralMomentClosed}, discharged by \lean{centralMomentClosed_holds}.)
\end{status}

\begin{remark}[What of the second VOA structure is formalised]\label{l15:rem:secondvoa}
The library does \emph{not} construct $(V,Y[\cdot,z],\lfifteenvac,\tilde\omega)$ as a \lean{VertexOperatorAlgebra}. It has
the modes $v[m]$, the Virasoro relation for $L[n]$, the $L[0]$-grading, and the specific commutator facts that the block
conditions need. The general statement ``$a[m]$ shifts the $L[0]$-weight by $\wt[a]-m-1$'' appeared as a carried datum
\lean{SqModeWeighted} in an earlier route; the route of the headline avoids it (Section~\ref{l15:sec:slash}). This is a
legitimate economy, not a gap in the headline.
\end{remark}

\section{Zhu's recursion and the genus-one block conditions}\label{l15:sec:recursion}

\subsection{Eisenstein series in Zhu's normalisation, and convolution}

\begin{definition}[Zhu's Eisenstein series; convolution]\label{l15:def:zhuG}
For $j\ge1$ let
\begin{equation*}
  G_{2j}(q)=-\frac{B_{2j}}{(2j)!}+\frac{2}{(2j-1)!}\sum_{n\ge1}\sigma_{2j-1}(n)q^n,
\end{equation*}
with $B_{2j}$ the Bernoulli numbers and $\sigma_s(n)=\sum_{d\mid n}d^s$. Equivalently $G_{2j}=\frac{2\zeta(2j)}{(2\pi i)^{2j}}E_{2j}$
with $E_{2j}=1+O(q)$ the normalised Eisenstein series; $G_2=-\frac1{12}E_2$. For a sequence $f:\N\to\C$ (a power series
$\sum_m f(m)q^m$) and a formal series $\chi$, the \emph{convolution} is $(f\star\chi)(e)=\sum_{m\ge0}f(m)\chi(e-m)$, the
coefficient of $q^e$ in the product of the two series.
\end{definition}
\noindent Files \leanpath{VOA/\allowbreak Modularity/\allowbreak ZhuNormalization.lean},
\leanpath{VOA/\allowbreak Modularity/\allowbreak ZhuRecursion.lean}, \leanpath{VOA/\allowbreak Modularity/\allowbreak EisensteinAnalytic.lean}
(Chapter~\ref{ch:L14}).
\begin{leancode}
noncomputable def zhuG (j : ℕ) : ℕ → ℂ :=
  fun m =>
    if m = 0 then -(((bernoulli (2 * j) : ℚ) : ℂ) / (Nat.factorial (2 * j) : ℂ))
    else (2 / (Nat.factorial (2 * j - 1) : ℂ)) * lambertCoeff (2 * j - 1) m

noncomputable def qConv (f : ℕ → ℂ) (g : FormalQSeries) : FormalQSeries :=
  fun e => ∑ᶠ n : ℕ, f n * g (e - (n : ℂ))

noncomputable def zhuGC (j : ℕ) (z : ℂ) : ℂ :=
  ∑' n : ℕ, zhuG j n * Complex.exp (2 * (Real.pi : ℂ) * Complex.I * (n : ℂ) * z)
\end{leancode}

\begin{dictionary}
\lean{bernoulli} is Mathlib's Bernoulli number (with $B_1=-\frac12$; only even indices occur here, where conventions
agree). \lean{lambertCoeff (2 * j - 1) m} is $\sigma_{2j-1}(m)$ (\lean{lambertCoeff_eq_sigma}). So \lean{zhuG j} is the
coefficient sequence of $G_{2j}$, and \lean{zhuGC j z} is the function $G_{2j}(e^{2\pi iz})$. \lean{qConv f g} is the
convolution (a \lean{finsum}: it is finite on the series that occur, which are supported on a coset $h_0+\N$).
\textbf{Fidelity.} This is Zhu's normalisation (Zhu 1996, \S3; Gannon \S5.3.5): $\frac{2\zeta(2j)}{(2\pi i)^{2j}}=-\frac{B_{2j}}{(2j)!}$.
The project first settled the constant by exact rational computations on the Heisenberg VOA, where other normalisations
were falsified; but that is not what certifies it here: the recursion below is \emph{proved} with this family.
\end{dictionary}

\begin{physics}
$G_{2k}$ are the lattice sums $\sum_{(m,n)\neq0}(m\tau+n)^{-2k}$ up to normalisation; they arise as the torus propagator: the
Weierstrass functions $\wp$ and $\zeta$ expanded at the origin have the $G_{2k}$ as Laurent coefficients. Zhu's recursion
is the torus version of the operator product expansion: moving the field $a$ around the torus and collecting the
singular part of its OPE with $v$ produces exactly these coefficients. $G_2$ is not modular; physically this is the
holomorphic anomaly of the torus propagator.
\end{physics}

\begin{status}
Definitions. \lean{zhuG_eq_}\allowbreak\lean{smul_zhuEisenstein} relates \lean{zhuG} to the $E$-normalised series under the hypothesis
$B_{2j}\neq0$ (true, carried there because Mathlib lacks it); nothing downstream uses that comparison.
\end{status}

\subsection{The recursion}

\begin{theorem}[Zhu's recursion; Zhu 1996, Prop.~4.3.2]\label{l15:thm:recursion}
Let $M$ be any ordinary module, $a\in V_{k}$ with $k\ge1$, and $v\in V$ arbitrary. Then, as formal $q$-series,
\begin{equation*}
  \lfifteenZ_M(a[-1]v;q)=\tr_M o(a)o(v)q^{L_0-c/24}+\sum_{j\ge1}G_{2j}(q)\,\lfifteenZ_M(a[2j-1]v;q),
\end{equation*}
the sum being finite.
\end{theorem}
\noindent Files \leanpath{VOA/\allowbreak Modularity/\allowbreak OnePointFunctions.lean},
\leanpath{VOA/\allowbreak Modularity/\allowbreak ZhuRecursion.lean}, \leanpath{VOA/\allowbreak Modularity/\allowbreak ZhuNormalization.lean},
\leanpath{VOA/\allowbreak Modularity/\allowbreak ZhuRecursionClose.lean}.
\begin{leancode}
noncomputable def torusTwoPoint (N : OrdVModule W M) (a v : 𝒱) : FormalQSeries :=
  fun e =>
    haveI := N.finrank (e + W.centralCharge / 24)
    LinearMap.trace ℂ _
      (weightBlock N ((oMod N a).comp (oMod N v)) (e + W.centralCharge / 24))

def ZhuRecursionAt (N : OrdVModule W M) (G : ℕ → (ℕ → ℂ)) (a v : 𝒱) (ka : ℕ) : Prop :=
  torusOnePoint N (vSqZ W a ka (-1) v)
    = torusTwoPoint N a v
      + ∑ᶠ k : ℕ, qConv (G (k + 1))
          (torusOnePoint N (vSqZ W a ka (2 * (k + 1) - 1 : ℤ) v))

def ZhuRecursionPinned (N : OrdVModule W M) (a v : 𝒱) (ka : ℕ) : Prop :=
  ZhuRecursionAt N zhuG a v ka

theorem zhuRecursionPinned_holds {ka : ℕ} (hk : 1 ≤ ka) {a : 𝒱} (ha : a ∈ W.grading (ka : ℤ))
    (v : 𝒱) : ZhuRecursionPinned N a v ka := by
\end{leancode}

\begin{dictionary}
\lean{weightBlock N A α} is the block of the operator $A$ on the weight space $M_\alpha$ and \lean{LinearMap.trace} its
trace (Chapter~\ref{ch:L12}); \lean{(oMod N a).comp (oMod N v)} is $o(a)o(v)$. So \lean{torusTwoPoint N a v} is
$\tr_M o(a)o(v)q^{L_0-c/24}$ with the exponent convention of Definition~\ref{l15:def:torusonepoint}. In
\lean{ZhuRecursionAt}, \lean{G (k + 1)} stands for $G_{2(k+1)}$ and \lean{vSqZ W a ka (2 * (k + 1) - 1 : ℤ) v} for
$a[2k+1]v$; reindexing $j=k+1$ gives the displayed sum. The theorem quantifies over every module \lean{N} (a section
variable), every $a$ of weight $k_a\ge1$, and every $v$ --- there is no homogeneity hypothesis on $v$ and no hypothesis
on $V$. Same as Zhu's Proposition~4.3.2 (whose statement is for homogeneous $v$; the extension to all $v$ is by
linearity). The case $a=\lfifteenvac$, $k_a=0$ is the trivial \lean{zhuRecursionAt_}\allowbreak\lean{vacuum_left}.
\end{dictionary}

\begin{warning}
\textbf{A wrong, unused interface (documentation only).} The file \leanpath{VOA/\allowbreak Modularity/\allowbreak OnePointFunctions.lean}
still contains the earlier
\begin{leancode}
def ZhuRecursion (N : OrdVModule W M) (G : ℕ → FormalQSeries) (a v : 𝒱) (ka : ℕ) : Prop :=
  torusOnePoint N (vSqZ W a ka (-1) v)
    = torusTwoPoint N a v
      + ∑ᶠ k : ℕ, G (k + 1) * torusOnePoint N (vSqZ W a ka (2 * (k + 1) - 1 : ℤ) v)
\end{leancode}
in which \lean{G (k + 1) * torusOnePoint …} is the \emph{pointwise} product of two functions $\C\to\C$, not the product of
$q$-series. As a statement about $q$-series it is meaningless (the docstring of \lean{qConv} says so itself). It is used
nowhere (a repository-wide search finds no consumer), so no theorem is affected; but a reader must not take it for
Zhu's recursion. Also stale: the docstring of \leanpath{ZhuNormalization.lean} says that ``\lean{ZhuRecursionAt N zhuG}
actually holds remains the open theorem'', and its final ``Residual'' section repeats it; \lean{zhuRecursionPinned_holds}
proves it.
\end{warning}

\begin{proof}[Proof on paper]
Take $a\in V_k$, $v\in V_l$ homogeneous. By \eqref{l15:eq:sqmode}, $a[-1]v=\sum_{i\ge-1}c(k,-1,i)a_{(i)}v$, with
$c(k,-1,-1)=1$. So the left side is $\lfifteenZ(a_{(-1)}v)+\sum_{i\ge0}c(k,-1,i)\lfifteenZ(a_{(i)}v)$. The heart of the matter
is the computation of $\lfifteenZ(a_{(i)}v)=\tr o(a_{(i)}v)q^{L_0-c/24}$. The Jacobi identity at $m=0$ (the
``iterate formula'') writes $o(a_{(i)}v)$ as a finite combination of products $a_{(p)}v_{(k+l-p-2)}$ and
$v_{(k+l-p-2)}a_{(p)}$ of modes, each of total degree zero. For such a ``compensated pair'' one uses cyclicity of the trace
together with $q^{L_0}a_{(p)}=q^{p-k+1}a_{(p)}q^{L_0}$:
\begin{equation*}
\tr\,a_{(p)}v_{(r)}q^{L_0}=\tr\,v_{(r)}q^{L_0}a_{(p)}=q^{\,k-1-p}\,\tr\,v_{(r)}a_{(p)}q^{L_0},
\end{equation*}
and the commutator $[a_{(p)},v_{(r)}]=\sum_j\binom pj(a_{(j)}v)_{(p+r-j)}$. Solving the resulting geometric relation gives,
for $p\neq k-1$, $\tr a_{(p)}v_{(r)}q^{L_0}=\frac{1}{1-q^{k-1-p}}\sum_j\binom pj\tr o(a_{(j)}v)q^{L_0}$, an expansion in
$q$ whose coefficients are the divisor-type sums $\sum_n n^{s}q^{nm}$. Collecting these geometric series over all $p$ and
all $i$ is a purely combinatorial identity, and the result is exactly the Lambert series
$\frac{2}{(2j-1)!}\sum_n\sigma_{2j-1}(n)q^n$ multiplying $\lfifteenZ(a[2j-1]v)$; the constant terms collect into
$-B_{2j}/(2j)!$ because the coefficients of $u(w)$ in the variable $z$ are the Bernoulli generating function. The only
term not of this form is $p=k-1$, which gives $\tr o(a)o(v)q^{L_0-c/24}$ with coefficient exactly $1$. The library's
proof is this computation, organised as an ``abstract assembly'' identity for the coefficient bookkeeping
(\leanpath{ZhuCoefficientIdentity.lean}) and the trace computation (\leanpath{ZhuRecursionClose.lean}).
\end{proof}

\begin{physics}
Zhu's recursion is the torus Ward identity: the one-point function of the descendant $a[-1]v$ (the normal-ordered product
of the fields $a$ and $v$ on the cylinder) equals the two-point function of the zero modes plus a sum of torus propagator
coefficients times one-point functions of the OPE coefficients $a[2j-1]v$. It is the higher-genus analogue of the
statement that a correlation function on the sphere is determined by its OPE singularities.
\end{physics}

\begin{status}
\lean{VOA.zhuRecursionPinned_holds}: \texttt{propext, Classical.choice, Quot.sound}. Unconditional: every VOA, every
module, every homogeneous $a$ of positive weight, every $v$.
\end{status}

\subsection{The three block conditions}

\begin{definition}[Formal genus-one block; Zhu 1996, Def.~4.3, 5.1]\label{l15:def:formalblock}
A linear map $S:V\to\{\text{formal }q\text{-series}\}$ is a \emph{formal genus-one block} if for every $k\ge1$, every
$a\in V_k$ and every $v\in V$:
\begin{enumerate}
\item[(i)] $S(a[0]v)=0$;
\item[(ii)] $S(a[-2]v)+\sum_{j\ge1}(2j-1)\,G_{2j}\star S(a[2j-2]v)=0$;
\item[(iii)] and for every $v$: $S(L[-2]v)=q\frac{d}{dq}S(v)+\sum_{j\ge1}G_{2j}\star S(L[2j-2]v)$.
\end{enumerate}
Here $q\frac d{dq}$ acts on $q^e$ by multiplication by $e$.
\end{definition}
\noindent File \leanpath{VOA/\allowbreak Modularity/\allowbreak ZhuBlockConditions.lean}.
\begin{leancode}
noncomputable def qDeriv (χ : FormalQSeries) : FormalQSeries := fun e => e * χ e

structure IsFormalBlock (W : VertexOperatorAlgebra 𝒱) (S : 𝒱 →ₗ[ℂ] FormalQSeries) : Prop where
  vanish : ∀ (ka : ℕ), 1 ≤ ka → ∀ a ∈ W.grading (ka : ℤ), ∀ v : 𝒱, S (vSqZ W a ka 0 v) = 0
  negTwo : ∀ (ka : ℕ), 1 ≤ ka → ∀ a ∈ W.grading (ka : ℤ), ∀ v : 𝒱,
    S (vSqZ W a ka (-2) v)
      + ∑ᶠ m : ℕ, ((2 * (m + 1) - 1 : ℕ) : ℂ)
          • qConv (zhuG (m + 1)) (S (vSqZ W a ka (2 * (m + 1) - 2 : ℤ) v)) = 0
  omega : ∀ v : 𝒱, S (Lsq W (-2) v)
    = qDeriv (S v) + ∑ᶠ m : ℕ, qConv (zhuG (m + 1)) (S (Lsq W (2 * (m + 1) - 2 : ℤ) v))
\end{leancode}

\begin{physics}
A formal genus-one block is a candidate assignment of torus one-point functions to all fields, written as power series in
$q$ (the torus seen from the degeneration limit of a long thin cylinder), and required to obey the torus Ward identities
term by term. It is the algebraic shadow of a conformal block before any question of convergence is asked.
\end{physics}

\begin{theorem}[The trace functions are formal blocks; Zhu 1996, Prop.~4.3.5--4.3.6]\label{l15:thm:formalblock}
For every ordinary module $M$ of every VOA, $v\mapsto\lfifteenZ_M(v;q)$ is a formal genus-one block.
\end{theorem}
\noindent File \leanpath{VOA/\allowbreak Modularity/\allowbreak ZhuBlockConditions.lean}.
\begin{leancode}
theorem torusOnePointHom_isFormalBlock (N : OrdVModule W M) : IsFormalBlock W (torusOnePointHom N) where
\end{leancode}

\begin{dictionary}
The three fields are the three conditions with $j=m+1$: \lean{(2 * (m + 1) - 1 : ℕ)} is $2j-1$ and
\lean{2 * (m + 1) - 2} is $2j-2$. \lean{qDeriv} is $q\frac d{dq}$ in the exponent convention ($e\cdot\chi(e)$). The
quantifier ``$a\in V_k$, $k\ge1$'' excludes only $V_0=\C\lfifteenvac$, for which (i) and (ii) are trivial. The $j=1$ term
of (ii) involves $G_2\star S(a[0]v)$, which vanishes by (i); the library keeps it, Zhu drops it --- equivalent. Same as
Zhu's definition.
\end{dictionary}

\begin{proof}[Proof on paper]
(i): for $a\in V_k$, $a[0]=\sum_i\binom{k-1}{i}a_{(i)}$, and the commutator formula gives
$[o(a),o(v)]=\sum_i\binom{k-1}{i}\,o(a_{(i)}v)=o(a[0]v)$ (the ``$\ell=0$'' case of the Jacobi identity). Since $o(a)$ and $o(v)$
preserve each finite-dimensional weight space $M_{h+n}$, the trace of their commutator there is $0$. (ii) is the recursion
(Theorem~\ref{l15:thm:recursion}) applied to $L[-1]a$ in place of $a$: in the square-bracket structure
$(L[-1]a)[m]=-m\,a[m-1]$, so $(L[-1]a)[-1]=a[-2]$ and $(L[-1]a)[2j-1]=-(2j-1)a[2j-2]$, which produces the factor $2j-1$;
and the two-point term vanishes because $o(L[-1]a)=o(L_{-1}a)+o(L_0a)=-k\,o(a)+k\,o(a)=0$ (apply the recursion separately to
the homogeneous pieces $L_{-1}a\in V_{k+1}$ and $L_0a\in V_k$). (iii) is the recursion with $a=\omega$: $o(\omega)=L_0$, so
$\tr o(\omega)o(v)q^{L_0-c/24}=q\frac d{dq}\lfifteenZ(v)+\frac c{24}\lfifteenZ(v)$, and the $\frac c{24}$ is exactly absorbed
by the vacuum term in $L[-2]=\omega[-1]-\frac c{24}$.
\end{proof}

\begin{physics}
Conditions (i)--(iii) are the Ward identities of the chiral algebra on the torus. (i) says that zero modes of currents act
by commutators, which are invisible in a trace (charge conservation around the cycle). (ii) is the OPE identity for the
descendant $a[-2]v$. (iii) is the Ward identity of the stress tensor: inserting $T$ on the torus is the same as
differentiating with respect to the modulus $\tau$ (the torus has a one-dimensional moduli space, and $T$ is its tangent
vector) plus OPE terms. Condition (iii) is the seed of every differential equation for characters.
\end{physics}

\begin{status}
\lean{VOA.torusOnePointHom_}\allowbreak\lean{isFormalBlock}: \texttt{propext, Classical.choice, Quot.sound}. Unconditional.
\end{status}

\section{Convergence from \texorpdfstring{$C_2$}{C2}-cofiniteness}\label{l15:sec:convergence}

\subsection{Subexponential growth}

A power series $\sum a_nq^n$ converges on the whole unit disc $|q|<1$ exactly when $\limsup|a_n|^{1/n}\le1$, i.e.\ when
the coefficients grow slower than every exponential $r^n$, $r>1$.

\begin{definition}[Subexponential sequences; the growth and support data]\label{l15:def:subexp}
A sequence $a:\N\to\mathbb{R}$ is \emph{subexponential} if for every $r>1$ there is $C$ with $a_n\le Cr^n$ for all $n$. A module
$M$ with exponent offset $h_0$ has \emph{subexponential one-point functions} if for every $v\in V$ the sequence
$n\mapsto|\lfifteenZ_M(v;q)[q^{h_0+n}]|$ is subexponential, and \emph{supported one-point functions} if
$\lfifteenZ_M(v;q)[q^{h_0+i}]=0$ for all $v$ and all negative integers $i$.
\end{definition}
\noindent Files \leanpath{VOA/\allowbreak Modularity/\allowbreak ConvergenceBridge.lean},
\leanpath{VOA/\allowbreak Modularity/\allowbreak AnalyticBlock.lean}.
\begin{leancode}
def Subexponential (a : ℕ → ℝ) : Prop :=
  ∀ r : ℝ, 1 < r → ∃ C : ℝ, ∀ n : ℕ, a n ≤ C * r ^ n

theorem qSummableAt_of_subexponential (χ : FormalQSeries) (h : ℂ)
    (hsub : Subexponential (fun n : ℕ => ‖χ (h + n)‖)) (τ : UpperHalfPlane) :
    QSummableAt χ h τ := by

def OnePointSupported (N : OrdVModule W M) (h : ℂ) : Prop :=
  ∀ (v : 𝒱) (i : ℤ), i < 0 → torusOnePoint N v (h + (i : ℂ)) = 0

def OnePointSubexp (N : OrdVModule W M) (h : ℂ) : Prop :=
  ∀ v : 𝒱, Subexponential (fun n : ℕ => ‖torusOnePoint N v (h + (n : ℂ))‖)
\end{leancode}

\begin{dictionary}
\lean{‖χ (h + n)‖} is the complex absolute value of the coefficient. The quantifier order in \lean{Subexponential} is
the right one ($C$ may depend on $r$); with $\exists C\,\forall r$ the notion would be bounded growth, which is false for
characters. \lean{qSummableAt_}\allowbreak\lean{of_subexponential} is the elementary comparison with a geometric series of ratio
$r|q|<1$. The offset $h$ in \lean{OnePointSubexp N h} is an \emph{exponent}; the theorems below always use
$h=h_M-c/24$, matching \lean{traceFun}.
\end{dictionary}

\begin{physics}
Subexponential growth of $\dim M_{h+n}$ is the statement that the density of states grows slower than $e^{\beta E}$ for
every inverse temperature $\beta>0$: the theory has no Hagedorn temperature. Rational CFTs (and, more generally, all
$C_2$-cofinite theories) have Cardy growth $e^{2\pi\sqrt{c_{\mathrm{eff}}E/6}}$.
\end{physics}

\begin{status}
Definitions and an elementary lemma: \texttt{propext, Classical.choice, Quot.sound}.
\end{status}

\subsection{The vacuum hypothesis}

\begin{definition}[Vacuum spans the degree-zero space]\label{l15:def:vgz}
$V_0\subseteq\C\lfifteenvac$.
\end{definition}
\noindent File \leanpath{VOA/\allowbreak Modularity/\allowbreak BlockZeroMode.lean}.
\begin{leancode}
def VacuumGradeZero : Prop := W.grading 0 ≤ Submodule.span ℂ ({W.vacuum} : Set 𝒱)
\end{leancode}

\begin{dictionary}
\lean{W.grading 0} is $V_0$ and \lean{Submodule.span ℂ {W.vacuum}} is $\C\lfifteenvac$; \lean{≤} is inclusion of subspaces.
Its docstring calls it ``the only part of `CFT type' that is not already a VOA axiom''.
\end{dictionary}

\begin{warning}
\textbf{\lean{VacuumGradeZero} is automatic (a redundant hypothesis, and a stale docstring).} The library's
\lean{VertexOperatorAlgebra} has the axiom \lean{vacuum_grade : W.grading 0 = Submodule.span ℂ {W.vacuum}}
(Chapter~\ref{ch:01}), so \lean{VacuumGradeZero W} holds for every $V$: it is \lean{le_of_eq W.vacuum_grade}. The library
knows this (\lean{QuasiLisseSl2.vacuumGradeZero_}\allowbreak\lean{of_vacuum_grade} in \leanpath{AKSl2Payoff.lean}, and the strongly
rational wrapper of Section~\ref{l15:sec:headline} discharges it that way), but the hypothesis \lean{hzero} is still
carried by \lean{onePointSubexp_of_c2}, by Zhu's 5.3.3 and by both headlines. Checked in Lean against the library:
\lean{VOA.L15Cert.vacuumGradeZero_always}
(compiled with \texttt{-DmaxSynthPendingDepth=3}; standard three axioms; the check is not part of this release). It costs nothing and hides nothing, but a
reader counting hypotheses should strike it.
\end{warning}

\begin{physics}
``CFT type'': the vacuum is the unique state of zero energy. Physically this is uniqueness of the ground state of the
chiral theory on the circle.
\end{physics}

\subsection{The growth theorem}

\begin{theorem}[Convergence of one-point functions; Zhu 1996 \S4.4, Dong--Li--Mason 2000 \S\S6--7]\label{l15:thm:growth}
Let $V$ be $C_2$-cofinite and $M$ an ordinary module with conformal weight $h$. Then for every $v\in V$ the coefficients
of $\lfifteenZ_M(v;q)$ along $h-c/24+\N$ grow subexponentially. Consequently every trace function $S_i(v,\tau)$ is given by
an absolutely convergent series on $\lfifteenH$.
\end{theorem}
\noindent File \leanpath{VOA/\allowbreak Modularity/\allowbreak BlockGrowth.lean}.
\begin{leancode}
theorem onePointSubexp_of_c2 {M : Type u} [AddCommGroup M] [Module ℂ M]
    (N : OrdVModule W M) {hw : ℂ} (hcw : HasConformalWeight N hw)
    (hzero : VacuumGradeZero W) (hC2 : C2Cofinite W) :
    OnePointSubexp N (hw - W.centralCharge / 24) := by
\end{leancode}

\begin{dictionary}
\lean{C2Cofinite W} is $\dim V/C_2(V)<\infty$ (Definition~\ref{l12:def:c2}). \lean{HasConformalWeight N hw}: $M_{hw}\neq0$
and every weight of $M$ lies in $hw+\N$ (Chapter~\ref{ch:L12}). \lean{hzero} is redundant (warning above). The
conclusion is for every $v$, not only $v=\lfifteenvac$. The ``consequently'' is
\lean{qSummableAt_}\allowbreak\lean{of_subexponential}; it is assembled in a separate Lean check, \lean{headline_}\allowbreak\lean{traces_are_genuine}. Exactly the
textbook statement, with no rationality hypothesis --- as in Miyamoto and Dong--Li--Mason, convergence needs only $C_2$.
\end{dictionary}

\begin{proof}[Proof on paper]
Let $R$ be the ring of power series $\sum f_mq^m$ whose coefficients grow subexponentially; it contains every $G_{2j}$
(coefficients $O(m^{2j})$) and is closed under products and under $q\frac d{dq}$. Fix a formal block $S$ (here
$S=\lfifteenZ_M$) whose values lie in $q^{h_0}\C[[q]]$, and let $\mathcal S=R\cdot S(V)$ be the $R$-span of its values.
\emph{Finite generation.} $C_2$-cofiniteness gives finitely many homogeneous $x_1,\dots,x_g$ spanning $V$ modulo $C_2(V)$
(\lean{exists_homog_}\allowbreak\lean{gen_of_C2Cofinite}). Condition (ii) of the block expresses $S(a[-2]v)$, and therefore $S$ of every
element of $C_2(V)$ up to strictly lower degree, as an $R$-combination of values of $S$ on vectors of lower degree. By
induction on the degree, $\mathcal S$ is generated over $R$ by $S(x_1),\dots,S(x_g),S(\lfifteenvac)$. (Finiteness over the
Noetherian ring $\C[G_2,G_4,G_6]$ is used to make this work uniformly.) \emph{Stability under $q\frac d{dq}$.} By
condition (iii), $q\frac d{dq}S(v)=S(L[-2]v)-\sum_jG_{2j}\star S(L[2j-2]v)\in\mathcal S$. So there is a $g'\times g'$ matrix
$A$ over $R$ with $q\frac d{dq}Y=AY$ for the vector $Y$ of generators. \emph{Growth.} Writing $Y=\sum_n y_nq^{h_0+n}$, the
system reads $(h_0+n)y_n=\sum_{j\le n}A_jy_{n-j}$; since $|h_0+n|\sim n$, a majorant argument shows that if the $A_j$ grow
subexponentially then so do the $y_n$ (\leanpath{SystemGrowth.lean}). Finally every $S(v)$ is an $R$-combination of the
generators, and convolution with a subexponential series preserves subexponential growth
(\lean{subexponential_qConv}).
\end{proof}

\begin{physics}
The first-order system $q\frac{d}{dq}Y=AY$ is the torus version of the Knizhnik--Zamolodchikov-type equations: the
genus-one conformal blocks form a flat vector bundle over the moduli space of tori, and $A$ is its connection in the
coordinate $q$. Convergence of the characters is then a statement about solutions of an ODE with a regular singularity at
the cusp $q=0$, which is how physicists usually think of it (Mathur--Mukhi--Sen).
\end{physics}

\begin{status}
\lean{VOA.onePointSubexp_of_c2}: \texttt{propext, Classical.choice, Quot.sound}. Conditional only on \lean{C2Cofinite}
(and the redundant \lean{hzero}) and \lean{HasConformalWeight}. Read exactly, it is a \emph{growth} statement: the
coefficients of $\lfifteenZ_M(v;q)$ are subexponential. Convergence on $\lfifteenH$ is the separate corollary
\lean{VOA.qSummableAt_of_subexponential} (standard three axioms, re-checked at an earlier commit (not re-run for this edition)); the two together give
the ``consequently'' of Theorem~\ref{l15:thm:growth}. Neither appears in the conclusion of the headline
(Section~\ref{l15:sec:headline}).
\end{status}

\begin{warning}
\textbf{A refuted route to convergence (dead end, recorded in the library).} Zhu's Lemma~4.4.4 converts the tower
$S(L[-2]^iv)$ into the tower $(q\frac d{dq})^jS(v)$ unitriangularly. Read naively, \emph{for every} $v$ --- the
hypothesis \lean{huni} of \lean{onePointSubexp_}\allowbreak\lean{of_unitriangular} and \lean{onePointSubexp_}\allowbreak\lean{of_unitriangular_cw}
(\leanpath{ZhuODEAssembly.lean}) --- it is \emph{false}: at $i=1$ and $v=L[-1]x$ the right side is built from
$\lfifteenZ(L[-1]x)=0$, while the left side is $\lfifteenZ(L[-2]L[-1]x)=-\lfifteenZ(L[-3]x)$ by the Virasoro relation, so it would
force $\lfifteenZ(L[-n]x)=0$ for all $n\ge3$. Certificate in the library: \lean{torusOnePoint_}\allowbreak\lean{Lsq_neg_}\allowbreak\lean{three_eq_}\allowbreak\lean{zero_of_uniStep}
in \leanpath{VOA/\allowbreak Modularity/\allowbreak UniCollapse.lean}:
\begin{leancode}
theorem torusOnePoint_Lsq_neg_three_eq_zero_of_uniStep
    (huni1 : ∀ v : 𝒱, ∃ hc : ℕ → PowerSeries ℂ,
      torusOnePoint N ((⇑(Lsq W (-2)))^[1] v)
        = ∑ j ∈ Finset.range 2, psSMul (hc j) (qDeriv^[j] (torusOnePoint N v)))
    (x : 𝒱) : torusOnePoint N (Lsq W (-3) x) = 0 := by
\end{leancode}
Consequently the two theorems named above have an unsatisfiable hypothesis for any module with a nonzero
$\lfifteenZ(L[-3]x)$; they are ``dead'' (the library's own word) but still present. The correct statement holds for
square-bracket \emph{primary} $v$ only (Arakawa--Kawasetsu, Lemma~5.5; Chapter~\ref{ch:L16}). The headline uses neither:
it uses the first-order system above.
\end{warning}

\begin{warning}
\textbf{Stale documentation (documentation only).} The header of \leanpath{ConvergenceBridge.lean} says that ``\lean{SubexponentialCoeffs}
is carried as a named \lean{Prop}'' because ``the ordering and the count are not built''; the header of
\leanpath{OnePointGrowth.lean} says the repository's growth theorem ``only covers $v=\lfifteenvac$''; the docstring of
\lean{ZhuModularityWeightedAt} (\leanpath{ZhuTheorem.lean}) says the convergence input is ``present in neither this repo nor
Mathlib''; and the header of \leanpath{OnePointODE.lean} says \lean{hode} ``is not claimed to be discharged''. All four
predate \lean{onePointSubexp_of_c2}, which proves the growth datum for every $v$ from $C_2$-cofiniteness.
\end{warning}

\subsection{What about the modular differential equation?}\label{l15:sec:mlde}

Physicists know Zhu's convergence argument as the existence of a \emph{modular linear differential equation} (MLDE):
the characters of a rational VOA satisfy
\begin{equation*}
  \Bigl(D^{s}+\sum_{i<s}\phi_i(\tau)D^{i}\Bigr)\chi=0,\qquad D=q\tfrac{d}{dq}-\tfrac{k}{12}E_2\ \text{(Serre derivative on weight }k),
\end{equation*}
with $\phi_i$ holomorphic modular forms of weight $2(s-i)$ (Mathur--Mukhi--Sen 1988; Zhu \S4.4; Kaneko--Zagier for $s=2$).
The library does \emph{not} prove this for $C_2$-cofinite $V$. What exists:
\begin{itemize}
\item \lean{exists_monic_}\allowbreak\lean{Lsq_relation} (\leanpath{FormalBlockFinite.lean}): Zhu's Lemma~4.4.2, a monic relation among
  $S(L[-2]^iv)$ with coefficients in $\C[G_2,G_4,G_6]$;
\item \lean{character_}\allowbreak\lean{formalODE_}\allowbreak\lean{of_etingofFinite} (Chapter~\ref{ch:L16}): for the \emph{character} ($v=\lfifteenvac$),
  a monic linear ODE in $q\frac d{dq}$ with \emph{subexponential} (not modular) coefficients, under \lean{EtingofFinite}
  (finite-dimensional zeroth Poisson homology; that quasi-lisse implies Etingof-finite is Etingof--Schedler and is not
  formalised); since $C_2$-cofinite implies \lean{EtingofFinite} (\lean{EtingofFinite_}\allowbreak\lean{of_C2Cofinite}), this
  covers $C_2$-cofinite $V$ too. By Chapter~\ref{ch:L16}, Proposition~\ref{l16:prop:weak}, this equation is near-vacuous:
  every convergent $q$-series of the relevant kind satisfies one.
\end{itemize}
Neither is the MLDE with modular-form coefficients. The headline does not need it.

\section{The space of genus-one blocks}\label{l15:sec:blocks}

\begin{definition}[Analytic genus-one block; the block space $\lfifteenB$]\label{l15:def:analyticblock}
An \emph{analytic genus-one block} is a map $F:V\times\lfifteenH\to\C$, $F(v,z)$, linear in $v$ and holomorphic in $z$,
such that for every $k\ge1$, $a\in V_k$, $v\in V$ and $z\in\lfifteenH$:
\begin{enumerate}
\item[(i)] $F(a[0]v,z)=0$;
\item[(ii)] $F(a[-2]v,z)+\sum_{j\ge1}(2j-1)\,G_{2j}(z)F(a[2j-2]v,z)=0$;
\item[(iii)] $F(L[-2]v,z)=\frac1{2\pi i}\frac{\partial}{\partial z}F(v,z)+\sum_{j\ge1}G_{2j}(z)F(L[2j-2]v,z)$.
\end{enumerate}
The \emph{block space} $\lfifteenB$ is the vector space of such $F$, regarded as functions on $V\times\lfifteenH$.
\end{definition}
\noindent Files \leanpath{VOA/\allowbreak Modularity/\allowbreak QEvalHolomorphic.lean} (Chapter~\ref{ch:L14}),
\leanpath{VOA/\allowbreak Modularity/\allowbreak AnalyticBlock.lean}, \leanpath{VOA/\allowbreak Modularity/\allowbreak BlockSpace.lean},
\leanpath{VOA/\allowbreak Modularity/\allowbreak BlockRank.lean}.
\begin{leancode}
def imGT (c : ℝ) : Set ℂ := {z : ℂ | c < z.im}

structure IsAnalyticBlock (W : VertexOperatorAlgebra 𝒱) (F : 𝒱 → ℂ → ℂ) : Prop where
  holo : ∀ v : 𝒱, DifferentiableOn ℂ (F v) (imGT 0)
  addv : ∀ (v w : 𝒱) (z : ℂ), F (v + w) z = F v z + F w z
  smulv : ∀ (c : ℂ) (v : 𝒱) (z : ℂ), F (c • v) z = c * F v z
  vanish : ∀ (ka : ℕ), 1 ≤ ka → ∀ a ∈ W.grading (ka : ℤ), ∀ v : 𝒱, F (vSqZ W a ka 0 v) = 0
  negTwo : ∀ (ka : ℕ), 1 ≤ ka → ∀ a ∈ W.grading (ka : ℤ), ∀ (v : 𝒱), ∀ z ∈ imGT 0,
    F (vSqZ W a ka (-2) v) z
      + ∑ᶠ m : ℕ, ((2 * (m + 1) - 1 : ℕ) : ℂ)
          * (zhuGC (m + 1) z * F (vSqZ W a ka (2 * (m + 1) - 2 : ℤ) v) z) = 0
  omega : ∀ (v : 𝒱), ∀ z ∈ imGT 0,
    F (Lsq W (-2) v) z
      = (2 * (Real.pi : ℂ) * Complex.I)⁻¹ * _root_.deriv (F v) z
        + ∑ᶠ m : ℕ, zhuGC (m + 1) z * F (Lsq W (2 * (m + 1) - 2 : ℤ) v) z

def analyticBlocks (W : VertexOperatorAlgebra 𝒱) : Submodule ℂ (𝒱 → ℂ → ℂ) where
  carrier := {F : 𝒱 → ℂ → ℂ | IsAnalyticBlock W F}

noncomputable def restrictImGT : (𝒱 → ℂ → ℂ) →ₗ[ℂ] (𝒱 → ℂ → ℂ) where
  toFun F := fun v => Set.indicator (imGT 0) (F v)
\end{leancode}

\begin{dictionary}
\textbf{Decoding.} A block is a function \lean{F : 𝒱 → ℂ → ℂ}, defined on all of $\C$ in the second variable;
\lean{imGT 0} is $\lfifteenH$ as a subset of $\C$. \lean{DifferentiableOn ℂ (F v) (imGT 0)} is holomorphy on $\lfifteenH$.
\lean{_root_.deriv} is the complex derivative $\partial_z$; since $\lfifteenH$ is open, it only sees values on
$\lfifteenH$. Fields (ii) and (iii) are imposed only at $z\in\lfifteenH$. Field (i) is stated as equality of functions on
all of $\C$, i.e.\ also off $\lfifteenH$; this is harmless (the trace blocks are extended by $0$ off $\lfifteenH$).
\lean{Set.indicator (imGT 0) f} is $f$ on $\lfifteenH$ and $0$ elsewhere.

\textbf{The block space is $\mathrm{restrictImGT}(\lfifteenB)$, not $\lfifteenB$.} Since nothing constrains a block off
$\lfifteenH$ except linearity and (i), the submodule \lean{analyticBlocks W} is infinite-dimensional for a trivial reason
(arbitrary junk values off the half-plane). All dimension statements are therefore about its image under
\lean{restrictImGT}, which identifies blocks agreeing on $\lfifteenH$. That image is Zhu's space $\lfifteenB$.

\textbf{Fidelity with Zhu.} Zhu's Definition~5.1.1 imposes the conditions on the $\C$-span of the functions, uses
$q\frac{d}{dq}=\frac1{2\pi i}\partial_\tau$ in (iii), and requires holomorphy; all three are matched. Some later formulations (Dong--Li--Mason) also prescribe the shape of the
$q$-expansion at the cusp; the library's $\lfifteenB$ imposes no condition at the cusp, which can only make $\lfifteenB$
larger, so the upper bound $\dim\lfifteenB\le\#\Phi(V)$ proved below is at least as strong as the literature's.
\end{dictionary}

\begin{physics}
$\lfifteenB$ is the space of torus one-point conformal blocks of the chiral algebra: all consistent assignments of a
``one-point function'' to fields, compatible with the torus Ward identities. In a rational CFT it is spanned by the
characters of the sectors (with insertions), and its dimension is the number of primaries. The statement $\dim\lfifteenB=\#\Phi$
is the genus-one case of the Verlinde dimension formula.
\end{physics}

\begin{status}
Definitions; the closure of \lean{IsAnalyticBlock} under sums and scalars is proved (\lean{IsAnalyticBlock.add},
\lean{IsAnalyticBlock.smul}).
\end{status}

\begin{theorem}[The trace functions are blocks; Zhu 1996, Thm.~5.3.1]\label{l15:thm:traceblock}
Let $M$ have conformal weight $h$ and subexponential one-point functions. Then $(v,z)\mapsto S_M(v,z)$ (the trace
function, extended by $0$ off $\lfifteenH$) is an analytic genus-one block. In particular this holds for every ordinary
module with a conformal weight over a $C_2$-cofinite $V$.
\end{theorem}
\noindent Files \leanpath{VOA/\allowbreak Modularity/\allowbreak AnalyticBlock.lean},
\leanpath{VOA/\allowbreak Modularity/\allowbreak RhoS.lean} (Chapter~\ref{ch:L14}),
\leanpath{VOA/\allowbreak Modularity/\allowbreak SlashAmbientObstruction.lean}.
\begin{leancode}
noncomputable def anEval (S : 𝒱 →ₗ[ℂ] FormalQSeries) (h : ℂ) (v : 𝒱) : ℂ → ℂ :=
  Set.indicator (imGT 0) (qEvalC (S v) h)

theorem torusOnePoint_isAnalyticBlock (N : OrdVModule W M) (h : ℂ)
    (hsupp : OnePointSupported N h) (hsub : OnePointSubexp N h) :
    IsAnalyticBlock W (anEval (torusOnePointHom N) h) :=

noncomputable def traceBlock (Mfam : ι → IrrModule W) (h : ι → ℂ) (i : ι) : 𝒱 → ℂ → ℂ :=
  anEval (torusOnePointHom (Mfam i).isMod) (h i - W.centralCharge / 24)

theorem traceBlock_coe (i : ι) (v : 𝒱) (τ : UpperHalfPlane) :
    traceBlock Mfam h i v (τ : ℂ) = traceFun Mfam h i v τ := by

theorem traceBlock_isAnalyticBlock
    (hcw : ∀ i, HasConformalWeight (Mfam i).isMod (h i))
    (hgr : ∀ i, OnePointSubexp (Mfam i).isMod (h i - W.centralCharge / 24)) (i : ι) :
    IsAnalyticBlock W (traceBlock Mfam h i) :=
\end{leancode}

\begin{dictionary}
\lean{qEvalC χ h z} is \lean{qEvalAt} in the plain complex variable $z$ (Chapter~\ref{ch:L14}). \lean{traceBlock_coe}
says that on $\lfifteenH$ the block \emph{is} the trace function \lean{traceFun} of Definition~\ref{l15:def:traceFun}. The
support hypothesis is free from \lean{HasConformalWeight} (\lean{onePointSupported_}\allowbreak\lean{of_hasConformalWeight}), and the growth
hypothesis from $C_2$ (Theorem~\ref{l15:thm:growth}). Same as Zhu's 5.3.1.
\end{dictionary}

\begin{proof}[Proof on paper]
The formal identities of Theorem~\ref{l15:thm:formalblock} are transported to functions: under subexponential growth
the evaluation map sends $q\frac d{dq}$ to $\frac1{2\pi i}\partial_z$ (termwise differentiation of a locally uniformly
convergent series), sends the convolution $G_{2j}\star\chi$ to the product $G_{2j}(z)\cdot\chi(z)$ (Cauchy product of
two absolutely convergent series), and is linear. Holomorphy is locally uniform convergence.
\end{proof}

\begin{physics}
The characters with insertions solve the torus Ward identities: they are genuine conformal blocks. This is the ``easy''
inclusion $\mathrm{span}\{\text{characters}\}\subseteq\lfifteenB$.
\end{physics}

\begin{status}
\lean{VOA.torusOnePoint_}\allowbreak\lean{isAnalyticBlock}: \texttt{propext, Classical.choice, Quot.sound}; conditional on the displayed
support and growth data, both of which are theorems under \lean{HasConformalWeight} and \lean{C2Cofinite}.
\end{status}

\section{Finiteness of the block space and Zhu's count}\label{l15:sec:finite}

\subsection{Zhu's Theorem 5.3.2: \texorpdfstring{$\dim\lfifteenB<\infty$}{dim B finite}}

\begin{theorem}[Zhu 1996, Thm.~5.3.2 in the form proved]\label{l15:thm:532}
Let $V$ be $C_2$-cofinite and let $x_1,\dots,x_g$ be homogeneous vectors spanning $V$ modulo $C_2(V)+\C\lfifteenvac$. Then
for any $z_0\in\lfifteenH$ the map $F\mapsto(F(x_1,z_0),\dots,F(x_g,z_0),F(\lfifteenvac,z_0))$ is injective on $\lfifteenB$.
Hence $\dim\lfifteenB\le g+1<\infty$.
\end{theorem}
\noindent Files \leanpath{VOA/\allowbreak Modularity/\allowbreak C2Extraction.lean},
\leanpath{VOA/\allowbreak Modularity/\allowbreak BlockFiniteDim.lean}.
\begin{leancode}
theorem exists_homog_gen_of_C2Cofinite {𝒱 : Type u} [AddCommGroup 𝒱] [Module ℂ 𝒱]
    (W : VertexOperatorAlgebra 𝒱) (hC2 : C2Cofinite W) :
    ∃ (g : ℕ) (gen : Fin g → 𝒱) (wt : Fin g → ℕ),
      (∀ i, 1 ≤ wt i) ∧ (∀ i, gen i ∈ W.grading (wt i : ℤ)) ∧
      (⊤ : Submodule ℂ 𝒱) = Submodule.span ℂ (Set.range gen ∪ {W.vacuum}) ⊔ O₂ W := by

theorem exists_finrank_analyticBlocks_le (hC2 : C2Cofinite W) {z₀ : ℂ} (hz₀ : z₀ ∈ imGT 0) :
    ∃ g : ℕ, FiniteDimensional ℂ (Submodule.map restrictImGT (analyticBlocks W)) ∧
      Module.finrank ℂ (Submodule.map restrictImGT (analyticBlocks W)) ≤ g + 1 := by
\end{leancode}

\begin{dictionary}
\lean{O₂ W} is $C_2(V)=\mathrm{span}\{a_{(-2)}b\}$ (Definition~\ref{l12:def:c2}); \lean{⊔} is the sum of subspaces;
\lean{(⊤ : Submodule ℂ 𝒱) = …} says the right side is all of $V$. \lean{Module.finrank} is the dimension. Since
\lean{finrank} is $0$ on infinite-dimensional spaces, a bound on \lean{finrank} alone would say nothing; the conjunct
\lean{FiniteDimensional} is what makes the bound meaningful, and it is present. The pointwise injectivity is the
internal lemma \lean{IsAnalyticBlock.eq_}\allowbreak\lean{zero_of_}\allowbreak\lean{vanish_at}.
\end{dictionary}

\begin{proof}[Proof on paper]
Suppose $F\in\lfifteenB$ vanishes at $z_0$ on $x_1,\dots,x_g,\lfifteenvac$. \emph{Step 1: $F(\cdot,z_0)=0$ on all of $V$.}
By induction on the degree: modulo $C_2(V)$, $V_n$ is spanned by the $x_i$, $\lfifteenvac$, and elements $a_{(-2)}b$; by
\eqref{l15:eq:sqmode}, $a[-2]b=a_{(-2)}b+(\text{lower degree})$, and condition (ii) at the single point $z_0$ gives
$F(a[-2]b,z_0)=-\sum_j(2j-1)G_{2j}(z_0)F(a[2j-2]b,z_0)$, where all vectors on the right have lower degree. \emph{Step 2: all
derivatives vanish at $z_0$.} Condition (iii) gives $\partial_zF(v,z)=2\pi i\bigl(F(L[-2]v,z)-\sum_jG_{2j}(z)F(L[2j-2]v,z)\bigr)$:
the derivative of $F(v,\cdot)$ is a finite combination, with holomorphic coefficients, of the functions $F(w,\cdot)$. Since
$F(w,z_0)=0$ for all $w$ by step 1, $\partial_zF(v,z_0)=0$ for all $v$; differentiating the same identity, by induction all
$\partial_z^kF(v,z_0)=0$. \emph{Step 3.} $F(v,\cdot)$ is holomorphic on the connected open set $\lfifteenH$ with all
derivatives zero at $z_0$, hence identically zero (identity theorem).
\end{proof}

\begin{physics}
A torus conformal block is determined by finitely many numbers: its values on the finitely many ``generating fields''
at one value of the modulus. Equivalently, the genus-one blocks form a vector bundle of finite rank over the moduli space,
with a flat connection given by the stress-tensor Ward identity (iii). Finite rank is the hallmark of a rational (or at
least $C_2$-cofinite) theory.
\end{physics}

\begin{status}
\lean{VOA.exists_}\allowbreak\lean{finrank_}\allowbreak\lean{analyticBlocks_le}: \texttt{propext, Classical.choice, Quot.sound}; conditional only on
\lean{C2Cofinite}.
\end{status}

\subsection{\texorpdfstring{$A(V)$}{A(V)} is finite-dimensional}

\begin{theorem}[Zhu; Abe--Buhl--Dong]\label{l15:thm:Afin}
If $V$ is $C_2$-cofinite then $\dim A(V)<\infty$.
\end{theorem}
\noindent File \leanpath{VOA/\allowbreak Modularity/\allowbreak ZhuAlgebraFinite.lean}.
\begin{leancode}
theorem finiteDimensional_A_of_C2Cofinite (hC2 : C2Cofinite W) :
    FiniteDimensional ℂ (A W) := by
\end{leancode}

\begin{dictionary}
\lean{A W} is Zhu's algebra $V/O(V)$ (Chapter~\ref{ch:L12}, Section~\ref{l12:sec:zhu}). Textbook statement.
\end{dictionary}

\begin{proof}[Proof on paper]
For homogeneous $a\in V_m$, $a\circ b=\sum_{i=0}^m\binom mia_{(i-2)}b\in O(V)$. Its $i=0$ term is $a_{(-2)}b$ and the others
have strictly lower degree. So modulo $O(V)$ every element of $C_2(V)$ is congruent to an element of lower degree, and by
induction the $C_2$-generators and $\lfifteenvac$ span $V$ modulo $O(V)$.
\end{proof}

\begin{physics}
$A(V)$ is the algebra of zero modes acting on ground states (primary states). Finite-dimensionality means finitely many
species of primary states, the first sign of rationality.
\end{physics}

\begin{status}
\texttt{propext, Classical.choice, Quot.sound}; conditional only on \lean{C2Cofinite}.
\end{status}

\subsection{Zhu's Theorem 5.3.3: \texorpdfstring{$\dim\lfifteenB\le\#\Phi(V)$}{dim B at most the number of modules}}

The idea of Zhu's count is to look at a block at the cusp $q=0$. Expanding $F(v,\tau)=q^{\lambda}\sum_n F_n(v)q^n$, the
leading coefficient $F_0$ is a linear functional on $V$ which, by (i)--(iii) at $q=0$, factors through $A(V)$ and is
\emph{symmetric}: $F_0(a*b)=F_0(b*a)$. Symmetric functionals on a finite-dimensional semisimple algebra
$A\cong\bigoplus_{i=1}^m\mathrm{Mat}_{n_i}(\C)$ form a space of dimension $m$ (the codimension of $[A,A]$), and each simple
$A(V)$-module is the bottom of an irreducible $V$-module, distinct for distinct simple modules. So $m\le\#\Phi(V)$. The
library carries this out with a ``relation space'' at the cusp and transports it to an interior point.

\begin{definition}[Relation spaces]\label{l15:def:blockrel}
For $z_0\in\C$, $R(z_0)\subseteq V$ is the span of all $a[0]v$ and of all
$a[-2]v+\sum_{j\ge2}(2j-1)G_{2j}(z_0)a[2j-2]v$ ($a\in V_k$, $k\ge1$). $R(\infty)$ is the same with $G_{2j}(z_0)$ replaced by
the constant term $-B_{2j}/(2j)!$ of $G_{2j}$, i.e.\ its value at the cusp. $[A,A]$ is the span of commutators.
\end{definition}
\noindent Files \leanpath{VOA/\allowbreak Modularity/\allowbreak BlockFiniteDim.lean}, \leanpath{VOA/\allowbreak Modularity/\allowbreak CuspLimit.lean},
\leanpath{VOA/\allowbreak Modularity/\allowbreak SemisimpleTrace.lean}.
\begin{leancode}
def blockRelSpan (z₀ : ℂ) : Submodule ℂ 𝒱 :=
  Submodule.span ℂ
    ({x : 𝒱 | ∃ (ka : ℕ) (a v : 𝒱), 1 ≤ ka ∧ a ∈ W.grading (ka : ℤ) ∧ x = vSqZ W a ka 0 v} ∪
     {x : 𝒱 | ∃ (ka K : ℕ) (a v : 𝒱), 1 ≤ ka ∧ a ∈ W.grading (ka : ℤ) ∧
        (∀ m : ℕ, K ≤ m → vSqZ W a ka (2 * (m + 1) - 2 : ℤ) v = 0) ∧
        x = vSqZ W a ka (-2) v + ∑ m ∈ Finset.range K,
              (((2 * (m + 2) - 1 : ℕ) : ℂ) * zhuGC (m + 2) z₀)
                • vSqZ W a ka (2 * (m + 2) - 2 : ℤ) v})

def blockRelSpanAtCusp : Submodule ℂ 𝒱 :=
  Submodule.span ℂ
    ({x : 𝒱 | ∃ (ka : ℕ) (a v : 𝒱), 1 ≤ ka ∧ a ∈ W.grading (ka : ℤ) ∧ x = vSqZ W a ka 0 v} ∪
     {x : 𝒱 | ∃ (ka K : ℕ) (a v : 𝒱), 1 ≤ ka ∧ a ∈ W.grading (ka : ℤ) ∧
        (∀ m : ℕ, K ≤ m → vSqZ W a ka (2 * (m + 1) - 2 : ℤ) v = 0) ∧
        x = vSqZ W a ka (-2) v + ∑ m ∈ Finset.range K,
              (((2 * (m + 2) - 1 : ℕ) : ℂ) * zhuG (m + 2) 0)
                • vSqZ W a ka (2 * (m + 2) - 2 : ℤ) v})

def commSpan : Submodule K A := Submodule.span K {x : A | ∃ a b : A, x = a * b - b * a}
\end{leancode}

\begin{dictionary}
In both relation spaces the second family is condition (ii) with the sum started at $j=m+2\ge2$ (the $j=1$ term is
$G_2\cdot a[0]v$, already in the first family), truncated at a $K$ beyond which $a[2j-2]v=0$. \lean{zhuGC (m + 2) z₀} is
$G_{2(m+2)}(z_0)$ and \lean{zhuG (m + 2) 0} is its constant term. Every block kills $R(z_0)$ at $z_0$
(\lean{IsAnalyticBlock.blockRelSpan_}\allowbreak\lean{le_kerAt}).
\end{dictionary}

\begin{physics}
$V/R(z_0)$ is the space of ``fields modulo torus Ward identities at modulus $z_0$''; its dual contains the space of
blocks evaluated at $z_0$. At the cusp ($q\to0$, the torus degenerating to a long cylinder) only the ground states
propagate, and the Ward identities reduce to the zero-mode algebra $A(V)$ modulo commutators: a block becomes a trace over
ground states.
\end{physics}

\begin{status}
Definitions.
\end{status}

\begin{theorem}[Zhu 1996, Thm.~5.3.3, upper bound]\label{l15:thm:533}
Let $V$ be $C_2$-cofinite with $A(V)$ semisimple and $\Phi(V)$ finite, assume every simple $A(V)$-module is the bottom of an
irreducible module, and assume the datum \lean{Pr}. Then
\begin{enumerate}
\item $\dim V/R(\infty)=\dim A(V)/[A(V),A(V)]\le\#\Phi(V)$;
\item $\lfifteenB$ is finite-dimensional and $\dim\lfifteenB\le\#\Phi(V)$.
\end{enumerate}
\end{theorem}
\noindent Files \leanpath{VOA/\allowbreak Modularity/\allowbreak CuspLimit.lean}, \leanpath{VOA/\allowbreak Modularity/\allowbreak ZhuCount.lean},
\leanpath{VOA/\allowbreak Modularity/\allowbreak ZhuCountC2.lean}, \leanpath{VOA/\allowbreak Modularity/\allowbreak ZhuBlockCount.lean}.
\begin{leancode}
theorem finrank_quot_blockRelSpanAtCusp_eq :
    letI := instRingA W
    letI := instAlgebraA W
    Module.finrank ℂ (𝒱 ⧸ blockRelSpanAtCusp W)
      = Module.finrank ℂ (A W ⧸ commSpan ℂ (A W)) := by

theorem finrank_quot_commSpan_le_card_Φ [Finite (Φ W)]
    (hfd : letI := instRingA W; FiniteDimensional ℂ (A W))
    (hss : letI := instRingA W; IsSemisimpleRing (A W))
    (hreal : SimpleAModsRealized W)
    (Pr : BottomProjection (parabolicRepOfAMod W (regularAMod W))) :
    letI := instRingA W; letI := instAlgebraA W;
    Module.finrank ℂ (A W ⧸ commSpan ℂ (A W)) ≤ Nat.card (Φ W) := by

theorem finiteDimensional_and_finrank_analyticBlocks_le_card_Φ [Finite (Φ W)]
    (hzero : VacuumGradeZero W) (hC2 : C2Cofinite W)
    (hss : letI := instRingA W; IsSemisimpleRing (A W))
    (hreal : SimpleAModsRealized W)
    (Pr : BottomProjection (parabolicRepOfAMod W (regularAMod W))) :
    FiniteDimensional ℂ (Submodule.map restrictImGT (analyticBlocks W)) ∧
      Module.finrank ℂ (Submodule.map restrictImGT (analyticBlocks W)) ≤ Nat.card (Φ W) := by
\end{leancode}

\begin{dictionary}
\textbf{Hypotheses.} \lean{[Finite (Φ W)]}: finitely many irreducible ordinary modules up to isomorphism.
\lean{letI := instRingA W} installs the ring structure of $A(V)$ (Zhu's product) so that \lean{IsSemisimpleRing (A W)}
(Mathlib: every module is semisimple, equivalently Artin--Wedderburn) can be stated. \lean{hreal : SimpleAModsRealized W}:
every simple $A(V)$-module is the bottom of an irreducible module (Chapter~\ref{ch:L13},
Definition~\ref{l13:def:simpleamod}); it is \emph{discharged} in the headline from $C_2$ and \lean{Pr}
(Section~\ref{l15:sec:gn}). \lean{Pr} is the datum of Chapter~\ref{ch:L13} (Definition~\ref{l13:def:bottomprojection}); it is
used through Zhu's bijection \lean{zhuBijective Pr} to make distinct simple $A(V)$-modules give non-isomorphic $V$-modules.
\lean{Nat.card (Φ W)} is the number of elements of the finite type $\Phi(V)$.

\textbf{Conclusion.} \lean{Submodule.map restrictImGT (analyticBlocks W)} is $\lfifteenB$ (dictionary of
Definition~\ref{l15:def:analyticblock}). Both finite-dimensionality and the bound are stated, so the \lean{finrank}
trap does not apply. The name \lean{finrank_}\allowbreak\lean{analyticBlocks_}\allowbreak\lean{le_card_Φ} denotes the bound alone; read alone it would be
vacuous on an infinite-dimensional space, which is why we quote the conjunction.

\textbf{Fidelity.} Zhu proves $\dim\lfifteenB=\#\Phi(V)$ (with the traces as a basis) for rational $C_2$-cofinite $V$. Here
only the \emph{upper bound} $\le$ is proved, together with finite-dimensionality. The reverse inequality follows on paper
from the linear independence of $\#\Phi(V)$ trace functions (Section~\ref{l15:sec:independence}), so the equality is a
consequence of the chapter's results under the headline's hypotheses; but it is \emph{not} a Lean theorem, and the
release must not announce ``$\dim\lfifteenB=\#\Phi(V)$''. When this result is cited, cite the conjunction
\lean{finiteDimensional_}\allowbreak\lean{and_finrank_}\allowbreak\lean{analyticBlocks_}\allowbreak\lean{le_card_Φ}, not
\lean{finrank_}\allowbreak\lean{analyticBlocks_}\allowbreak\lean{le_card_Φ}.

\textbf{The five displayed hypotheses are all dischargeable.} \lean{hzero} holds for every VOA of the library
(Section~\ref{l15:sec:convergence}); \lean{hreal} is \lean{simpleAModsRealized_of_C2Cofinite}
(Theorem~\ref{l15:thm:realized}); \lean{Pr} is \lean{voaPBWBottomProjection}; \lean{[Finite (Φ W)]} is
\lean{finite_Φ_of_C2Cofinite_pbw} (Theorem~\ref{l15:thm:finitephi}). So in substance the bound holds under $C_2$ and
\lean{hss} alone; the theorem as stated still displays all five.
\end{dictionary}

\begin{proof}[Proof on paper]
(1) A careful computation at $q=0$ shows that $V/R(\infty)\cong A(V)/[A(V),A(V)]$: modulo $R(\infty)$, $O(V)$ dies (this uses
condition (ii) at the cusp, whose constants $-B_{2j}/(2j)!$ reproduce the binomial coefficients of $a\circ b$), and
$a[0]b\equiv a*b-b*a$ (Zhu's commutator lemma). For a finite-dimensional semisimple $\C$-algebra
$A\cong\bigoplus_{i=1}^m\mathrm{Mat}_{n_i}(\C)$, $[A,A]=\bigoplus\mathfrak{sl}_{n_i}$ has codimension $m$, the number of simple
modules (Mathlib's Artin--Wedderburn over an algebraically closed field). Each simple module is realised as a bottom
(\lean{hreal}) and by the uniqueness half of Zhu's correspondence (\lean{zhuBijective Pr}) non-isomorphic simple modules
give non-isomorphic $V$-modules; so $m\le\#\Phi(V)$.

(2) From (1), finitely many vectors $T$, $\#T\le\#\Phi(V)$, span $V$ modulo $R(\infty)$. \emph{Transport to the interior.}
The relation module over the ring $\C[E_4,E_6]$ of modular forms (all $G_{2j}$, $j\ge2$, lie in it) is finitely generated
(Zhu's Lemma~4.4.1); evaluation at the cusp is a ring homomorphism $\C[E_4,E_6]\to\C$, and Nakayama's lemma at its kernel
lifts the spanning family $T$ from the cusp to a neighbourhood, hence to some interior point $z_0\in\lfifteenH$:
$V=\mathrm{span}(T)+R(z_0)$. \emph{Dual descent.} A block $F$ kills $R(z_0)$ at $z_0$, so $F(\cdot,z_0)$ is determined by its
values on $T$; by steps 1--3 of Theorem~\ref{l15:thm:532} (run with $T$ in place of the generators) the evaluation
$F\mapsto(F(t,z_0))_{t\in T}$ is injective on $\lfifteenB$. So $\dim\lfifteenB\le\#T\le\#\Phi(V)$.
\end{proof}

\begin{physics}
``The number of independent torus blocks is at most the number of sectors.'' At the cusp each block becomes a trace
functional on the primary states; a semisimple zero-mode algebra has exactly one trace per irreducible representation;
each irreducible representation of $A(V)$ is the ground-state multiplet of exactly one sector.
\end{physics}

\begin{status}
\lean{VOA.finrank_}\allowbreak\lean{quot_blockRelSpanAtCusp_}\allowbreak\lean{le_card_Φ_of_C2}, \lean{VOA.finiteDimensional_}\allowbreak\lean{and_finrank_}\allowbreak\lean{analyticBlocks_}\allowbreak\lean{le_card_Φ}:
\texttt{propext, Classical.choice, Quot.sound} (the conjunction, and the bound alone, re-checked at an earlier commit (not re-run for this edition)).
Conditional on \lean{C2Cofinite}, \lean{IsSemisimpleRing (A W)}, \lean{Finite (Φ W)}, \lean{SimpleAModsRealized},
\lean{Pr} (and the redundant \lean{hzero}); all but the first two are dischargeable (dictionary). No \lean{VOA.pbw}. An
upper bound only.
\end{status}

\section[Existence of modules: Gaberdiel--Neitzke]{Existence of modules: Gaberdiel--Neitzke and the realisation of simple \texorpdfstring{$A(V)$}{A(V)}-modules}\label{l15:sec:gn}

Zhu's count needs every simple $A(V)$-module to occur as the bottom of an irreducible ordinary $V$-module. Chapter~\ref{ch:L13}
constructs, for a simple $A(V)$-module $U$, Zhu's (Dong--Jiang's) irreducible module $L(U)$ with bottom $U$; what it does not
supply is that the weight spaces of $L(U)$ are \emph{finite-dimensional}, which is needed for $L(U)$ to be an ordinary module.
That finiteness comes from a spanning theorem of Gaberdiel and Neitzke.

\subsection{Strictly decreasing spanning}

\begin{definition}[Mode words; strict spanning; $C_m(V)$]\label{l15:def:strict}
Fix homogeneous generators $x_1,\dots,x_g$ of positive weights. A \emph{mode word} is
$x^{i_1}_{(-p_1)}x^{i_2}_{(-p_2)}\cdots x^{i_k}_{(-p_k)}\lfifteenvac$ (outermost mode first), with $p_j\ge1$; it is
\emph{strictly decreasing} if $p_1>p_2>\cdots>p_k\ge1$. $V$ has \emph{strict spanning} if for some such generators the
strictly decreasing words span $V$. $C_m(V)=\mathrm{span}\{u_{(-m)}v:u,v\in V\}$.
\end{definition}
\noindent Files \leanpath{VOA/\allowbreak Modularity/\allowbreak StrictSpanning.lean}, \leanpath{VOA/\allowbreak Modularity/\allowbreak AnFiniteDim.lean}.
\begin{leancode}
noncomputable def modeWord (W : VertexOperatorAlgebra 𝒱) {g : ℕ} (gen : Fin g → 𝒱) :
    List (Fin g × ℕ) → 𝒱 → 𝒱
  | [], x => x
  | p :: l, x => W.smul (gen p.1) (-(p.2 : ℤ)) (modeWord W gen l x)

def StrictDec {g : ℕ} (l : List (Fin g × ℕ)) : Prop :=
  List.Pairwise (fun a b : Fin g × ℕ => b.2 < a.2) l ∧ ∀ p ∈ l, 1 ≤ p.2

structure StrictSpanning (W : VertexOperatorAlgebra 𝒱) where
  g : ℕ
  gen : Fin g → 𝒱
  wt : Fin g → ℕ
  gen_mem : ∀ i, gen i ∈ W.grading (wt i : ℤ)
  spans : (⊤ : Submodule ℂ 𝒱) = Submodule.span ℂ
    {x | ∃ l : List (Fin g × ℕ), StrictDec l ∧ x = modeWord W gen l W.vacuum}

def CGen (W : VertexOperatorAlgebra 𝒱) (M : ℤ) : Submodule ℂ 𝒱 :=
  Submodule.span ℂ {x | ∃ u v : 𝒱, x = W.smul u M v}
\end{leancode}

\begin{dictionary}
\lean{W.smul u n v} is the product $u_{(n)}v$ (Chapter~\ref{ch:01}). A word is a list of pairs (generator index, $p$),
head outermost, so \lean{modeWord W gen l 𝟙} is the displayed word. \lean{List.Pairwise (fun a b => b.2 < a.2) l} says
that every later index is smaller than every earlier one: strict decrease. \lean{StrictSpanning W} is a structure in
\lean{Type} (data, not a proposition): it records the generators. \lean{CGen W (-m)} is $C_m(V)$; \lean{CGen W (-2)} is
\lean{O₂ W}.

\textbf{Convention.} The library's words use \emph{vertex-algebra} modes $x_{(-p)}$, $p\ge1$, and require the index $p$
to decrease strictly across \emph{all} generators. Whether Gaberdiel--Neitzke's Proposition~8 is stated in these modes or
in \emph{weight} modes $x_{-n}=x_{(-n+\wt x-1)}$ ordered by $n$ has \textbf{not been checked against the paper} (the release
audit flags this as open). If GN use weight modes, the two statements differ whenever the generators have different
weights, and neither implies the other. This is a question of labelling, not of correctness: the library \emph{proves}
its own form from $C_2$-cofiniteness (below), and uses it only through the consequence
\lean{cofinite_}\allowbreak\lean{of_strictSpanning}, which is the content of GN's Theorem~11. When citing, say
``strict-decrease spanning in vertex-algebra mode indices $x_{(-p)}$, $p\ge1$'', not ``GN Proposition~8'' without
qualification.
\end{dictionary}

\begin{physics}
A PBW-type basis of the vacuum sector built from finitely many generating fields, with each mode used at most once: an
``exclusion principle'' for the modes. For the free fermion (not a VOA in this library's sense, but the paradigm) this is
literally the Pauli principle; for $C_2$-cofinite theories it is a consequence of the null vectors that make $V/C_2(V)$
finite, and it is the algebraic reason for Cardy-type growth of the spectrum.
\end{physics}

\begin{status}
Definitions.
\end{status}

\begin{theorem}[Strict spanning; after Gaberdiel--Neitzke 2003, Prop.~8 and Thm.~11]\label{l15:thm:gn}
Let $V$ be $C_2$-cofinite and $x_1,\dots,x_g$ homogeneous of positive weight spanning $V$ modulo $C_2(V)+\C\lfifteenvac$.
Then the mode words in the $x_i$ whose vertex-algebra mode indices $p$ (modes $x_{(-p)}$, $p\ge1$) strictly decrease
span $V$. Consequently, for every $m\ge1$, $V/C_m(V)$ is
finite-dimensional, spanned by the finitely many strictly decreasing words with all $p_j<m$.
\end{theorem}
\noindent Files \leanpath{VOA/\allowbreak Modularity/\allowbreak GNProp8Closure.lean},
\leanpath{VOA/\allowbreak Modularity/\allowbreak StrictSpanning.lean}.
\begin{leancode}
theorem gnProp8_GNReach (W : VertexOperatorAlgebra 𝒱)
    {g : ℕ} (gen : Fin g → 𝒱) (wt : Fin g → ℕ)
    (hgen : ∀ i, gen i ∈ W.grading (wt i : ℤ))
    (hwt : ∀ i, 1 ≤ wt i)
    (hdec : (⊤ : Submodule ℂ 𝒱) =
      Submodule.span ℂ (Set.range gen ∪ {W.vacuum}) ⊔ O₂ W) :
    ∀ G N, GNReach W gen wt G N :=

noncomputable def strictSpanning_of_C2Cofinite (W : VertexOperatorAlgebra 𝒱)
    (hC2 : C2Cofinite W) : StrictSpanning W :=
  Classical.choice (nonempty_strictSpanning_of_C2Cofinite W hC2)

theorem cofinite_of_strictSpanning {W : VertexOperatorAlgebra 𝒱} (d : StrictSpanning W)
    (m : ℕ) (hm : 1 ≤ m) :
    ∃ S : Finset 𝒱, (∀ s ∈ (S : Set 𝒱), IsHomog W s) ∧
      (⊤ : Submodule ℂ 𝒱) = Submodule.span ℂ (S : Set 𝒱) ⊔ CGen W (-(m : ℤ)) := by
\end{leancode}

\begin{dictionary}
\lean{GNReach W gen wt G N} is the induction statement of GN's proof: the part of $V$ of total generator weight $\le G$ is
spanned by ordered words with no repeated mode index at or below $N$; Proposition~8 is its limit, packaged into a
\lean{StrictSpanning} by \lean{strictSpanning_}\allowbreak\lean{of_GNReachBounded}. \lean{strictSpanning_}\allowbreak\lean{of_C2Cofinite} is a \lean{def}
(it produces data) defined by \lean{Classical.choice} from the proposition \lean{nonempty_}\allowbreak\lean{strictSpanning_}\allowbreak\lean{of_C2Cofinite};
this is the standard way to extract data from an existence proof and adds no assumption. \lean{IsHomog W s}: $s$ is
homogeneous. The conclusion of \lean{cofinite_}\allowbreak\lean{of_strictSpanning} is ``$V=\mathrm{span}(S)+C_m(V)$ with $S$ finite'', i.e.\
$\dim V/C_m(V)<\infty$.
\end{dictionary}

\begin{proof}[Proof on paper]
\emph{Prop.~8.} Start from words with weakly ordered modes (Li's filtration gives spanning by such words modulo lower
filtration). A word with a repeated index $p$ contains $x_{(-p)}y_{(-p)}$; using the commutator formula, the iterate
formula, and the fact that $x_{(-1)}y$ is a combination of generators modulo $C_2(V)$, one rewrites
$x_{(-p)}y_{(-p)}w$ as a combination of words of lower total weight or with fewer repetitions (GN's Lemma~7). A
lexicographic induction on (total weight, repetition bound) concludes. \emph{Thm.~11.} In a strictly decreasing word the
outermost index is the largest. If it is $\ge m$, the word is $x_{(-p)}w$ with $p\ge m$, which lies in $C_p(V)\subseteq
C_m(V)$ (because $u_{(-p-1)}v=\frac1p(L_{-1}u)_{(-p)}v$). Otherwise all indices are $<m$; strictly decreasing lists of
positive integers $<m$ have length $<m$, so there are finitely many such words.
\end{proof}

\begin{physics}
$C_2$-cofiniteness (``finitely many fields modulo derivatives and normal-ordered products'') propagates to all
$C_m$-cofiniteness: the higher-level ``zero-mode algebras'' $A_n(V)$ that control excited states are all finite-dimensional.
\end{physics}

\begin{status}
\lean{VOA.gnProp8_GNReach}, \lean{VOA.nonempty_}\allowbreak\lean{strictSpanning_}\allowbreak\lean{of_C2Cofinite}, \lean{VOA.cofinite_}\allowbreak\lean{of_strictSpanning}:
\texttt{propext, Classical.choice, Quot.sound} (\lean{nonempty_}\allowbreak\lean{strictSpanning_}\allowbreak\lean{of_C2Cofinite}
re-checked at an earlier commit (not re-run for this edition)). Conditional only on \lean{C2Cofinite}. The Prop-level results are
\lean{nonempty_}\allowbreak\lean{strictSpanning_}\allowbreak\lean{of_C2Cofinite} and \lean{gnProp8_GNReach};
\lean{strictSpanning_}\allowbreak\lean{of_C2Cofinite} is a \lean{def} producing data from them. The match of the mode
convention with Gaberdiel--Neitzke's paper is unverified (dictionary).
\end{status}

\subsection{Every simple \texorpdfstring{$A(V)$}{A(V)}-module is realised}

\begin{theorem}[Zhu's existence theorem under $C_2$; Zhu 1996, Thm.~2.2.1, with Dong--Li--Mason]\label{l15:thm:realized}
Let $V$ be $C_2$-cofinite and assume the datum \lean{Pr}. Then every simple $A(V)$-module is the bottom of an
irreducible ordinary $V$-module.
\end{theorem}
\noindent Files \leanpath{VOA/\allowbreak Modularity/\allowbreak ZhuModularityGN.lean},
\leanpath{VOA/\allowbreak Modularity/\allowbreak GNProp8Closure.lean}.
\begin{leancode}
theorem simpleAModsRealized_of_strictSpanning (hC2 : C2Cofinite W)
    (Pr : BottomProjection (parabolicRepOfAMod W (regularAMod W)))
    (d : StrictSpanning W) : SimpleAModsRealized W :=

theorem simpleAModsRealized_of_C2Cofinite {W : VertexOperatorAlgebra 𝒱}
    (hC2 : C2Cofinite W)
    (Pr : BottomProjection (parabolicRepOfAMod W (regularAMod W))) :
    SimpleAModsRealized W :=
\end{leancode}

\begin{dictionary}
\lean{SimpleAModsRealized W} (Chapter~\ref{ch:L13}, Definition~\ref{l13:def:simpleamod}): for every $A(V)$-module $U$ with
\lean{IsSimpleAMod U}, there is an \lean{IrrRealization W U}: an irreducible graded module with finite-dimensional weight
spaces and an $A(V)$-isomorphism of its bottom with $U$. Both hypotheses are explicit: \lean{hC2} and \lean{Pr}
(there is no \lean{_voaPBW} wrapper for this theorem; one instantiates \lean{Pr} with \lean{voaPBWBottomProjection}, as
Chapter~\ref{ch:L13}, Theorem~\ref{l13:thm:zhucorr} does). \textbf{Comparison with Zhu.} Zhu's existence theorem (part~(2) of
Theorem~\ref{l13:thm:zhu-lit}) needs \emph{no} $C_2$ hypothesis: it produces an irreducible \emph{admissible} module $L(U)$
for every simple $U$. The library's statement asks for more in the conclusion (finite-dimensional weight spaces, i.e.\ an
ordinary module) and pays for it with $C_2$-cofiniteness, which is what makes the weight spaces finite-dimensional. Read
correctly: ``every simple $A(V)$-module is the bottom of an irreducible \emph{ordinary} module, for $C_2$-cofinite $V$''.
\end{dictionary}

\begin{proof}[Proof on paper]
Let $U$ be simple. $A(V)$ is finite-dimensional (Theorem~\ref{l15:thm:Afin}), so $U$ is finite-dimensional and $[\omega]$
acts on it by a scalar $h$ (it is central). Chapter~\ref{ch:L13} builds the irreducible module $L(U)=\bigoplus_nL(U)_{h+n}$
with $L(U)_h\cong U$; each $L(U)_{h+n}$ is a quotient of $A_{n,0}(V)\otimes_{A(V)}U$, where $A_{n,0}(V)$ is a Dong--Jiang
bimodule, and $\dim A_{n,0}(V)<\infty$ as soon as $\dim V/C_{n+2}(V)<\infty$
(\lean{finiteDimensional_}\allowbreak\lean{AnmFull_}\allowbreak\lean{of_cofinite}, \leanpath{Anm0FiniteDim.lean}). By Theorem~\ref{l15:thm:gn} the
latter holds for every $n$, so $L(U)$ has finite-dimensional weight spaces and is an ordinary module.
\end{proof}

\begin{physics}
Every primary multiplet allowed by the zero-mode algebra generates a sector with finitely many states at each energy
level. In a $C_2$-cofinite theory there are no ``infinite-multiplicity'' sectors.
\end{physics}

\begin{status}
\lean{VOA.simpleAModsRealized_}\allowbreak\lean{of_C2Cofinite}: \texttt{propext, Classical.choice, Quot.sound}; conditional on
\lean{C2Cofinite} and \lean{Pr}. With \lean{Pr} supplied by \lean{voaPBWBottomProjection} it is conditional on $C_2$ alone, with
the standard three axioms. At the snapshot \lean{Pr} was available only through the then-axiom \lean{VOA.pbw}
(Chapter~\ref{ch:L13}); that dependence was recorded as a library defect at the time and is now resolved, since
\lean{VOA.pbw} is a theorem.
\end{status}

\begin{warning}
\textbf{Stale docstring (documentation only).} The header of \leanpath{RealizationCofinite.lean} lists ``$C_{n+2}$-cofinite $\Leftarrow$
$C_2$-cofinite'' as ``Gaberdiel--Neitzke Prop.~8 --- NOT BUILT''. It is built (\leanpath{GNProp8Closure.lean}).
\end{warning}

\section{The action of \texorpdfstring{$SL_2(\mathbb{Z})$}{SL2(Z)} on blocks}\label{l15:sec:slash}

\subsection{The slash operator}

For $\gamma=\left(\begin{smallmatrix}a&b\\c&d\end{smallmatrix}\right)\in\lfifteenSL$ and $k\in\Z$, the weight-$k$ slash of a
function $f$ on $\lfifteenH$ is $(f|_k\gamma)(z)=(cz+d)^{-k}f(\gamma z)$. Zhu's $\lfifteenSL$-action on $\lfifteenB$ is
\begin{equation*}
  (F|\gamma)(v,z)=(cz+d)^{-n}F(v,\gamma z)\qquad (v\in V_{[n]}),
\end{equation*}
extended linearly in $v$ using the decomposition $V=\bigoplus_nV_{[n]}$ of Theorem~\ref{l15:thm:sqvir}.

\begin{definition}[Slash; the pointwise slash of a family]\label{l15:def:slash}
$\mathrm{denom}_\gamma(z)=cz+d$, $\gamma z=(az+b)/(cz+d)$, $(f|_k\gamma)(z)=f(\gamma z)\,\mathrm{denom}_\gamma(z)^{-k}$. For a
family $F$, the \emph{pointwise slash} is $(F|\gamma)(v,\cdot)=F(v,\cdot)|_{\wt[v]}\gamma$, where $\wt[v]$ is the $L[0]$-weight
of $v$ if $v$ is a nonzero $L[0]$-eigenvector and $0$ otherwise.
\end{definition}
\noindent Files \leanpath{VOA/\allowbreak Modularity/\allowbreak SerreDerivative.lean} (Chapter~\ref{ch:L14}),
\leanpath{VOA/\allowbreak Modularity/\allowbreak GradedBlock.lean}.
\begin{leancode}
noncomputable def denomC (γ : SL(2, ℤ)) (z : ℂ) : ℂ := (γ 1 0 : ℂ) * z + (γ 1 1 : ℂ)

noncomputable def mobiusC (γ : SL(2, ℤ)) (z : ℂ) : ℂ :=
  ((γ 0 0 : ℂ) * z + (γ 0 1 : ℂ)) / denomC γ z

noncomputable def slashC (k : ℤ) (γ : SL(2, ℤ)) (f : ℂ → ℂ) (z : ℂ) : ℂ :=
  f (mobiusC γ z) * (denomC γ z) ^ (-k)

noncomputable def sqWtOf (W : VertexOperatorAlgebra 𝒱) (v : 𝒱) : ℤ :=
  if h : ∃ n : ℤ, v ∈ sqWeightSpace W (n : ℂ) ∧ v ≠ 0 then h.choose else 0

noncomputable def blockSlash (W : VertexOperatorAlgebra 𝒱) (γ : SL(2, ℤ)) (F : 𝒱 → ℂ → ℂ) :
    𝒱 → ℂ → ℂ := fun v => slashC (sqWtOf W v) γ (F v)
\end{leancode}

\begin{dictionary}
\lean{SL(2, ℤ)} is Mathlib's $\lfifteenSL$; \lean{γ 1 0} is the entry $c$, \lean{γ 1 1} is $d$ (rows and columns indexed
from $0$). \lean{sqWtOf W v} is the $L[0]$-weight of $v$ when $v$ is a nonzero eigenvector (\lean{h.choose} picks it; it
is unique since eigenspaces for distinct eigenvalues meet in $0$) and the junk value $0$ otherwise. So
\lean{blockSlash W γ F} agrees with Zhu's action on every $L[0]$-homogeneous $v$, but is \emph{not} linear in $v$: on a sum
$v+w$ of vectors of different weights it uses the junk weight $0$.
\end{dictionary}

\begin{warning}
\textbf{The pointwise slash is not Zhu's action, and a theorem stated with it was vacuous (recorded in the library).}
Because \lean{blockSlash W γ F} is not additive in $v$, it is not a block, and an earlier reduction of Zhu's theorem
(\lean{zhu_modularity_}\allowbreak\lean{of_analytic_}\allowbreak\lean{finrank_le}, \leanpath{BlockFiniteDim.lean}, marked ``do not use'') assumed a rank bound
on an ambient space containing it. The library proves that this rank hypothesis is contradictory as soon as $V$ has
nonzero vectors of two different $L[0]$-weights $n\neq0$ and $0$ and some trace function is nonzero at $Sz$
(\lean{not_rank_}\allowbreak\lean{analyticSlashAmbient_le}, \leanpath{VOA/\allowbreak Modularity/\allowbreak SlashAmbientObstruction.lean}):
\begin{leancode}
theorem not_rank_analyticSlashAmbient_le
    (hcw : ∀ i, HasConformalWeight (Mfam i).isMod (h i))
    (hgr : ∀ i, OnePointSubexp (Mfam i).isMod (h i - W.centralCharge / 24))
    (hind : TraceIndependent Mfam h)
    {z : ℂ} (hz : z ∈ imGT 0) {n : ℤ} (hn : n ≠ 0) {v w : 𝒱}
    (hv : v ∈ sqWeightSpace W (n : ℂ)) (hv0 : v ≠ 0)
    (hw : w ∈ sqWeightSpace W ((0 : ℤ) : ℂ)) (hw0 : w ≠ 0)
    {i : ι} (hT : traceBlock Mfam h i v (mobiusC ModularGroup.S z) ≠ 0)
    (hD : denomC ModularGroup.S z ^ (-n) ≠ 1) :
    ¬ (Module.rank ℂ (analyticSlashAmbient Mfam h) ≤ (Fintype.card ι : Cardinal)) := by
\end{leancode}
The repair, used by the headline, is to restrict all functions to $L[0]$-\emph{homogeneous} vectors before comparing them
(next definition), where the pointwise slash is additive, and to use the genuinely linear slash \lean{sqBlockSlash}
(\leanpath{SqGradedBlockSlash.lean}) to show that the restricted slash is the restriction of a block.
\end{warning}

\begin{physics}
The modular group acts on the torus modulus; a field of conformal weight $n$ on the torus transforms like a modular form
of weight $n$. The pointwise slash attaches the weight to each field separately, which is correct for fields of definite
weight; for a superposition of fields of different weights one must transform each component separately, which is what
the linear extension does.
\end{physics}

\begin{status}
Definitions; \lean{VOA.not_}\allowbreak\lean{rank_analyticSlashAmbient_le}: \texttt{propext, Classical.choice, Quot.sound}.
\end{status}

\subsection{Homogeneous points and the slash of a trace function}

\begin{definition}[Restriction to homogeneous points]\label{l15:def:sqpt}
$\mathrm{Pt}(V)=\{(n,v):n\in\Z,\ v\in V_{[n]}\}$. For $F:V\times\C\to\C$, $\mathrm{res}(F)$ is the function on
$\mathrm{Pt}(V)\times\C$ given by $(n,v,z)\mapsto F(v,z)$ for $z\in\lfifteenH$ and $0$ otherwise.
\end{definition}
\noindent File \leanpath{VOA/\allowbreak Modularity/\allowbreak SlashAmbientObstruction.lean}.
\begin{leancode}
def SqPt (W : VertexOperatorAlgebra 𝒱) : Type u := Σ n : ℤ, (sqWeightSpace W (n : ℂ))

noncomputable def sqRestrict (W : VertexOperatorAlgebra 𝒱) :
    (𝒱 → ℂ → ℂ) →ₗ[ℂ] (SqPt W → ℂ → ℂ) where
  toFun F := fun p z => restrictImGT F (p.2 : 𝒱) z
\end{leancode}

\begin{dictionary}
\lean{Σ n : ℤ, (sqWeightSpace W (n : ℂ))} is the type of pairs $(n,v)$ with $v\in V_{[n]}$ (a dependent sum); it contains
$(n,0)$ for every $n$, harmlessly. \lean{sqRestrict W} is linear. Only integer weights are recorded; by
Theorem~\ref{l15:thm:sqvir} the weights are in $\N$, so nothing is lost. Since Zhu's law
(Theorem~\ref{l15:thm:lit}(2)) is only asserted for homogeneous $v$, comparing functions after \lean{sqRestrict} loses
nothing that the theorem needs.
\end{dictionary}

\begin{physics}
Only fields of definite torus weight have a definite modular transformation law; restricting to them is restricting to
the fields for which ``the one-point function is a modular form of weight $n$'' makes sense.
\end{physics}

\begin{status}
Definition.
\end{status}

\begin{theorem}[The $S$-slash of a trace function is a block; Zhu 1996, Thm.~5.3.2(ii)/5.3.3 step]\label{l15:thm:slashblock}
Let $V$ be $C_2$-cofinite and $M_i$ have conformal weight $h_i$. Then the restriction to homogeneous points of the pointwise
$S$-slash of the trace block of $M_i$ equals the restriction of some analytic block. More generally, for every analytic
block $F$ and every $\gamma$, $\mathrm{res}(F|\gamma)\in\mathrm{res}(\lfifteenB)$.
\end{theorem}
\noindent File \leanpath{VOA/\allowbreak Modularity/\allowbreak SqSlashBridge.lean}.
\begin{leancode}
theorem sqRestrict_blockSlash_mem_analyticBlocks {F : V → ℂ → ℂ}
    (hF : IsAnalyticBlock W F) (gamma : SL(2, ℤ)) :
    sqRestrict W (blockSlash W gamma F) ∈
      Submodule.map (sqRestrict W) (analyticBlocks W) := by

theorem traceBlock_S_sqRestrict_mem_analyticBlocks
    (hcw : ∀ i, HasConformalWeight (Mfam i).isMod (h i))
    (hzero : VacuumGradeZero W) (hC2 : C2Cofinite W) (i : iota) :
    sqRestrict W (blockSlash W ModularGroup.S (traceBlock Mfam h i)) ∈
      Submodule.map (sqRestrict W) (analyticBlocks W) := by
\end{leancode}

\begin{dictionary}
In this file the space is called \lean{V} and the index type \lean{iota}; they are the $\mathcal V$ and $\iota$ of the
other files. \lean{ModularGroup.S} is $S=\left(\begin{smallmatrix}0&-1\\1&0\end{smallmatrix}\right)$. \lean{Submodule.map
(sqRestrict W) (analyticBlocks W)} is $\mathrm{res}(\lfifteenB)$. The first theorem holds for every VOA and every block; the
second specialises it, using Theorems~\ref{l15:thm:growth} and~\ref{l15:thm:traceblock} to know that the trace function is a
block. Textbook content: ``$\lfifteenB$ is $\lfifteenSL$-stable'' (Zhu's Theorem~5.3.2 includes this), read on homogeneous
vectors.
\end{dictionary}

\begin{proof}[Proof on paper]
Define the linear slash $F\|\gamma$ by $(F\|\gamma)(v)=\sum_n F(v_n)|_n\gamma$, where $v=\sum_nv_n$, $v_n\in V_{[n]}$. On
homogeneous $v$ it agrees with the pointwise slash, so $\mathrm{res}(F|\gamma)=\mathrm{res}(F\|\gamma)$; it remains to show
$F\|\gamma\in\lfifteenB$. Linearity and holomorphy are clear. For (i) and (ii) one decomposes into square-bracket weights,
using that $a[m]$ shifts the $L[0]$-weight by $\wt[a]-m-1$ (proved here for the needed modes); then (ii) for $F\|\gamma$ follows
from (ii) for $F$ and the modularity $G_{2j}|_{2j}\gamma=G_{2j}$ for $j\ge2$ (the $j=1$ term is absent because of (i)).
For (iii) one uses the Serre derivative $\vartheta_kf=\frac1{2\pi i}f'+kG_2f=\frac1{2\pi i}f'-\frac{k}{12}E_2f$ (in Zhu's normalisation $G_2=-E_2/12$),
which satisfies $(\vartheta_kf)|_{k+2}\gamma=\vartheta_k(f|_k\gamma)$ (Chapter~\ref{ch:L14}): the non-modular part of
$G_2$ is exactly the defect of $\frac1{2\pi i}\frac{d}{dz}$ under slashing, and condition (iii) contains the term
$G_2(z)F(L[0]v,z)=nG_2(z)F(v,z)$, which turns the derivative into a Serre derivative. So (iii) is equivalent to an
$\lfifteenSL$-covariant identity.
\end{proof}

\begin{physics}
The torus Ward identities are invariant under large diffeomorphisms of the torus. The derivative with respect to the
modulus is not covariant by itself (the holomorphic anomaly of $E_2$), but the trace of the stress tensor supplies exactly
the compensating term: the Serre derivative is the covariant derivative on the line bundle of weight-$n$ modular forms.
\end{physics}

\begin{status}
\lean{VOA.traceBlock_}\allowbreak\lean{S_sqRestrict_}\allowbreak\lean{mem_analyticBlocks}: \texttt{propext, Classical.choice, Quot.sound}; conditional on
\lean{HasConformalWeight}, \lean{C2Cofinite} (and the redundant \lean{hzero}). This removed the last ``residual (e1)'' of the
project's earlier reductions.
\end{status}

\section{Linear independence of the trace functions}\label{l15:sec:independence}

\begin{definition}[Trace independence; fibre independence]\label{l15:def:indep}
The trace functions of a family $(M_i,h_i)_{i\in\iota}$ are \emph{independent} if $\sum_jc_jS_j(v,\tau)=0$ for all
$n\in\N$, $v\in V_{[n]}$, $\tau\in\lfifteenH$ forces all $c_j=0$. Write $t_i=e^{2\pi i(h_i-c/24)}$ for the $T$-eigenvalue. They
are \emph{fibre-independent} if the same holds for each class $\{j:t_j=t_i\}$ separately.
\end{definition}
\noindent Files \leanpath{VOA/\allowbreak Modularity/\allowbreak BlockSpanning.lean},
\leanpath{VOA/\allowbreak Modularity/\allowbreak TraceIndependence.lean}.
\begin{leancode}
def TraceIndependent (Mfam : ι → IrrModule W) (h : ι → ℂ) : Prop :=
  ∀ c : ι → ℂ, (∀ (n : ℕ) (v : 𝒱), v ∈ sqWeightSpace W (n : ℂ) → ∀ τ : UpperHalfPlane,
    ∑ j, c j * traceFun Mfam h j v τ = 0) → c = 0

def FibreIndependent (Mfam : ι → IrrModule W) (h : ι → ℂ) : Prop :=
  ∀ (i : ι) (c : ι → ℂ),
    (∀ (n : ℕ) (v : 𝒱), v ∈ sqWeightSpace W (n : ℂ) → ∀ τ : UpperHalfPlane,
      ∑ j ∈ Finset.univ.filter (fun j => tEigen W h j = tEigen W h i),
        c j * traceFun Mfam h j v τ = 0) →
    ∀ j ∈ Finset.univ.filter (fun j => tEigen W h j = tEigen W h i), c j = 0

theorem traceIndependent_of_fibreIndependent (hcw : ∀ i, HasConformalWeight (Mfam i).isMod (h i))
    (hfib : FibreIndependent Mfam h) : TraceIndependent Mfam h := by
\end{leancode}

\begin{dictionary}
\lean{tEigen W h j} is $e^{2\pi i(h_j-c/24)}$ (\leanpath{TEigenInjective.lean}). \lean{Finset.univ.filter P} is the finite
set of indices satisfying $P$. \lean{TraceIndependent} quantifies over homogeneous $v$ only, which is the form needed (and
weaker than independence over all $v$, hence a weaker hypothesis to prove and a sufficient one to use). Textbook notions.
\end{dictionary}

\begin{physics}
Distinct sectors have distinguishable torus one-point functions: no linear combination of the characters-with-insertions of
distinct sectors vanishes. The $T$-eigenvalue fibres are the sectors with the same spin $h-c/24 \bmod 1$.
\end{physics}

\begin{status}
\lean{VOA.traceIndependent_}\allowbreak\lean{of_fibreIndependent}: standard three axioms; no hypothesis on $V$.
\end{status}

\begin{theorem}[Zhu 1996, Thm.~5.3.3, independence]\label{l15:thm:indep}
Let $M_1,\dots,M_r$ have conformal weights $h_i$ and subexponential one-point functions, with simple bottoms
$(M_i)_{h_i}$, such that $(M_j)_{h_j}\not\cong(M_k)_{h_k}$ as $A(V)$-modules whenever $j\neq k$ and $h_j=h_k$. Then their trace
functions are independent.
\end{theorem}
\noindent File \leanpath{VOA/\allowbreak Modularity/\allowbreak FibreIndependence.lean}.
\begin{leancode}
theorem fibreIndependent_of_bottomSeparation (Mfam : ι → IrrModule W) (h : ι → ℂ)
    (hcw : ∀ i, HasConformalWeight (Mfam i).isMod (h i))
    (hgr : ∀ i, OnePointSubexp (Mfam i).isMod (h i - W.centralCharge / 24))
    (hsimp : ∀ j, IsSimpleAMod (lowestAMod (hcw j).2))
    (hsep : ∀ j k : ι, j ≠ k → h j = h k →
      IsEmpty (AModEquiv W (lowestAMod (hcw j).2) (lowestAMod (hcw k).2))) :
    FibreIndependent Mfam h := by
\end{leancode}

\begin{dictionary}
\lean{(hcw j).2} is the \lean{IsLowestWeight} half of \lean{HasConformalWeight}, and \lean{lowestAMod} of it is the bottom
$(M_j)_{h_j}$ as an $A(V)$-module (Chapter~\ref{ch:L12}). \lean{AModEquiv W U U'} is the type of $A(V)$-module isomorphisms
and \lean{IsEmpty} says there are none. So \lean{hsep} is ``non-isomorphic bottoms at equal weight'', \lean{hsimp} is
``simple bottoms''. Same as Zhu's hypotheses (which follow from ``pairwise non-isomorphic irreducible modules'' for rational
$V$).
\end{dictionary}

\begin{proof}[Proof on paper]
\emph{Reduction to fibres.} If $\sum_jc_jS_j=0$, apply $\tau\mapsto\tau+k$ for $k=0,1,2,\dots$: $S_j$ is multiplied by
$t_j^k$, so $\sum_jt_j^k(c_jS_j)=0$ for all $k$; a Vandermonde (Lagrange interpolation) argument isolates each
fibre $\{t_j=t\}$. \emph{Within a fibre} the $h_j$ differ by integers. Let $h_0$ be the smallest weight among the $j$
with $c_j\neq0$ and expand at the cusp: the combination is a single $q$-series vanishing on $\lfifteenH$, so its lowest
coefficient vanishes; that coefficient is $\sum_{j:h_j=h_0}c_j\tr_{(M_j)_{h_0}}o(v)$ for all $v$, i.e.\ a vanishing combination
of the characters $v\mapsto\tr_{U_j}\rho_j([v])$ of the pairwise non-isomorphic simple finite-dimensional $A(V)$-modules
$U_j=(M_j)_{h_0}$. By the Jacobson density theorem these characters are linearly independent (there is $x\in A(V)$ acting as
an idempotent of trace $\ne0$ on $U_{j_0}$ and as $0$ on the others), so some $c_j=0$ with $h_j=h_0$, a contradiction.
\end{proof}

\begin{physics}
At low temperature ($q\to0$) each character is dominated by its ground states, and the zero-mode algebra distinguishes
non-isomorphic ground-state multiplets by their traces; the $T$-twist separates sectors of different spin.
\end{physics}

\begin{status}
\lean{VOA.fibreIndependent_}\allowbreak\lean{of_bottomSeparation}: \texttt{propext, Classical.choice, Quot.sound}; conditional on the
displayed \lean{hcw}, \lean{hgr}, \lean{hsimp}, \lean{hsep}. No hypothesis on $V$ itself.
\end{status}

\section{Reading off the representation}\label{l15:sec:rho}

\subsection{Modular data and the analytic modularity law}

\begin{definition}[Modular data; Zhu's analytic law]\label{l15:def:modulardata}
For a finite index set $\iota$, \emph{modular data} for $V$ consist of a family $(M_i)_{i\in\iota}$ of irreducible modules, numbers
$h_i$ which are their conformal weights, a monoid homomorphism $\rho:\lfifteenSL\to\mathrm{Mat}_\iota(\C)$, and a matrix $S$
with $\rho\bigl(\begin{smallmatrix}0&-1\\1&0\end{smallmatrix}\bigr)=S$ and $\rho\bigl(\begin{smallmatrix}1&1\\0&1\end{smallmatrix}\bigr)=\mathrm{diag}(e^{2\pi i(h_i-c/24)})$.
They satisfy \emph{Zhu's analytic law} if for all $\gamma\in\lfifteenSL$, $n\in\N$, $v\in V_{[n]}$, $i\in\iota$, $\tau\in\lfifteenH$:
$S_i(v,\gamma\tau)=(c\tau+d)^n\sum_j\rho(\gamma)_{ij}S_j(v,\tau)$.
\end{definition}
\noindent Files \leanpath{VOA/\allowbreak Modularity/\allowbreak ZhuTheorem.lean},
\leanpath{VOA/\allowbreak Modularity/\allowbreak BlockSpanning.lean}.
\begin{leancode}
structure ModularData (W : VertexOperatorAlgebra 𝒱) (ι : Type*) [Fintype ι] [DecidableEq ι] where
  Mfam : ι → IrrModule W
  h : ι → ℂ
  hasWeight : ∀ i, HasConformalWeight (Mfam i).isMod (h i)
  rep : Matrix.SpecialLinearGroup (Fin 2) ℤ →* Matrix ι ι ℂ
  repS : ι → ι → ℂ
  rep_S : rep ModularGroup.S = Matrix.of repS
  rep_T : rep ModularGroup.T = Matrix.of (modularRepT W h)

noncomputable def autFactor (A : Matrix.SpecialLinearGroup (Fin 2) ℤ) (τ : UpperHalfPlane) : ℂ :=
  ((A : Matrix (Fin 2) (Fin 2) ℤ) 1 0 : ℂ) * (τ : ℂ)
    + ((A : Matrix (Fin 2) (Fin 2) ℤ) 1 1 : ℂ)

def ZhuModularityAnalytic (D : ModularData W ι) : Prop :=
  ∀ (A : SL(2, ℤ)) (n : ℕ) (v : 𝒱), v ∈ sqWeightSpace W (n : ℂ) → ∀ (i : ι) (τ : UpperHalfPlane),
    traceFun D.Mfam D.h i v (A • τ)
      = autFactor A τ ^ n * ∑ j, D.rep A i j * traceFun D.Mfam D.h j v τ
\end{leancode}

\begin{dictionary}
\textbf{Decoding.} \lean{[Fintype ι]}: $\iota$ is finite. \lean{Matrix.SpecialLinearGroup (Fin 2) ℤ} $=$ \lean{SL(2, ℤ)}
$=\lfifteenSL$. \lean{→*} is a monoid homomorphism: $\rho(1)=1$ and $\rho(AB)=\rho(A)\rho(B)$; as $\lfifteenSL$ is a group,
$\rho(A)$ is automatically invertible with inverse $\rho(A^{-1})$, so $\rho$ is a representation $\lfifteenSL\to GL_\iota(\C)$.
\lean{Matrix.of f} is the matrix with entries $f\,i\,j$. \lean{ModularGroup.S}, \lean{ModularGroup.T} are $S$ and $T$.
\lean{A • τ} is Mathlib's Möbius action $\gamma\tau$. \lean{autFactor A τ} is $c\tau+d$. \lean{traceFun D.Mfam D.h i v} is the
trace function $S_i(v,\cdot)$ of Definition~\ref{l15:def:traceFun}, built from the family and weights \emph{stored in $D$}.

\textbf{What is and is not constrained.} The field \lean{rep_T} pins $\rho(T)$ to the diagonal matrix proved in
Theorem~\ref{l15:thm:T}; \lean{repS} is free data. Because $S$ and $T$ generate $\lfifteenSL$, $\rho$ is determined by
\lean{repS} and $h$ (\lean{ModularData.rep_}\allowbreak\lean{eq_of_repS_eq}, via \lean{monoidHom_ext_SL2Z}). Nothing in \lean{ModularData}
requires $\rho(S)$ to be unitary or symmetric, or $\rho(S)^2$ to be a permutation matrix.

\textbf{\lean{ZhuModularityAnalytic} vs.\ Theorem~\ref{l15:thm:lit}(2).} Same quantifiers, same automorphy factor
$(c\tau+d)^n$ with $n$ the $L[0]$-weight, same index placement $\rho(\gamma)_{ij}$. Two points to keep in mind:
(a) the predicate reads its whole content off \lean{D.Mfam} and \lean{D.h}; a theorem concluding only
$\exists D,\ \lean{ZhuModularityAnalytic}\ D$ says nothing about any given family (Section~\ref{l15:sec:headline});
(b) \textbf{the predicate does not assert convergence and is not self-contained.} \lean{traceFun} uses \lean{tsum}, which is
$0$ on a divergent series; read alone, \lean{ZhuModularityAnalytic D} would be satisfied by junk-zero trace functions for any
\lean{rep} with the right $\rho(T)$. Convergence is a separate theorem (Theorem~\ref{l15:thm:growth},
\lean{onePointSubexp_of_c2}), and a separate Lean check against the library, \lean{headline_traces_are_genuine} (not
part of this release), proves, under the
headline's own hypotheses $C_2$ and \lean{hcw}, that every series in the conclusion converges, every trace function is
holomorphic on $\lfifteenH$, and none is identically zero. A referee must read the headline together with that check.
\end{dictionary}

\begin{physics}
The modular data $(h,S)$ of a rational CFT: the conformal weights (through $T$) and the modular $S$-matrix. The
transformation law says that the vector $(\langle v\rangle_{M_i,\tau})_i$ of one-point functions is a vector-valued modular
form of weight $n$ for the representation $\rho$. For $v=\lfifteenvac$ ($n=0$) it is the statement that the characters form a
vector-valued modular function, which makes $\sum_{ij}N_{ij}\chi_i\bar\chi_j$ modular invariant whenever the matrix $N$
commutes with $\rho$.
\end{physics}

\begin{status}
Definitions; no axioms beyond Mathlib's.
\end{status}

\subsection{From spanning and independence to modular data}

\begin{theorem}[$\rho$ from spanning and independence]\label{l15:thm:spanning}
Suppose (spanning) for every $\gamma$ there is a matrix $\rho_\gamma$ with $S_i(v,\gamma\tau)=(c\tau+d)^n\sum_j(\rho_\gamma)_{ij}S_j(v,\tau)$
for all homogeneous $v$ of weight $n$, and (independence) the $S_j$ are linearly independent. Then $\rho_\gamma$ is unique,
$\gamma\mapsto\rho_\gamma$ is a homomorphism, $\rho_T=\mathrm{diag}(e^{2\pi i(h_i-c/24)})$, and the resulting modular data, with
the given family and weights, satisfy Zhu's analytic law. Moreover spanning at the single matrix $S$ implies spanning
at every $\gamma$.
\end{theorem}
\noindent Files \leanpath{VOA/\allowbreak Modularity/\allowbreak BlockSpanning.lean},
\leanpath{VOA/\allowbreak Modularity/\allowbreak TraceSpanningS.lean}.
\begin{leancode}
def TraceSpanning (Mfam : ι → IrrModule W) (h : ι → ℂ) : Prop :=
  ∀ A : SL(2, ℤ), ∃ ρA : ι → ι → ℂ, ∀ (n : ℕ) (v : 𝒱), v ∈ sqWeightSpace W (n : ℂ) →
    ∀ (i : ι) (τ : UpperHalfPlane),
      traceFun Mfam h i v (A • τ) = autFactor A τ ^ n * ∑ j, ρA i j * traceFun Mfam h j v τ

noncomputable def modularDataOf (hcw : ∀ i, HasConformalWeight (Mfam i).isMod (h i)) :
    ModularData W ι where
  Mfam := Mfam
  h := h
  hasWeight := hcw
  rep := modularRepOf Mfam h hsp hind
  repS := rhoOf Mfam h hsp ModularGroup.S
  rep_S := rfl

theorem zhu_modularity_of_spanning (hcw : ∀ i, HasConformalWeight (Mfam i).isMod (h i)) :
    ZhuModularityAnalytic (modularDataOf Mfam h hsp hind hcw) := by

theorem traceSpanning_of_S (hcw : ∀ i, HasConformalWeight (Mfam i).isMod (h i))
    (hS : ∃ ρS : ι → ι → ℂ, RealisesLaw Mfam h ModularGroup.S ρS) :
    TraceSpanning Mfam h := by
\end{leancode}

\begin{dictionary}
\lean{hsp : TraceSpanning Mfam h} and \lean{hind : TraceIndependent Mfam h} are section variables of
\lean{modularDataOf} and \lean{zhu_modularity_}\allowbreak\lean{of_spanning}. \lean{RealisesLaw Mfam h A ρ} is the law at the single matrix
$A$ with matrix $\rho$. \lean{modularDataOf} visibly stores the \emph{given} \lean{Mfam} and \lean{h}; its \lean{rep} is
$\gamma\mapsto\rho_\gamma$ (\lean{rhoOf} chooses $\rho_\gamma$ from \lean{hsp} with \lean{Classical.choose}, and independence
makes the choice irrelevant). Same as the textbook argument.
\end{dictionary}

\begin{proof}[Proof on paper]
Uniqueness: if $\rho,\rho'$ both realise the law at $\gamma$, then $\sum_j(\rho_{ij}-\rho'_{ij})S_j=0$ (divide by $(c\tau+d)^n\ne0$),
so $\rho=\rho'$ by independence. Homomorphism: $j(\gamma\delta,\tau)=j(\gamma,\delta\tau)j(\delta,\tau)$ for
$j(\gamma,\tau)=c\tau+d$ (the cocycle \lean{autFactor_mul}), so $\rho_\gamma\rho_\delta$ realises the law at $\gamma\delta$, and
by uniqueness equals $\rho_{\gamma\delta}$; likewise $\rho_1=1$. $T$: by Theorem~\ref{l15:thm:T} the diagonal matrix realises
the law at $T$. From $S$ to all $\gamma$: the set of $\gamma$ at which the law is realisable is closed under products and contains
$S$ and $T$; since $S^{-1}=S^3$ and $T^{-1}=(ST)^5S$, the \emph{monoid} generated by $S,T$ is all of $\lfifteenSL$.
\end{proof}

\begin{physics}
Modular invariance need only be checked on the two generators: $T$ (the Dehn twist, which is automatic) and $S$ (the
exchange of space and time, which is the real content). This is how physicists verify modular invariance of a proposed
torus partition function.
\end{physics}

\begin{status}
\lean{VOA.zhu_}\allowbreak\lean{modularity_}\allowbreak\lean{of_spanning}: \texttt{propext, Classical.choice, Quot.sound}; conditional on the two named data.
\lean{VOA.traceSpanning_of_S}: same axioms; conditional on the $S$-law.
\end{status}

\subsection{Spanning at \texorpdfstring{$S$}{S} from the dimension count}

\begin{theorem}[$S$-spanning from a rank bound]\label{l15:thm:rank}
Let $\mathfrak A$ be any space of functions on homogeneous points containing $\mathrm{res}(S_i)$ and $\mathrm{res}(S_i|S)$ for all $i$,
of dimension $\le\#\iota$, and suppose the $S_i$ are independent. Then each $\mathrm{res}(S_i|S)$ is a linear combination of the
$\mathrm{res}(S_j)$. Under the hypotheses of Theorem~\ref{l15:thm:533} the space $\mathfrak A=\mathrm{res}(\lfifteenB)$ has dimension
$\le\#\Phi(V)$.
\end{theorem}
\noindent Files \leanpath{VOA/\allowbreak Modularity/\allowbreak RhoS.lean} (Chapter~\ref{ch:L14}),
\leanpath{VOA/\allowbreak Modularity/\allowbreak SlashAmbientObstruction.lean}, \leanpath{VOA/\allowbreak Modularity/\allowbreak SqAmbientRank.lean}.
\begin{leancode}
def SlashSpannedAtS (Mfam : ι → IrrModule W) (h : ι → ℂ) : Prop :=
  ∀ i : ι, ∃ c : ι → ℂ, ∀ (n : ℤ) (v : 𝒱), v ∈ sqWeightSpace W (n : ℂ) → ∀ z ∈ imGT 0,
    blockSlash W ModularGroup.S (traceBlock Mfam h i) v z
      = ∑ j, c j * traceBlock Mfam h j v z

theorem slashSpannedAtS_of_sqRank_le (𝔄 : Submodule ℂ (SqPt W → ℂ → ℂ))
    (hind : TraceIndependent Mfam h)
    (htr : ∀ i : ι, sqRestrict W (traceBlock Mfam h i) ∈ 𝔄)
    (hsl : ∀ i : ι, sqRestrict W (blockSlash W ModularGroup.S (traceBlock Mfam h i)) ∈ 𝔄)
    (hrank : Module.rank ℂ 𝔄 ≤ (Fintype.card ι : Cardinal)) :
    SlashSpannedAtS Mfam h := by

theorem traceSpanning_of_slashSpannedAtS (hcw : ∀ i, HasConformalWeight (Mfam i).isMod (h i))
    (hsub : ∀ i, OnePointSubexp (Mfam i).isMod (h i - W.centralCharge / 24))
    (hS : SlashSpannedAtS Mfam h) : TraceSpanning Mfam h := by

theorem rank_map_sqRestrict_analyticBlocks_le [Finite (Φ W)]
    (hzero : VacuumGradeZero W) (hC2 : C2Cofinite W)
    (hss : letI := instRingA W; IsSemisimpleRing (A W))
    (hreal : SimpleAModsRealized W)
    (Pr : BottomProjection (parabolicRepOfAMod W (regularAMod W)))
    (hcard : Fintype.card ι = Nat.card (Φ W)) :
    Module.rank ℂ (Submodule.map (sqRestrict W) (analyticBlocks W))
      ≤ (Fintype.card ι : Cardinal) := by
\end{leancode}

\begin{dictionary}
\lean{Module.rank} is the dimension as a cardinal (so no \lean{finrank}-is-$0$ trap). \lean{Fintype.card ι} is $\#\iota$;
\lean{hcard} says the family has exactly $\#\Phi(V)$ members (in the headline it forces one representative per
isomorphism class; Section~\ref{l15:sec:headline}). \lean{SlashSpannedAtS} uses the pointwise slash, but only on homogeneous $v$,
where it is Zhu's action. \lean{traceSpanning_}\allowbreak\lean{of_slashSpannedAtS} converts the $S$-statement about blocks into the
$S$-law for \lean{traceFun} (using \lean{traceBlock_coe} and that $\mathrm{denom}_S(\tau)^{-n}$ is $(c\tau+d)^{-n}$) and then applies
\lean{traceSpanning_of_S}. \lean{rank_map_}\allowbreak\lean{sqRestrict_}\allowbreak\lean{analyticBlocks_le} is Theorem~\ref{l15:thm:533} pushed through the linear
map $\mathrm{res}$ (dimension cannot increase).
\end{dictionary}

\begin{proof}[Proof on paper]
The $\#\iota$ functions $\mathrm{res}(S_i)$ are independent (independence of the $S_i$ is already a statement on homogeneous
vectors) and lie in $\mathfrak A$ of dimension $\le\#\iota$; so they are a basis of $\mathfrak A$, and
$\mathrm{res}(S_i|S)\in\mathfrak A$ is a combination of them.
\end{proof}

\begin{physics}
Since there are as many characters as independent torus blocks, the characters \emph{are} the blocks; the $S$-transform of a
character is a block, hence a combination of characters: the coefficients are the $S$-matrix.
\end{physics}

\begin{status}
All three theorems: \texttt{propext, Classical.choice, Quot.sound}. \lean{rank_map_}\allowbreak\lean{sqRestrict_}\allowbreak\lean{analyticBlocks_le} is
conditional on the hypotheses of Theorem~\ref{l15:thm:533} and \lean{hcard}.
\end{status}

\section{The modularity theorem}\label{l15:sec:headline}

The section has four parts. First the general \emph{pinned} form, for an arbitrary complete list of irreducible modules
(Theorem~\ref{l15:thm:headline}); it is the theorem that does the work. Then its specialisation to the canonical list of
representatives, which has only the two hypotheses $C_2$ and ``$A(V)$ semisimple'' (the engine,
Section~\ref{l15:sec:canonical}). Then the chapter's headline, in Zhu's own form ``$C_2$-cofinite and rational'', obtained
from the engine through the Lean theorem ``rational $\Rightarrow A(V)$ semisimple'' (Section~\ref{l15:sec:rational}).
Finally the older strongly rational packaging, now superseded.

\subsection{The statement for a given family}

\begin{theorem}[Zhu's modularity theorem, pinned to a given family]\label{l15:thm:headline}
Let $V$ be a vertex operator algebra of the library such that
\begin{enumerate}
\item[(a)] $V$ is $C_2$-cofinite;
\item[(b)] Zhu's algebra $A(V)$ is semisimple;
\item[(c)] $\Phi(V)$, the set of isomorphism classes of irreducible ordinary modules, is finite;
\item[(d)] the datum \lean{Pr : BottomProjection (parabolicRepOfAMod W (regularAMod W))} is given (Chapter~\ref{ch:L13}).
\end{enumerate}
Let $\iota$ be a finite set with $\#\iota=\#\Phi(V)$ (hypothesis \lean{hcard}: one index per isomorphism class), and $(M_i,h_i)_{i\in\iota}$ irreducible ordinary modules with conformal
weights $h_i$, whose bottoms $(M_i)_{h_i}$ are simple $A(V)$-modules, pairwise non-isomorphic whenever $h_j=h_k$, $j\neq k$.
Then there are modular data $D$ with $D.\mathrm{Mfam}=(M_i)$ and $D.h=(h_i)$ satisfying Zhu's analytic law: there is a
homomorphism $\rho:\lfifteenSL\to\mathrm{Mat}_\iota(\C)$ with $\rho(T)=\mathrm{diag}(e^{2\pi i(h_i-c/24)})$ such that for all
$\gamma\in\lfifteenSL$, $n\in\N$, $v\in V_{[n]}$, $i\in\iota$, $\tau\in\lfifteenH$,
\begin{equation*}
  \tr_{M_i}o(v)q^{L_0-c/24}\big|_{\gamma\tau}=(c\tau+d)^n\sum_{j\in\iota}\rho(\gamma)_{ij}\,\tr_{M_j}o(v)q^{L_0-c/24}\big|_{\tau}.
\end{equation*}
Since (d) is supplied for every $V$ by the theorem \lean{VOA.pbw} (through \lean{voaPBWBottomProjection}), the
conclusion holds under (a)--(c) (this is \lean{zhuModularity_pinned_of_C2Cofinite_voaPBW}).
\end{theorem}
\noindent Files \leanpath{VOA/\allowbreak Modularity/\allowbreak GNProp8Closure.lean} and
\leanpath{VOA/\allowbreak Modularity/\allowbreak GNProp8PBW.lean}.
\begin{leancode}
theorem zhuModularity_pinned_of_C2Cofinite [Finite (Φ W)]
    (Mfam : ι → IrrModule W) (h : ι → ℂ)
    (hC2 : C2Cofinite W)
    (hss : letI := instRingA W; IsSemisimpleRing (A W))
    (Pr : BottomProjection (parabolicRepOfAMod W (regularAMod W)))
    (hcard : Fintype.card ι = Nat.card (Φ W))
    (hcw : ∀ i, HasConformalWeight (Mfam i).isMod (h i))
    (hsep : ∀ j k : ι, j ≠ k → h j = h k →
      IsEmpty (AModEquiv W (lowestAMod (hcw j).2) (lowestAMod (hcw k).2))) :
    ∃ D : ModularData W ι, D.Mfam = Mfam ∧ D.h = h ∧ ZhuModularityAnalytic D :=
  zhuModularity_pinned_of_strictSpanning Mfam h hC2 hss Pr
    (strictSpanning_of_C2Cofinite W hC2) hcard hcw hsep

theorem zhuModularity_pinned_of_C2Cofinite_voaPBW [Finite (Φ W)]
    (Mfam : ι → IrrModule W) (h : ι → ℂ)
    (hC2 : C2Cofinite W)
    (hss : letI := instRingA W; IsSemisimpleRing (A W))
    (hcard : Fintype.card ι = Nat.card (Φ W))
    (hcw : ∀ i, HasConformalWeight (Mfam i).isMod (h i))
    (hsep : ∀ j k : ι, j ≠ k → h j = h k →
      IsEmpty (AModEquiv W (lowestAMod (hcw j).2) (lowestAMod (hcw k).2))) :
    ∃ D : ModularData W ι, D.Mfam = Mfam ∧ D.h = h ∧ ZhuModularityAnalytic D :=
  zhuModularity_pinned_of_C2Cofinite Mfam h hC2 hss
    (voaPBWBottomProjection (parabolicRepOfAMod W (regularAMod W)))
\end{leancode}

\begin{dictionary}
We go through the statement symbol by symbol; the section variables are
\lean{{𝒱 : Type u} [AddCommGroup 𝒱] [Module ℂ 𝒱] {W : VertexOperatorAlgebra 𝒱}} and
\lean{{ι : Type*} [Fintype ι] [DecidableEq ι]}.

\textbf{Hypotheses, one by one.}
\begin{itemize}
\item \lean{[Finite (Φ W)]}: (c). An instance argument, i.e.\ a hypothesis found automatically when available.
\item \lean{Mfam}, \lean{h}: the family and its weights; arbitrary data constrained only by the hypotheses below.
\item \lean{hzero : VacuumGradeZero W}: $V_0\subseteq\C\lfifteenvac$. \emph{Redundant}: true for every VOA of the library
  (Section~\ref{l15:sec:convergence}).
\item \lean{hC2 : C2Cofinite W}: (a), Definition~\ref{l12:def:c2}.
\item \lean{hss}: (b). \lean{IsSemisimpleRing} is Mathlib's notion for the ring \lean{A W} with Zhu's product.
\item \lean{Pr}: (d). In the \lean{_voaPBW} form it is supplied by \lean{voaPBWBottomProjection}, which uses the theorem \lean{VOA.pbw}.
\item \lean{hcard : Fintype.card ι = Nat.card (Φ W)}: the index set has \emph{exactly} $\#\Phi(V)$ elements. With \lean{hcw}
  and \lean{hsep} this forces the family to be one representative of each isomorphism class of irreducible modules, i.e.\ a
  complete irredundant list (argued in ``The family is a complete list'' below). It is the hypothesis that makes
  ``the traces span the blocks'' possible: $\#\iota$ independent traces in a space of dimension $\le\#\Phi(V)=\#\iota$.
\item \lean{hcw}: each $M_i$ has conformal weight $h_i$ (nonzero bottom $(M_i)_{h_i}$, all weights in $h_i+\N$).
  It ties the free data \lean{h} to the modules. For irreducible modules such an $h_i$ always exists and is unique
  (\lean{exists_hasConformalWeight_}\allowbreak\lean{of_irreducible}, \leanpath{LowestWeightExistence.lean}).
\item (removed) \lean{hsimp}: each bottom is a simple $A(V)$-module. It was \emph{redundant given \lean{hss} and \lean{Pr}} and the
  current headline no longer carries it: it follows from
  Chapter~\ref{ch:L12}'s \lean{isSimpleAMod_}\allowbreak\lean{of_irreducible_}\allowbreak\lean{of_semisimple} (Theorem~\ref{l12:thm:simplebottom}).
\item \lean{hsep}: bottoms of distinct members with equal weight are non-isomorphic.
\end{itemize}

\textbf{The family is a complete list.} The hypotheses force $(M_i)$ to be a system of representatives of $\Phi(V)$: if
$M_j\cong M_k$ with $j\neq k$, then $h_j=h_k$ (the conformal weight of an irreducible module is an isomorphism invariant) and
the bottoms are isomorphic, contradicting \lean{hsep}; so $i\mapsto[M_i]$ is injective, and as $\#\iota=\#\Phi(V)$ is finite,
bijective. (This is an argument on paper; the library does not state it as a lemma, and does not need to.) Thus the
theorem is exactly about \emph{all} irreducible modules, as Zhu's is.

\textbf{The conclusion is pinned.} \lean{∃ D : ModularData W ι, D.Mfam = Mfam ∧ D.h = h ∧ ZhuModularityAnalytic D}: the
modular data are those \emph{of the given family}. Without the two equations the statement would be much weaker (see the
warning below).

\textbf{Same as, stronger or weaker than Zhu?} Compared with Theorem~\ref{l15:thm:lit}:
\begin{itemize}
\item \emph{Hypotheses on $V$}: ``rational'' is replaced by ``$A(V)$ semisimple and $\Phi(V)$ finite''; and finiteness of
  $\Phi(V)$ is itself a theorem under $C_2$ (\lean{finite_Φ_of_C2Cofinite_pbw}, Section~\ref{l15:sec:canonical}). For
  rational $V$ semisimplicity follows (Zhu; Dong--Li--Mason), and the library proves that implication for its predicate
  \lean{AdmissiblyRational} (\lean{isSemisimpleRing_}\allowbreak\lean{A_of_}\allowbreak\lean{admissiblyRational}, Theorem~\ref{l15:thm:zhuss}). Under $C_2$,
  Dong--Li--Mason rationality implies \lean{AdmissiblyRational} in Lean (\lean{admissiblyRational_of_dlmRational}); only the
  converse is open (Remark~\ref{l15:rem:hyp}). In the
  other direction we make no claim that the library's hypotheses imply rationality (Dong and Jiang characterise
  rationality by semisimplicity of \emph{all} the higher Zhu algebras $A_n(V)$, not of $A(V)$ alone), so as a theorem
  about $V$ the library's form is at least as general as Zhu's.
\item \emph{Conclusion}: (2) and (3) of Theorem~\ref{l15:thm:lit} for the one-point functions, with $\rho$ a representation.
  Convergence (1) is proved by \lean{onePointSubexp_of_c2} under the same hypotheses but is not part of the conclusion
  (next item). Zhu's theorem also covers $n$-point functions on the torus; the library does not.
\item \emph{Convergence}: \lean{ZhuModularityAnalytic} uses \lean{traceFun}, whose \lean{tsum} would be $0$ on a divergent
  series. Read alone, the conclusion would be compatible with all trace functions being junk zeros. It is not: under
  \lean{hC2} and \lean{hcw}, a separate Lean check, \lean{headline_}\allowbreak\lean{traces_are_genuine} (below), proves that every series in the
  conclusion converges, that every trace function is holomorphic on $\lfifteenH$, and that each is not identically zero.
\end{itemize}

\textbf{How it could have been vacuous, and why it is not.} (1) If the hypotheses were jointly unsatisfiable: they are
satisfied by $V=\C\lfifteenvac$ (Section~\ref{l15:sec:nonvac}). (2) If the conclusion could be met by an unrelated $D$: it is
pinned. (3) If the trace functions were junk: they are not (checked in Lean; the check is not part of this release). (4) If $\rho$ were unconstrained: by uniqueness
(independence of the traces) $\rho$ is determined by the family, and $\rho(T)$ is the correct diagonal matrix. What remains
\emph{weak} is the hypothesis bundle as displayed: \lean{[Finite (Φ W)]} and the family data are dischargeable but are
displayed; the canonical form of Section~\ref{l15:sec:canonical} discharges them. The hypotheses are now checked at
nontrivial $V$: the lattice VOAs $V_{\sqrt{2m}\Z}$ and the $c=1$ orbifolds of Section~\ref{l15:sec:instances} (before this
release, only at the trivial VOA);
and \lean{Pr}, which had no axiom-free witness while \lean{VOA.pbw} was an axiom, is supplied for every $V$ by
\lean{voaPBWBottomProjection}.
\end{dictionary}

\begin{warning}
\textbf{Two redundant hypotheses (moderate).} A separate Lean check against the library (not part of this release)
proves, with the standard three axioms only, that \lean{VacuumGradeZero W} holds for every $W$
(\lean{vacuumGradeZero_always}, by \lean{le_of_eq W.vacuum_grade}), that \lean{hsimp} follows from
\lean{hss} and \lean{Pr} (\lean{hsimp_of_semisimple}), and that the \lean{Pr}-form theorem holds without these two
hypotheses (\lean{zhuModularity_pinned_lean}).
So the honest hypothesis list of the \lean{Pr}-form theorem is: $C_2$, $A(V)$ semisimple, $\Phi(V)$ finite, \lean{Pr}, and the
family data \lean{hcard}, \lean{hcw}, \lean{hsep}. \emph{Correction (release audit).} An earlier version of this box said that
\lean{[Finite (Φ W)]} was ``plausibly derivable \ldots{} not certified''. That was stale: it is the library theorem
\lean{VOA.Huang.ZhuFiniteIrr.finite_Φ_of_C2Cofinite_pbw} (Theorem~\ref{l15:thm:finitephi}), which derives it from $C_2$
alone, without \lean{hss}. The family data are discharged by the canonical family. Both discharges are carried out in
the library's headline, Theorem~\ref{l15:thm:canonical}, whose only hypotheses are $C_2$ and \lean{hss}.
\end{warning}

\begin{warning}
\textbf{Unpinned variants: do not cite them as Zhu's theorem.} Several theorems on the way conclude only
\lean{∃ D : ModularData W ι, ZhuModularityAnalytic D}: \lean{exists_zhuModularity_of_c2}, \lean{exists_zhuModularity_of_ode},
\lean{exists_zhuModularity_}\allowbreak\lean{of_strictSpanning}, \lean{exists_zhuModularity_}\allowbreak\lean{of_cofinite},
\lean{exists_zhuModularity_}\allowbreak\lean{of_c2_sqSlash}. Such a statement is \emph{not} about the given family: \lean{ZhuModularityAnalytic}
reads everything off \lean{D.Mfam} and \lean{D.h}, so any single $D$ whatsoever (e.g.\ for the trivial VOA) witnesses it.
The library's own docstrings say so (``Conclusion-strength warning'') and the README forbids advertising them as the result.
The pinned forms are \lean{zhuModularity_}\allowbreak\lean{canonical_of_}\allowbreak\lean{C2Cofinite_semisimple} (the headline),
\lean{zhuModularity_}\allowbreak\lean{pinned_of_}\allowbreak\lean{C2Cofinite}, \lean{zhuModularity_}\allowbreak\lean{pinned_of_}\allowbreak\lean{C2Cofinite_voaPBW},
\lean{zhuModularity_}\allowbreak\lean{pinned_of_}\allowbreak\lean{strictSpanning}, \lean{zhuModularity_}\allowbreak\lean{pinned_of_real}, \lean{zhu_modularity_}\allowbreak\lean{of_bottomSeparation}
and \lean{zhuModularity_}\allowbreak\lean{stronglyRational_voaPBW}.
\end{warning}

\begin{proof}[Proof on paper (the assembly)]
Fix the family. (1) Growth: $\lfifteenZ_{M_i}$ has subexponential coefficients (Theorem~\ref{l15:thm:growth}), so the trace
functions $S_i$ are analytic blocks (Theorem~\ref{l15:thm:traceblock}). (2) Existence of modules: $V$ has strict spanning
(Theorem~\ref{l15:thm:gn}), hence every simple $A(V)$-module is realised (Theorem~\ref{l15:thm:realized}). (3) Count:
$\dim\mathrm{res}(\lfifteenB)\le\#\Phi(V)=\#\iota$ (Theorems~\ref{l15:thm:533} and~\ref{l15:thm:rank}). (4) Independence: the
$S_i$ are independent (Theorem~\ref{l15:thm:indep}, using \lean{hsimp}, \lean{hsep} and the growth). (5) Slash: each
$\mathrm{res}(S_i|S)$ lies in $\mathrm{res}(\lfifteenB)$ (Theorem~\ref{l15:thm:slashblock}). (6) By (3)--(5) and
Theorem~\ref{l15:thm:rank}, each $S_i|S$ is a combination of the $S_j$ on homogeneous vectors; by Theorem~\ref{l15:thm:spanning}
this gives the law at every $\gamma$, a unique representation $\rho$ with the correct $\rho(T)$, and modular data whose
family and weights are the given ones.
\end{proof}

\begin{physics}
This is the theorem behind every modular-invariant partition function of a rational CFT: the characters of the finitely
many sectors span a representation of the modular group of the torus, and so does every family of torus one-point
functions of a field of definite conformal weight. The matrix $\rho(S)$ is the modular $S$-matrix of the theory (the
datum consumed by the Verlinde formula $N_{ij}^{k}=\sum_m S_{im}S_{jm}\overline{S_{km}}/S_{0m}$, proved by Huang for
rational $C_2$-cofinite VOAs; not in this chapter, see Chapter~\ref{ch:L21}); $\rho(T)$ encodes the spins of the primaries. The characters form a vector-valued
modular form, which is the starting point of the modular bootstrap and of the classification of rational CFTs by the
Mathur--Mukhi--Sen differential equation.
\end{physics}

\begin{status}
Independent run on a scratch file, \texttt{lake env lean -DmaxSynthPendingDepth=3} (never a bare \texttt{lake env lean}):
\begin{plaincode}
'VOA.zhuModularity_pinned_of_C2Cofinite' depends on axioms: [propext, Classical.choice, Quot.sound]
'VOA.zhuModularity_pinned_of_C2Cofinite_voaPBW' depends on axioms: [propext, Classical.choice, Quot.sound]
\end{plaincode}
So: \lean{zhuModularity_}\allowbreak\lean{pinned_of_}\allowbreak\lean{C2Cofinite} is \texttt{\#print axioms}-clean, \emph{conditional on} \lean{C2Cofinite},
\lean{IsSemisimpleRing (A W)}, \lean{Finite (Φ W)}, \lean{Pr}, and the family data \lean{hcard}, \lean{hcw}, \lean{hsep}
(the formerly redundant \lean{hzero}, \lean{hsimp} have been removed). \lean{zhuModularity_}\allowbreak\lean{pinned_of_}\allowbreak\lean{C2Cofinite_voaPBW} uses the theorem
\lean{VOA.pbw}; it is conditional on $C_2$, semisimplicity, finiteness of $\Phi(V)$ and the family data (no longer on
\lean{Pr}). Neither is unconditional. No full \texttt{lake build} was run for this chapter; the oleans of
the snapshot were used (not re-run for this edition), and the \lean{_voaPBW} line was re-probed in a later run, after
\lean{VOA.pbw} became a theorem (in the earlier run it also listed \lean{VOA.pbw}, then an axiom).
\end{status}

\subsection{The engine: the canonical family, from \texorpdfstring{$C_2$}{C2} and semisimplicity alone}\label{l15:sec:canonical}

Theorem~\ref{l15:thm:headline} asks the user to supply a finiteness instance and a family with three properties. All of
them can be discharged once and for all. Finiteness of $\Phi(V)$ is a consequence of $C_2$-cofiniteness; and the family can
be taken to be ``one chosen representative of each isomorphism class, with its conformal weight''. Doing both gives the
form of Zhu's theorem from which the headline (Section~\ref{l15:sec:rational}) follows.

\begin{theorem}[Finitely many irreducible modules; Dong--Li--Mason, Abe--Buhl--Dong]\label{l15:thm:finitephi}
If $V$ is $C_2$-cofinite, then $\Phi(V)$ is finite.
\end{theorem}
\noindent File \leanpath{VOA/\allowbreak Classification/\allowbreak COne/\allowbreak McRaeO1FinitePhi.lean}
(Chapter~\ref{ch:L22}), in the namespace \lean{VOA.Huang.ZhuFiniteIrr}, with section variable
\lean{{W : VertexOperatorAlgebra 𝒱}}.
\begin{leancode}
theorem finite_Φ_of_C2Cofinite_pbw (hC2 : C2Cofinite W) : Finite (Φ W) :=
  finite_Φ_of_C2Cofinite' (voaPBWBottomProjection (parabolicRepOfAMod W (regularAMod W))) hC2
\end{leancode}

\begin{dictionary}
\lean{Finite (Φ W)} is Mathlib's ``the type $\Phi(V)$ has finitely many elements'' (Chapter~\ref{ch:L12} for $\Phi$). The
one-line proof term is informative: the theorem is the general \lean{finite_Φ_of_C2Cofinite'}, which takes the datum
\lean{Pr}, applied to \lean{voaPBWBottomProjection}, i.e.\ to the theorem \lean{VOA.pbw}. The only hypothesis is $C_2$; in
particular no semisimplicity of $A(V)$ and no rationality is assumed. Same as the literature (for $C_2$-cofinite $V$,
$A(V)$ is finite-dimensional, hence has finitely many simple modules, and each irreducible module is determined by its
bottom). The docstring still calls \lean{VOA.pbw} ``the sanctioned axiom''; it is a theorem (documentation only). The
analogue for \emph{admissible} modules, Dong--Li--Mason's ``finitely many irreducible admissible modules'' (the predicate
\lean{DLMFinitelyManyIrredAdmissible}), follows for $C_2$-cofinite admissibly rational $V$, because there every irreducible
admissible module is ordinary; it is the theorem \lean{dlmFinitelyManyIrredAdmissible_of_c2Cofinite} of
Chapter~\ref{ch:L12} (Theorem~\ref{l12:thm:litproved}).
\end{dictionary}

\begin{proof}[Proof on paper]
$A(V)$ is finite-dimensional (Theorem~\ref{l15:thm:Afin}). A finite-dimensional algebra has only finitely many simple
modules up to isomorphism (each is a quotient of $A(V)$ by a maximal left ideal, and they are the simple modules of the
semisimple quotient $A(V)/\mathrm{rad}A(V)$, of which there are finitely many by Artin--Wedderburn). The bottom of an
irreducible module is a simple $A(V)$-module (Chapter~\ref{ch:L13}, Theorem~\ref{l13:thm:simplebottom-all}; the
library's proof of this theorem instead goes through the datum \lean{Pr}), and by the uniqueness half of Zhu's correspondence
(Chapter~\ref{ch:L13}) non-isomorphic irreducible modules have non-isomorphic bottoms. So $[M]\mapsto[M_h]$ is an
injection of $\Phi(V)$ into a finite set.
\end{proof}

\begin{physics}
A theory whose chiral algebra has finitely many independent fields modulo derivatives and normal ordering has only finitely
many sectors. This is the algebraic shadow of ``finitely many primaries''.
\end{physics}

\begin{status}
\lean{VOA.Huang.ZhuFiniteIrr.finite_Φ_of_C2Cofinite_pbw}: \texttt{[propext, Classical.choice, Quot.sound]}, run at
an earlier commit (not re-run for this edition) (scratch file, \texttt{lake env lean -DmaxSynthPendingDepth=3}). Conditional only on \lean{C2Cofinite};
uses the theorem \lean{VOA.pbw}.
\end{status}

\begin{definition}[The canonical family]\label{l15:def:canonical}
For $c\in\Phi(V)$ let $M_c$ be a chosen representative of the class $c$, and $h_c$ its conformal weight (the unique $h$ with
$(M_c)_h\neq0$ and all weights in $h+\N$; it exists for every irreducible ordinary module).
\end{definition}
\noindent Files \leanpath{VOA/\allowbreak Modularity/\allowbreak StronglyRationalFamily.lean} and
\leanpath{VOA/\allowbreak Modularity/\allowbreak ConformalWeightChoice.lean}.
\begin{leancode}
noncomputable def canonicalIrrModule (c : Φ W) : IrrModule W := Quotient.out c

noncomputable def conformalWeightOf (V : IrrModule W) : ℂ :=
  (OrdVModule.exists_hasConformalWeight_of_irreducible V.isMod V.irreducible).choose
\end{leancode}

\begin{dictionary}
\lean{Quotient.out c} picks a representative of the class \lean{c} (by the axiom of choice). \lean{conformalWeightOf} picks the
$h$ whose existence is \lean{exists_hasConformalWeight_of_irreducible}; it is unique, so the choice is harmless. Which
representative is picked does not matter for the theorem below: the trace functions of isomorphic modules are equal.
\end{dictionary}

\begin{physics}
One primary field per sector, labelled by its conformal weight: the list of sectors of the theory.
\end{physics}

\begin{status}
Definitions.
\end{status}

\begin{theorem}[Zhu's modularity theorem for one-point functions, canonical family, semisimple form]\label{l15:thm:canonical}
Let $V$ be a vertex operator algebra of the library (so of CFT type) such that
\begin{enumerate}
\item[(a)] $V$ is $C_2$-cofinite, and
\item[(b)] Zhu's algebra $A(V)$ is a semisimple ring.
\end{enumerate}
Then $\Phi(V)$ is finite, and there are modular data $D$ indexed by $\Phi(V)$, whose family is the canonical family
$(M_c)_{c\in\Phi(V)}$ and whose weights are its conformal weights $(h_c)$, satisfying Zhu's analytic law: there is a group
homomorphism $\rho:\lfifteenSL\to\mathrm{Mat}_{\Phi(V)}(\C)$ with $\rho(T)=\mathrm{diag}(e^{2\pi i(h_c-c_V/24)})$ such that for
all $\gamma=\left(\begin{smallmatrix}a&b\\ c&d\end{smallmatrix}\right)\in\lfifteenSL$, $n\in\N$, $v\in V_{[n]}$,
$c'\in\Phi(V)$ and $\tau\in\lfifteenH$,
\begin{equation*}
  \tr_{M_{c'}}o(v)q^{L_0-c_V/24}\big|_{\gamma\tau}=(c\tau+d)^n\sum_{c''\in\Phi(V)}\rho(\gamma)_{c'c''}\,
  \tr_{M_{c''}}o(v)q^{L_0-c_V/24}\big|_{\tau},
\end{equation*}
where $c_V$ is the central charge.
\end{theorem}
\noindent File \leanpath{VOA/\allowbreak Modularity/\allowbreak ZhuModularityCanonical.lean}.
\begin{leancode}
theorem zhuModularity_canonical_of_C2Cofinite_semisimple (W : VertexOperatorAlgebra 𝒱)
    (hC2 : C2Cofinite W) (hss : letI := instRingA W; IsSemisimpleRing (A W)) :
    letI : Finite (Φ W) := VOA.Huang.ZhuFiniteIrr.finite_Φ_of_C2Cofinite_pbw hC2
    letI : Fintype (Φ W) := Fintype.ofFinite (Φ W)
    letI : DecidableEq (Φ W) := Classical.decEq (Φ W)
    ∃ D : ModularData W (Φ W), D.Mfam = canonicalIrrModule ∧
      D.h = (fun c => conformalWeightOf (canonicalIrrModule c)) ∧ ZhuModularityAnalytic D := by
\end{leancode}

\begin{dictionary}
\textbf{Hypotheses.} Exactly two, besides the ambient \lean{{𝒱 : Type u} [AddCommGroup 𝒱] [Module ℂ 𝒱]}.
\begin{itemize}
\item \lean{hC2 : C2Cofinite W}: $\dim V/C_2(V)<\infty$ (Definition~\ref{l12:def:c2}). Standard.
\item \lean{hss}: \lean{letI := instRingA W} installs Zhu's product on \lean{A W} (Chapter~\ref{ch:L12},
  Section~\ref{l12:sec:zhu}), and \lean{IsSemisimpleRing (A W)} is Mathlib's ``every module over this ring is semisimple''
  (equivalently, by Artin--Wedderburn, a finite product of matrix rings over division rings; over $\C$ and finite-dimensional,
  of matrix algebras $\mathrm{Mat}_{n_i}(\C)$). In this form it stands in for Zhu's ``rational''; the headline
  (Theorem~\ref{l15:thm:rational}) puts rationality back, and Remark~\ref{l15:rem:hyp} says exactly how the two compare.
\end{itemize}
CFT type is not a hypothesis because it is built into \lean{VertexOperatorAlgebra} (the axioms \lean{vacuum_grade} and
\lean{positivity}, Chapter~\ref{ch:01}); this is slightly narrower than Zhu's setting. No finiteness, family, cardinality or
separation hypothesis is assumed, and \lean{Pr} is supplied internally by \lean{VOA.pbw}.

\textbf{The three \lean{letI} lines.} They are not hypotheses; they \emph{construct} instances that the statement needs in
order to use $\Phi(V)$ as the index type of a \lean{ModularData}, which requires \lean{[Fintype ι] [DecidableEq ι]}. The
first is the \emph{proof} \lean{finite_Φ_of_C2Cofinite_pbw hC2} (Theorem~\ref{l15:thm:finitephi}); the second turns
finiteness into an enumeration (\lean{Fintype.ofFinite}); the third is classical decidable equality. Because they appear in
the statement, a user of the theorem sees exactly which instance is meant.

\textbf{Conclusion.} \lean{∃ D : ModularData W (Φ W), D.Mfam = canonicalIrrModule ∧ D.h = … ∧ ZhuModularityAnalytic D}.
The modular data (Definition~\ref{l15:def:modulardata}) are indexed by $\Phi(V)$ itself; their family is \emph{the} canonical
family and their weights are \emph{its} conformal weights (the two equations pin $D$, so the unpinned-variant trap of the
warning above does not apply); and \lean{ZhuModularityAnalytic D} is Zhu's law. By the structure \lean{ModularData},
$\rho=$ \lean{D.rep} is a monoid homomorphism, hence a representation, with \lean{rep_T} the diagonal matrix of
$e^{2\pi i(h_c-c_V/24)}$.

\textbf{Same as, stronger or weaker than Zhu?}
\begin{itemize}
\item \emph{Hypotheses}: $C_2$ and semisimplicity of $A(V)$, in place of $C_2$ and rationality. As a theorem about $V$ this is
  at least as general as Zhu's (Remark~\ref{l15:rem:hyp}). For a $V$ that is proved rational (\lean{AdmissiblyRational})
  the library applies it through \lean{isSemisimpleRing_}\allowbreak\lean{A_of_}\allowbreak\lean{admissiblyRational}; that is the headline,
  Theorem~\ref{l15:thm:rational}. (The first release revision of this chapter said that this step was not formalised. That
  was wrong.)
\item \emph{Which functions}: one-point functions $\tr_Mo(v)q^{L_0-c/24}$ for all $v$ of definite square-bracket weight,
  including the characters ($v=\lfifteenvac$). Not the $n$-point functions on the torus, which Zhu's theorem also treats.
\item \emph{Which modules}: all irreducible ordinary modules, one per isomorphism class, with no list to supply.
\item \emph{Convergence}: not in the conclusion; \lean{traceFun} uses \lean{tsum} (junk value $0$ on divergence). It is
  supplied by the separate theorems \lean{onePointSubexp_of_c2} and \lean{qSummableAt_of_subexponential}
  (Theorem~\ref{l15:thm:growth}); a separate Lean check, \lean{headline_}\allowbreak\lean{traces_are_genuine}, assembles them
  (not part of this release).
\item \emph{What is said about $\rho$}: existence, the homomorphism property and $\rho(T)$. Linear independence of the trace
  functions, hence uniqueness of $\rho$, is proved internally (Section~\ref{l15:sec:independence}) but is not exported in
  this statement. Nothing is said about $\rho(S)$ beyond existence.
\end{itemize}

\textbf{Could it be vacuous?} The hypotheses are satisfied by nontrivial vertex operator algebras
(Section~\ref{l15:sec:instances}) and by $\C\lfifteenvac$; the conclusion is pinned; the trace functions are genuine under
the hypotheses (checked in Lean); $\rho(T)$ is fixed by \lean{rep_T}.
\end{dictionary}

\begin{proof}[Proof on paper]
Apply Theorem~\ref{l15:thm:headline} (in its \lean{_voaPBW} form) to $\iota=\Phi(V)$, the canonical family and its
conformal weights, with the finiteness instance of Theorem~\ref{l15:thm:finitephi}. Its remaining hypotheses hold: $\#\iota
=\#\Phi(V)$ trivially (\lean{canonicalIrrFamily_card}); each $M_c$ has conformal weight $h_c$ by the choice of $h_c$
(\lean{canonicalIrrFamily_hasConformalWeight}); and if $c\neq c'$ have $h_c=h_{c'}$, then the bottoms of $M_c$ and $M_{c'}$
are not isomorphic, since otherwise the uniqueness half of Zhu's correspondence would give $M_c\cong M_{c'}$, i.e.\ $c=c'$
(\lean{canonicalIrrFamily_bottomSeparated}). The Lean proof is this composition, in \leanpath{ZhuModularityCanonical.lean}.
\end{proof}

\begin{physics}
This is the form in which a physicist states the theorem: \emph{the characters (and, more generally, the torus one-point
functions) of all the sectors of a $C_2$-cofinite chiral CFT whose zero-mode algebra is semisimple form a vector-valued
modular form for $\lfifteenSL$}; the $T$-matrix is the diagonal of $e^{2\pi i(h-c/24)}$ (the spins and the Casimir energy) and
the $S$-matrix exists. Semisimplicity of the zero-mode algebra on ground states is the algebraic statement that the
primary multiplets do not form indecomposable, non-split ``logarithmic'' pairs at the level of ground states; without it
(e.g.\ the triplet algebras) the characters alone are not closed under $S$, and one needs Miyamoto's pseudo-traces.
\end{physics}

\begin{status}
\lean{VOA.zhuModularity_}\allowbreak\lean{canonical_of_}\allowbreak\lean{C2Cofinite_semisimple}:
\texttt{[propext, Classical.choice, Quot.sound]}, run at an earlier commit (not re-run for this edition) on a scratch file with \texttt{lake env lean
-DmaxSynthPendingDepth=3}. The module \leanpath{ZhuModularityCanonical.lean} had no olean at that commit; the scratch file
therefore contained the file's text verbatim (its imports all had oleans, built 2026-10-25) followed by the
\texttt{\#print axioms} lines, so the theorem was elaborated and checked in that run. \textbf{\texttt{\#print
axioms}-clean, conditional on} \lean{C2Cofinite W} and \lean{IsSemisimpleRing (A W)}; uses the theorem \lean{VOA.pbw}; one-point
functions only. Per \texttt{Notes/RELEASE\_PLAN.md} (``Caveats removed in Lean'') the umbrellas containing it were built
green when it was merged; that build was not re-run for this chapter.
\end{status}

\subsection{The headline: Zhu's theorem for rational \texorpdfstring{$C_2$}{C2}-cofinite \texorpdfstring{$V$}{V}}\label{l15:sec:rational}

Zhu states his theorem for rational $C_2$-cofinite $V$. The engine of the previous subsection replaces ``rational'' by
``$A(V)$ is semisimple''. To return to Zhu's form one needs one more theorem, Zhu's semisimplicity theorem: rationality
implies that $A(V)$ is semisimple. The library proves it. The theorem lives in the Huang directory (Chapter~\ref{ch:L21}),
because it was proved there as a stage of the modular-tensor-category track. Since this chapter consumes it, we present it in
full here.

\begin{theorem}[Zhu's semisimplicity theorem; Zhu 1996, Theorem~2.2.3; Dong--Li--Mason]\label{l15:thm:zhuss}
Let $V$ be a vertex operator algebra of the library, and let the datum \lean{Pr} of Chapter~\ref{ch:L13} be given. If $V$ is
rational in the sense of \lean{AdmissiblyRational} (every admissible $V$-module is completely reducible;
Chapter~\ref{ch:L12}, Definition~\ref{l12:def:admrat}), then Zhu's algebra $A(V)$ is a semisimple ring. Since \lean{Pr} is
supplied for every $V$ by the theorem \lean{VOA.pbw}, the implication holds for every $V$; in particular the named predicate
\lean{ZhuAdmissibleSemisimple W} of Chapter~\ref{ch:L12} holds for every $V$.
\end{theorem}
\noindent File \leanpath{VOA/\allowbreak Huang/\allowbreak ZhuSemisimple.lean} (Chapter~\ref{ch:L21}), namespace
\lean{VOA.Huang.ZhuSemisimple}, with section variables \lean{{𝒱 : Type u} [AddCommGroup 𝒱] [Module ℂ 𝒱] {W : VertexOperatorAlgebra 𝒱}}
and \lean{(Pr : BottomProjection (parabolicRepOfAMod W (regularAMod W)))}; each theorem is preceded by the line
\lean{include Pr in}.
\begin{leancode}
theorem isSemisimpleRing_A_of_admissiblyRational (hrat : AdmissiblyRational W) :
    letI := instRingA W; IsSemisimpleRing (A W) := by

theorem zhuAdmissibleSemisimple_of_Pr : ZhuAdmissibleSemisimple W :=
  fun hrat => isSemisimpleRing_A_of_admissiblyRational Pr hrat
\end{leancode}

\begin{dictionary}
\textbf{The arguments.} The statement shows only \lean{hrat}, but the line \lean{include Pr in} makes the section variable
\lean{Pr} an explicit argument, placed first. So the full form is
\lean{isSemisimpleRing_}\allowbreak\lean{A_of_}\allowbreak\lean{admissiblyRational Pr hrat}, and this is how the headline below calls it, with
\lean{Pr} the term \lean{voaPBWBottomProjection} \lean{(parabolicRepOfAMod W (regularAMod W))}. \lean{Pr} is the bottom-projection datum of
Chapter~\ref{ch:L13}: the regular $A(V)$-module embeds as the bottom of the module induced from it. It is a consequence of
the Poincar\'e--Birkhoff--Witt theorem, which is the library theorem \lean{VOA.pbw}. So in effect the implication has no
hypothesis besides \lean{hrat}.

\textbf{The hypothesis and the conclusion.} \lean{AdmissiblyRational W} is Definition~\ref{l12:def:admrat} of
Chapter~\ref{ch:L12}: for every admissible module, every weak submodule has a weak complement. The conclusion
\lean{letI := instRingA W; IsSemisimpleRing (A W)} is literally the hypothesis \lean{hss} of the engine
(Theorem~\ref{l15:thm:canonical}): the same ring structure (\lean{instRingA W}, Zhu's product) and Mathlib's semisimplicity.
So the two theorems compose with no conversion.

\textbf{The named predicate.} \lean{ZhuAdmissibleSemisimple W} is the \lean{Prop}
\lean{AdmissiblyRational W → letI := instRingA W; IsSemisimpleRing (A W)} of Chapter~\ref{ch:L12}, Section~\ref{l12:sec:rationality}.
Earlier revisions of this book and of the library's docstrings called it an unproved bridge. \lean{zhuAdmissibleSemisimple_of_Pr}
proves it from \lean{Pr}, and with \lean{voaPBWBottomProjection} for every $V$. We checked this composition in the
scratch run of the status box below:
\begin{plaincode}
example : ∀ {𝒱 : Type} [AddCommGroup 𝒱] [Module ℂ 𝒱] (W : VOA.VertexOperatorAlgebra 𝒱),
    VOA.ZhuAdmissibleSemisimple W :=
  fun W => VOA.Huang.ZhuSemisimple.zhuAdmissibleSemisimple_of_Pr
    (VOA.VHat.DLM3.voaPBWBottomProjection
      (VOA.VHat.DLM3.parabolicRepOfAMod W (VOA.regularAMod W)))
\end{plaincode}
(This is a check, not a library declaration; the universe of \lean{𝒱} was fixed to \lean{Type} for the check.)

\textbf{Same as, stronger or weaker than the literature?} Zhu (Theorem~2.2.3, for his three-clause rationality) and
Dong--Li--Mason (for their one-clause rationality) prove: if $V$ is rational, $A(V)$ is semisimple. (The library's docstring
cites DLM, \emph{Math.\ Ann.} 310, Theorem~4.10(a); Chapter~\ref{ch:L12} cites DLM, \emph{J.~Algebra} 206, Theorem~4.10. We have
not resolved which numbering is meant; the statement is the same.) The Lean theorem has the same conclusion, with ``rational''
read as \lean{AdmissiblyRational}. DLM-rationality implies \lean{AdmissiblyRational}: in Lean for $C_2$-cofinite $V$
(Theorem~\ref{l12:thm:dlmrat}), on paper in general (Chapter~\ref{ch:L12}, Section~\ref{l12:warn:admcr}). So the Lean
theorem implies the literature's (for $C_2$-cofinite $V$ with no paper step): its hypothesis is
\emph{at most} as strong as DLM's. The Lean proof does not invoke any theorem of Dong--Li--Mason as a hypothesis, so the open
question of whether \lean{AdmissiblyRational} implies DLM-rationality does not affect it.

\textbf{Could it be vacuous?} Only if \lean{AdmissiblyRational} were unsatisfiable. It is not: the library proves it for the
trivial VOA and for every lattice VOA $V_{\sqrt{2m}\Z}$ (\lean{V2m_admissiblyRational}, Chapter~\ref{ch:L22}). And the
conclusion is not automatic: the triplet algebra $W(2)$ is $C_2$-cofinite, and the library proves that it is not
\lean{AdmissiblyRational} (Chapter~\ref{ch:L12}); on paper its Zhu algebra is not semisimple (Adamovi\'c--Milas), which is
not a Lean theorem.
\end{dictionary}

\begin{proof}[Proof on paper]
Let $U=A(V)$, viewed as the regular left module, and let $M=M(U)$ be the generalised Verma module induced from $U$
(Dong--Jiang). It is an admissible module whose degree-zero part is $U$; that the bottom is all of $U$, and not a quotient,
is the datum \lean{Pr}. Let $T\subseteq A(V)$ be a left ideal. Let $P\subseteq M$ be the submodule generated by $T\subseteq
M(0)$. By bottom recovery (Chapter~\ref{ch:L21}), $P$ is graded and $P\cap M(0)=T$. By rationality, $P$ has a complement $Q$
with $M=P\oplus Q$. Let $\pi_0$ be the projection onto degree zero, and put
\begin{equation*}
  T'=\{u\in A(V) : \exists q\in Q,\ \pi_0(q)=u,\ q-u\in P\}.
\end{equation*}
Zero modes preserve degrees, so they commute with $\pi_0$. They act on $M(0)=U$ by left multiplication in $A(V)$. Hence $T'$ is
a left ideal. $T\cap T'=0$: if $u\in T\cap T'$ with witness $q$, then $q=(q-u)+u\in P\cap Q=0$, so $u=\pi_0(q)=0$. $T+T'=A(V)$:
write $u=p+q$ with $p\in P$, $q\in Q$. Then $u':=\pi_0(q)$ lies in $T'$ with witness $q$ (indeed
$q-u'=\pi_0(p)-p\in P$, as $P$ is graded), and $u-u'=\pi_0(p)\in P\cap M(0)=T$.
So every left ideal has a complementary left ideal, i.e.\ $A(V)$ is a semisimple ring. The Lean proof is this argument, in
\leanpath{ZhuSemisimple.lean}.
\end{proof}

\begin{physics}
$A(V)$ is the algebra of zero modes acting on the ground states of the sectors. Rationality says that every representation of
the chiral algebra splits into irreducible sectors. The theorem says that this splitting is already visible on the ground
states: the zero-mode algebra is a direct sum of matrix algebras, one block per sector, of size the ground-state degeneracy.
A logarithmic theory, in which $L_0$ has Jordan blocks between sectors, fails both.
\end{physics}

\begin{status}
\lean{VOA.Huang.ZhuSemisimple.}\allowbreak\lean{isSemisimpleRing_}\allowbreak\lean{A_of_}\allowbreak\lean{admissiblyRational} and
\lean{VOA.Huang.ZhuSemisimple.}\allowbreak\lean{zhuAdmissibleSemisimple_of_Pr}: \texttt{[propext, Classical.choice, Quot.sound]}, run at
an earlier commit (not re-run for this edition) (scratch file, \texttt{lake env lean -DmaxSynthPendingDepth=3}; the olean of
\leanpath{ZhuSemisimple.lean} is newer than its source). \texttt{\#print axioms}-clean, \emph{conditional on} \lean{Pr} and
\lean{AdmissiblyRational W}; with \lean{Pr} from \lean{voaPBWBottomProjection} (also standard three in the same run),
conditional only on \lean{AdmissiblyRational W}.
\end{status}

\begin{theorem}[Zhu's modularity theorem for rational $C_2$-cofinite $V$; the chapter's headline]\label{l15:thm:rational}
Let $V$ be a vertex operator algebra of the library (so of CFT type) such that
\begin{enumerate}
\item[(a)] $V$ is $C_2$-cofinite, and
\item[(r)] $V$ is rational, in the sense of \lean{AdmissiblyRational}.
\end{enumerate}
Then the conclusion of Theorem~\ref{l15:thm:canonical} holds: $\Phi(V)$ is finite, and the one-point functions
$\tr_{M_c}o(v)q^{L_0-c_V/24}$ of the canonical irreducible modules $M_c$, $c\in\Phi(V)$, with their conformal weights $h_c$,
transform under $\lfifteenSL$ by a matrix representation $\rho$ with $\rho(T)=\mathrm{diag}(e^{2\pi i(h_c-c_V/24)})$, with
weight $n$ for $v\in V_{[n]}$. In particular this holds for every strongly rational $V$.
\end{theorem}
\noindent File \leanpath{VOA/\allowbreak Modularity/\allowbreak ZhuModularityCanonical.lean}.
\begin{leancode}
theorem zhuModularity_canonical_of_C2Cofinite_admissiblyRational (W : VertexOperatorAlgebra 𝒱)
    (hC2 : C2Cofinite W) (hrat : AdmissiblyRational W) :
    letI : Finite (Φ W) := VOA.Huang.ZhuFiniteIrr.finite_Φ_of_C2Cofinite_pbw hC2
    letI : Fintype (Φ W) := Fintype.ofFinite (Φ W)
    letI : DecidableEq (Φ W) := Classical.decEq (Φ W)
    ∃ D : ModularData W (Φ W), D.Mfam = canonicalIrrModule ∧
      D.h = (fun c => conformalWeightOf (canonicalIrrModule c)) ∧ ZhuModularityAnalytic D :=
  zhuModularity_canonical_of_C2Cofinite_semisimple W hC2
    (VOA.Huang.ZhuSemisimple.isSemisimpleRing_A_of_admissiblyRational
      (voaPBWBottomProjection (parabolicRepOfAMod W (regularAMod W))) hrat)

theorem zhuModularity_canonical_of_stronglyRational (W : VertexOperatorAlgebra 𝒱)
    (hSR : StronglyRational W) :
    letI : Finite (Φ W) := VOA.Huang.ZhuFiniteIrr.finite_Φ_of_C2Cofinite_pbw hSR.c2Cofinite
    letI : Fintype (Φ W) := Fintype.ofFinite (Φ W)
    letI : DecidableEq (Φ W) := Classical.decEq (Φ W)
    ∃ D : ModularData W (Φ W), D.Mfam = canonicalIrrModule ∧
      D.h = (fun c => conformalWeightOf (canonicalIrrModule c)) ∧ ZhuModularityAnalytic D :=
  zhuModularity_canonical_of_C2Cofinite_admissiblyRational W hSR.c2Cofinite hSR.rational
\end{leancode}

\begin{dictionary}
\textbf{Hypotheses.} Besides the ambient \lean{{𝒱 : Type u} [AddCommGroup 𝒱] [Module ℂ 𝒱]}:
\begin{itemize}
\item \lean{hC2 : C2Cofinite W}: (a), Definition~\ref{l12:def:c2}.
\item \lean{hrat : AdmissiblyRational W}: (r), Definition~\ref{l12:def:admrat}. This is Zhu's ``rational'', in the library's
  reading (weak complements). See the comparison below.
\item In the second form, \lean{hSR : StronglyRational W} (Definition~\ref{l12:def:strong}) bundles \lean{simple},
  \lean{selfDual}, \lean{c2Cofinite} and \lean{rational : AdmissiblyRational W}. The proof uses only \lean{hSR.c2Cofinite}
  and \lean{hSR.rational}. Simplicity and self-duality are not used, so this form is a special case of the first.
\end{itemize}
Nothing else is assumed: no finiteness of $\Phi(V)$, no family, no \lean{Pr}, and no named literature hypothesis.

\textbf{The proof terms are shown} because they are the whole proof. The first is the engine
\lean{zhuModularity_}\allowbreak\lean{canonical_of_}\allowbreak\lean{C2Cofinite_semisimple} applied to \lean{hC2} and to the semisimplicity proof
\lean{isSemisimpleRing_}\allowbreak\lean{A_of_}\allowbreak\lean{admissiblyRational Pr hrat} of Theorem~\ref{l15:thm:zhuss}, with \lean{Pr} the PBW datum. The
second is the first applied to the two relevant fields of \lean{hSR}.

\textbf{The \lean{letI} lines and the conclusion} are those of the engine, word for word, except that in the second form the
finiteness instance is built from \lean{hSR.c2Cofinite}. The conclusion is pinned: the modular data have the canonical family
and its conformal weights (Theorem~\ref{l15:thm:canonical}, dictionary).

\textbf{Same as, stronger or weaker than Zhu?} Compare with Theorem~\ref{l15:thm:lit}.
\begin{itemize}
\item \emph{Hypotheses}: Zhu's, ``rational and $C_2$-cofinite'', with rational read as \lean{AdmissiblyRational}, and CFT
  type built into \lean{VertexOperatorAlgebra}. Every $C_2$-cofinite DLM-rational $V$ is \lean{AdmissiblyRational}, by
  the Lean theorem \lean{admissiblyRational_of_dlmRational} (Theorem~\ref{l12:thm:dlmrat}), so the theorem covers every
  rational $C_2$-cofinite $V$ of CFT type with no paper step. Whether \lean{AdmissiblyRational} conversely implies DLM-rationality is open. That converse is not
  needed here, because the Lean proof uses no theorem of Dong--Li--Mason as a hypothesis.
\item \emph{Conclusion}: as for the engine: (2) and (3) of Theorem~\ref{l15:thm:lit} for one-point functions, with $\rho$ a
  representation and $\rho(T)$ the correct diagonal matrix. Convergence (1) is the separate growth theorem. $n$-point
  functions are not covered. Nothing is said about $\rho(S)$ beyond existence.
\end{itemize}

\textbf{Could it be vacuous?} No. The hypotheses hold for every lattice VOA $V_{\sqrt{2m}\Z}$ (Section~\ref{l15:sec:instances},
\lean{zhuModularity_V2m}, which is literally the second form applied to \lean{V2m_stronglyRational}) and for $\C\lfifteenvac$
(\lean{trivial_stronglyRational}). The conclusion is pinned, and the trace functions are genuine (checked in Lean,
\lean{headline_}\allowbreak\lean{traces_are_genuine}; the check is not part of this release).
\end{dictionary}

\begin{proof}[Proof on paper]
By Theorem~\ref{l15:thm:zhuss} (with \lean{Pr} from the PBW theorem), (r) gives (b): $A(V)$ is semisimple. Then (a) and (b) are
the hypotheses of the engine, Theorem~\ref{l15:thm:canonical}. For strongly rational $V$, (a) and (r) are two of the four
clauses of the definition.
\end{proof}

\begin{physics}
This is the theorem in the form a physicist quotes it: \emph{for a rational chiral CFT whose chiral algebra is finitely
generated in the $C_2$ sense (finitely many independent fields modulo derivatives and normal ordering), the characters of the
finitely many sectors, and all torus one-point functions of fields of definite weight, form a vector-valued modular form for
$\lfifteenSL$.} The $T$-matrix is the diagonal of the phases $e^{2\pi i(h-c/24)}$, the spins of the primaries shifted by the
Casimir energy, and the $S$-matrix exists. For the strongly rational case, which also asks for a unique vacuum and a
nondegenerate two-point function, this is the starting point of the modular bootstrap and of Huang's modular tensor category
(Chapter~\ref{ch:L21}).
\end{physics}

\begin{status}
\lean{VOA.zhuModularity_}\allowbreak\lean{canonical_of_}\allowbreak\lean{C2Cofinite_}\allowbreak\lean{admissiblyRational},
\lean{VOA.zhuModularity_}\allowbreak\lean{canonical_of_}\allowbreak\lean{stronglyRational}: \texttt{[propext, Classical.choice,
Quot.sound]}, run at an earlier commit (not re-run for this edition) on a scratch file with \texttt{lake env lean -DmaxSynthPendingDepth=3} from the
repository root. The olean of \leanpath{ZhuModularityCanonical.lean} predates the version of the file that contains these
theorems, so the scratch file contained the file's text verbatim (after its imports, whose oleans were present), followed by
the \texttt{\#print axioms} lines. \textbf{\texttt{\#print axioms}-clean, conditional on} \lean{C2Cofinite W} and
\lean{AdmissiblyRational W} (respectively on \lean{StronglyRational W}); uses the theorem \lean{VOA.pbw}; one-point functions
only. Per \texttt{Notes/RELEASE\_PLAN.md} (``Zhu for rational VOAs'', 2026-10-28) the umbrella builds of
\lean{VOA.Modularity} and \lean{VOA.Classification} containing them were green (12,271 jobs) with standard axioms. That build
was not re-run for this chapter, and no aggregate \texttt{lake build} was run.
\end{status}

\subsection{The strongly rational packaging}

\begin{theorem}[Zhu's theorem for strongly rational $V$, the older packaging with two (now proved) literature hypotheses]\label{l15:thm:sr}
Let $V$ be strongly rational (simple, self-dual, $C_2$-cofinite, rational; Definition~\ref{l12:def:strong}), and assume
the Dong--Li--Mason finiteness theorem (rational $\Rightarrow$ finitely many irreducible admissible modules) and Zhu's
semisimplicity theorem (rational $\Rightarrow$ $A(V)$ semisimple) for $V$. Then, with $\iota=\Phi(V)$ and the canonical
family of representatives with their conformal weights, the pinned conclusion of Theorem~\ref{l15:thm:headline} holds. The
proof uses the theorem \lean{VOA.pbw}.
\end{theorem}
\noindent Files \leanpath{VOA/\allowbreak Modularity/\allowbreak ZhuStronglyRational.lean},
\leanpath{VOA/\allowbreak Modularity/\allowbreak StronglyRationalFamily.lean}, \leanpath{VOA/\allowbreak Modularity/\allowbreak ConformalWeightChoice.lean}.
\begin{leancode}
noncomputable def canonicalIrrModule (c : Φ W) : IrrModule W := Quotient.out c

noncomputable def conformalWeightOf (V : IrrModule W) : ℂ :=
  (OrdVModule.exists_hasConformalWeight_of_irreducible V.isMod V.irreducible).choose

theorem zhuModularity_stronglyRational_voaPBW
    (hSR : StronglyRational W)
    (hFin : DLMFinitelyManyIrredAdmissible W)
    (hZhu : ZhuAdmissibleSemisimple W) :
    letI : Finite (Φ W) := finite_Φ_of_finite_Φadm W (hFin hSR.rational)
    letI : Fintype (Φ W) := Fintype.ofFinite (Φ W)
    letI : DecidableEq (Φ W) := Classical.decEq (Φ W)
    ∃ D : ModularData W (Φ W),
      D.Mfam = canonicalIrrModule ∧
      D.h = (fun c => conformalWeightOf (canonicalIrrModule c)) ∧
      ZhuModularityAnalytic D := by
\end{leancode}

\begin{dictionary}
\lean{Quotient.out c} chooses a representative of the class $c\in\Phi(V)$. \lean{conformalWeightOf} chooses the conformal weight
(it exists for irreducible ordinary modules and is unique). \lean{StronglyRational W} bundles
\lean{simple}, \lean{selfDual}, \lean{c2Cofinite}, \lean{rational : AdmissiblyRational W} (Chapter~\ref{ch:L12}).
\lean{DLMFinitelyManyIrredAdmissible W} is \lean{AdmissiblyRational W → Finite (Φadm W)} and \lean{ZhuAdmissibleSemisimple W} is
\lean{AdmissiblyRational W → IsSemisimpleRing (A W)} (Chapter~\ref{ch:L12}, Section~\ref{l12:sec:rationality}): both are
theorems of the literature for Dong--Li--Mason-rational $V$ (DLM 1998, Thm.~8.1; DLM 1998 \emph{J.~Algebra}, Thm.~4.10,
with Zhu 1996, Thm.~2.2.3 for Zhu's stronger notion), and in this declaration they are \emph{carried as hypotheses}. They
are, however, \textbf{theorems of the library} for every $V$ to which the declaration applies:
\begin{itemize}
\item \lean{hFin}: \lean{StronglyRational.dlmFinitelyManyIrredAdmissible hSR} (more generally
  \lean{dlmFinitelyManyIrredAdmissible_of_c2Cofinite} for any $C_2$-cofinite $V$; Chapter~\ref{ch:L12},
  Theorems~\ref{l12:thm:litproved} and~\ref{l12:thm:srlit}). It is proved from $C_2$-cofiniteness, the semisimplicity of
  $A(V)$ and the bottom projection, and the predicate's own \lean{AdmissiblyRational} antecedent. For an admissibly rational
  $V$ that is not $C_2$-cofinite the predicate is not proved, but such a $V$ is not strongly rational, so this declaration
  never needs that case.
\item \lean{hZhu}: \lean{zhuAdmissibleSemisimple_holds} (equivalently \lean{zhuAdmissibleSemisimple_of_Pr} applied to
  \lean{voaPBWBottomProjection}) proves \lean{ZhuAdmissibleSemisimple W} for every $V$ (Theorem~\ref{l15:thm:zhuss}).
\end{itemize}
Since both are proved directly from the library's \lean{AdmissiblyRational} and \lean{C2Cofinite}, their truth does not depend
on the comparison of \lean{AdmissiblyRational} with Dong--Li--Mason rationality (Chapter~\ref{ch:L12},
Section~\ref{l12:warn:admcr}). The library contains no restatement of this declaration with the two hypotheses discharged;
the discharge is the composition just described. \textbf{This packaging is superseded} anyway by the headline
\lean{zhuModularity_}\allowbreak\lean{canonical_of_}\allowbreak\lean{stronglyRational} (Theorem~\ref{l15:thm:rational}), which
reaches the same conclusion from \lean{hSR} alone, with finiteness of $\Phi(V)$ from Theorem~\ref{l15:thm:finitephi}. The \lean{letI} lines construct the finiteness, enumeration and decidability instances needed to use
$\Phi(V)$ as an index set. The other inputs of the headline are discharged: \lean{hzero} by \lean{le_of_eq W.vacuum_grade},
\lean{hcard} tautologically, \lean{hcw} by the choice of weights, \lean{hsimp} from semisimplicity and \lean{VOA.pbw}, \lean{hsep}
from the injectivity half of \lean{zhuBijective_of_voaPBW}. Simplicity and self-duality of $V$ are not used.
\end{dictionary}

\begin{proof}[Proof on paper]
Specialise Theorem~\ref{l15:thm:headline} to the canonical family; the family inputs are discharged as described.
\end{proof}

\begin{physics}
``Strongly rational'' is the class of VOAs that are chiral algebras of the nicest rational CFTs: unique vacuum, nondegenerate
two-point function, finitely many sectors, completely reducible representation theory. The theorem says that for these
the vector of characters of all sectors is a vector-valued modular form. In this older form the two classical structure
theorems (finitely many sectors; the zero-mode algebra is a sum of matrix blocks) are assumed as hypotheses, although the
library now proves both for these theories; the headline (Theorem~\ref{l15:thm:rational}) does not take them.
\end{physics}

\begin{status}
\lean{VOA.zhuModularity_}\allowbreak\lean{stronglyRational_voaPBW}: \texttt{[propext, Classical.choice, Quot.sound]} (re-probed
at a later commit, after \lean{VOA.pbw} became a theorem, not re-run for this edition; at the snapshot, when \lean{VOA.pbw} was an axiom, the list also contained \lean{VOA.pbw}). Conditional on
\lean{StronglyRational}, \lean{DLMFinitelyManyIrredAdmissible}, \lean{ZhuAdmissibleSemisimple}; uses the theorem
\lean{VOA.pbw}. Both literature hypotheses are now dischargeable for every strongly rational $V$:
\lean{ZhuAdmissibleSemisimple} by \lean{zhuAdmissibleSemisimple_holds} (or \lean{zhuAdmissibleSemisimple_of_Pr}), and
\lean{DLMFinitelyManyIrredAdmissible} by \lean{StronglyRational.dlmFinitelyManyIrredAdmissible} (proved for $C_2$-cofinite $V$;
\texttt{\#print axioms} of both: standard three, run at an earlier commit (not re-run for this edition), Chapter~\ref{ch:L12}, Theorem~\ref{l12:thm:litproved}).
\emph{Correction:} the previous revision called \lean{DLMFinitelyManyIrredAdmissible} ``unproved but unnecessary''; it is
proved for every $C_2$-cofinite $V$, which covers every strongly rational $V$. So the declaration, with its hypotheses
discharged, is unconditional given \lean{StronglyRational W}; the discharged form is a composition, not a named declaration.
It is still strictly worse than the headline \lean{zhuModularity_}\allowbreak\lean{canonical_of_}\allowbreak\lean{stronglyRational} and should not be
cited as the release form of Zhu's theorem.
\end{status}

\section{Non-vacuity: the trivial vertex operator algebra}\label{l15:sec:nonvac}

\texttt{\#print axioms} cannot detect an unsatisfiable hypothesis bundle. The library therefore instantiates the headline at
the simplest VOA, $V=\C\lfifteenvac$, realised as the lattice VOA of the rank-zero lattice.

\begin{theorem}[Zhu's theorem at $\C\lfifteenvac$]\label{l15:thm:trivial}
The trivial VOA $V=\C\lfifteenvac$ (central charge $0$) satisfies all hypotheses of Theorem~\ref{l15:thm:headline} with the
one-element family $\{V\}$, $h=0$; hence there are pinned modular data satisfying Zhu's law. $V$ is also strongly rational,
and the strongly rational packaging applies to it.
\end{theorem}
\noindent Files \leanpath{VOA/\allowbreak Examples/\allowbreak ZhuInstantiation.lean},
\leanpath{VOA/\allowbreak Examples/\allowbreak TrivialRationalityLattice.lean}, \leanpath{VOA/\allowbreak Examples/\allowbreak TrivialStronglyRational.lean}.
\begin{leancode}
noncomputable def trivialVOA : VertexOperatorAlgebra (LatGen.LatModuleGen (stdLattice 0)) :=
  LatGen.latticeVOAGen (stdLattice 0) std0_even std0_posDef

noncomputable def trivialAdjointIrr : IrrModule trivialVOA where
  carrier := LatGen.LatModuleGen (stdLattice 0)
  isMod := adjointModule trivialVOA
  irreducible := trivial_adjoint_irreducible

theorem trivialVOA_zhuModularity [Finite (Φ trivialVOA)]
    (hcard : Nat.card (Φ trivialVOA) = 1) :
    ∃ D : ModularData trivialVOA (Fin 1),
      D.Mfam = (fun _ => trivialAdjointIrr) ∧ D.h = (fun _ => 0) ∧
        ZhuModularityAnalytic D := by

theorem trivial_card_Φ : Nat.card (Φ trivialVOA) = 1 := by

theorem trivial_stronglyRational : StronglyRational trivialVOA :=
  stronglyRational_of_span_vacuum trivialVacuum_ne_zero trivial_span_vacuum
\end{leancode}

\begin{dictionary}
\lean{stdLattice 0} is the lattice $\Z^0=\{0\}$, even and positive definite vacuously; \lean{LatGen.latticeVOAGen} is the
lattice VOA construction (Chapter~\ref{ch:L05}), whose space is here one-dimensional. \lean{adjointModule} is $V$ as a module
over itself. \lean{Fin 1} is a one-element index set. The theorem \lean{trivialVOA_zhuModularity} still carries
\lean{[Finite (Φ trivialVOA)]} and \lean{hcard} (its docstring calls them ``the classification $\#\Phi=1$''), but
\lean{trivial_card_Φ} proves \lean{hcard}, and finiteness follows from \lean{trivial_subsingleton_Φ}. A separate Lean check
against the library (not part of this release) composes them, in \lean{trivialVOA_}\allowbreak\lean{zhuModularity_uncond} (the
pinned statement for $\C\lfifteenvac$ with the index set \lean{Fin 1} and $h=0$) and \lean{stronglyRational_}\allowbreak\lean{wrapper_fires}
(\lean{zhuModularity_stronglyRational_voaPBW} at \lean{trivialVOA}).
Both go through the \lean{_voaPBW} headline, i.e.\ through the theorem \lean{VOA.pbw}. So the hypothesis bundle of the
\lean{Pr}-form headline is satisfiable (unconditionally since \lean{VOA.pbw} became a theorem); \lean{Pr} for $\C\lfifteenvac$ is
supplied by \lean{voaPBWBottomProjection}.

\textbf{The closed form is now in the library.} After the release audit the library states the closed instance itself
(file \leanpath{VOA/\allowbreak Modularity/\allowbreak ZhuModularityCanonical.lean}):
\begin{leancode}
theorem trivialVOA_zhuModularity_closed :
    ∃ D : ModularData trivialVOA (Fin 1),
      D.Mfam = (fun _ => trivialAdjointIrr) ∧ D.h = (fun _ => 0) ∧ ZhuModularityAnalytic D := by
\end{leancode}
It has no hypotheses: its proof supplies \lean{[Finite (Φ trivialVOA)]} from
\lean{finite_Φ_of_C2Cofinite_pbw trivial_C2Cofinite} (Theorem~\ref{l15:thm:finitephi}) and \lean{hcard} from
\lean{trivial_card_Φ}. It is the library counterpart of that check's \lean{trivialVOA_}\allowbreak\lean{zhuModularity_uncond}, with
the same statement.
\end{dictionary}

\begin{warning}
\textbf{What the witness does and does not show.} For $V=\C\lfifteenvac$ the theorem says $\chi_V(\tau)=1$ is invariant under
$\lfifteenSL$ with $\rho=1$: true and trivial. It shows that the hypotheses are consistent and that the pinned conclusion is
attainable; it exercises none of the hard steps (the block space is one-dimensional, $\Phi$ has one element, the slash is
trivial). Nontrivial instances are the subject of the next section (Section~\ref{l15:sec:instances}); before this
release there were none. The Lee--Yang model is proved $C_2$-cofinite (\lean{leeYang_stronglyFinite},
Chapter~\ref{ch:L19}), but semisimplicity of its Zhu algebra is not established, so it is not an instance.
(Stale documentation: \lean{trivialVOA_zhuModularity}'s docstring says it is ``conditional only on the classification
$\#\Phi=1$''; that is now the theorem \lean{trivial_card_Φ}, and the closed form is \lean{trivialVOA_zhuModularity_closed}.)
\end{warning}

\begin{physics}
The trivial chiral algebra: a theory with one state, the vacuum. Its torus partition function is $1$, trivially modular
invariant. Its role is that of a consistency check, like verifying a formula on the empty set.
\end{physics}

\begin{status}
\lean{VOA.trivialVOA_}\allowbreak\lean{zhuModularity}: \texttt{[propext, Classical.choice, Quot.sound]} (re-probed at a later commit, after \lean{VOA.pbw} became a theorem, not re-run for this edition;
at the snapshot, when \lean{VOA.pbw} was an axiom, the list also contained \lean{VOA.pbw}). Theorems of the separate Lean check (not part of this release)
\lean{VOA.L15Cert.trivialVOA_}\allowbreak\lean{zhuModularity_uncond}, \lean{VOA.L15Cert.stronglyRational_}\allowbreak\lean{wrapper_fires}: the
standard three (check re-run at that later commit).
\lean{VOA.trivial_}\allowbreak\lean{stronglyRational}, \lean{VOA.trivial_card_Φ}: standard three.
\lean{VOA.trivialVOA_}\allowbreak\lean{zhuModularity_closed}: \texttt{[propext, Classical.choice, Quot.sound]}, run at
an earlier commit (not re-run for this edition) (same scratch run as Theorem~\ref{l15:thm:canonical}); no hypotheses.
\end{status}

\section{Nontrivial instances: the \texorpdfstring{$c=1$}{c=1} orbifolds and lattice VOAs}\label{l15:sec:instances}

The release audit observed that, as long as the only vertex operator algebra satisfying the hypotheses was $\C\lfifteenvac$,
the library's Zhu theorem had not been shown to say anything about a single nonconstant character. This section records
instances with nontrivial content. They come from the $c=1$ classification of Chapter~\ref{ch:L22}. For two families of
orbifolds of the compact boson, both hypotheses of the engine are proved directly, with no rationality
(Theorem~\ref{l15:thm:instances}). For the rank-one lattice VOAs, and for every VOA proved \emph{nice}, the headline applies
through strong rationality (Theorem~\ref{l15:thm:instancesV2m}).

\begin{definition}[The two vertex operator algebras]\label{l15:def:orbifolds}
(i) For $m\ge1$, $V_{2m}^+$ is the fixed-point subalgebra of the rank-one lattice VOA $V_{2m}=V_{\sqrt{2m}\Z}$ under the lift
$\theta$ of $-1$ (Chapter~\ref{ch:L22}, Definition~\ref{l22:def:V2m}). (ii) The \emph{tetrahedral model} is the fixed-point
subalgebra of $V_8^+$ (realised as the fixed points of the Klein four-group \lean{klein} acting on $V_2$, with
$V_2^{\lean{klein}}\cong V_8^+$, \lean{klein_iso_plus_four}) under an explicit automorphism \lean{kleinCubic} of order
dividing three. It is a model of the orbifold $V_2^{A_4}$ by the tetrahedral group.
\end{definition}
\noindent File \leanpath{VOA/\allowbreak Classification/\allowbreak COne/\allowbreak OrbRat2Reduction.lean} (Chapter~\ref{ch:L22}).
\begin{leancode}
abbrev cubicModel := orbifoldVOA (Subgroup.zpowers kleinCubic)
\end{leancode}

\begin{dictionary}
\lean{V2mPlus m hm} is quoted and decoded in Chapter~\ref{ch:L22} (\lean{orbifoldVOA (Subgroup.zpowers (thetaV2m m hm))});
\lean{orbifoldVOA G} is the fixed-point VOA $W^G$ (Chapter~\ref{ch:04}); \lean{Subgroup.zpowers g} is the cyclic group
$\langle g\rangle$. \lean{kleinCubic : Aut (orbifoldVOA klein)} (\leanpath{OrbC2bKlein.lean}) is the restriction to
$V_2^{\lean{klein}}$ of a rotation of $V_2$ of order three that normalises the Klein group; \lean{kleinCubic_cube} proves
$\lean{kleinCubic}^3=1$. The library's docstring describes \lean{cubicModel} as ``$(V_8^+)^{\Z/3}$''. That it is isomorphic
to $V_2^{A_4}$ is the content of the tetrahedral case of Chapter~\ref{ch:L22}; we use only the definition.
\end{dictionary}

\begin{physics}
$V_{2m}^+$ is the chiral algebra of the $c=1$ boson on the orbifold line $S^1/\Z_2$ at the radius where the momentum lattice is
$\sqrt{2m}\Z$; it has $m+7$ sectors (the untwisted sectors of the orbifold, plus the four $\theta$-twisted ones). The tetrahedral
model is one of the three exceptional isolated $c=1$ theories, the gauging of the $SU(2)_1$ WZW model by the tetrahedral
subgroup $A_4\subset SO(3)$ of its symmetry.
\end{physics}

\begin{status}
Definition.
\end{status}

\begin{theorem}[Zhu's theorem for $V_{2m}^+$ and for the tetrahedral model, with no hypotheses]\label{l15:thm:instances}
(i) For every $m\ge2$, the conclusion of Theorem~\ref{l15:thm:canonical} holds for $V=V_{2m}^+$: $\Phi(V_{2m}^+)$ is finite, and
the one-point functions of its canonical irreducible modules, with their conformal weights, transform under $\lfifteenSL$ by a
matrix representation with diagonal $\rho(T)$. (ii) The same holds for the tetrahedral model.
\end{theorem}
\noindent File \leanpath{VOA/\allowbreak Classification/\allowbreak COne/\allowbreak ZhuModularityV2mPlus.lean},
namespace \lean{VOA.Classification.COne}.
\begin{leancode}
theorem zhuModularity_V2mPlus (m : ℕ) (hm : 1 ≤ m) (hm2 : 2 ≤ m) :
    letI : Finite (Φ (V2mPlus m hm)) :=
      VOA.Huang.ZhuFiniteIrr.finite_Φ_of_C2Cofinite_pbw (V2mPlus_c2Cofinite m hm)
    letI : Fintype (Φ (V2mPlus m hm)) := Fintype.ofFinite _
    letI : DecidableEq (Φ (V2mPlus m hm)) := Classical.decEq _
    ∃ D : VOA.ModularData (V2mPlus m hm) (Φ (V2mPlus m hm)), D.Mfam = canonicalIrrModule ∧
      D.h = (fun c => conformalWeightOf (canonicalIrrModule c)) ∧ VOA.ZhuModularityAnalytic D :=
  zhuModularity_canonical_of_C2Cofinite_semisimple (V2mPlus m hm) (V2mPlus_c2Cofinite m hm)
    (V2mPlus_zhuSemisimple m hm hm2)

theorem zhuModularity_cubicModel :
    letI : Finite (Φ OrbRat2.cubicModel) :=
      VOA.Huang.ZhuFiniteIrr.finite_Φ_of_C2Cofinite_pbw OrbA4.cubicModel_c2Cofinite_holds
    letI : Fintype (Φ OrbRat2.cubicModel) := Fintype.ofFinite _
    letI : DecidableEq (Φ OrbRat2.cubicModel) := Classical.decEq _
    ∃ D : VOA.ModularData OrbRat2.cubicModel (Φ OrbRat2.cubicModel),
      D.Mfam = canonicalIrrModule ∧
      D.h = (fun c => conformalWeightOf (canonicalIrrModule c)) ∧ VOA.ZhuModularityAnalytic D :=
  zhuModularity_canonical_of_C2Cofinite_semisimple OrbRat2.cubicModel
    OrbA4.cubicModel_c2Cofinite_holds OrbTop.cubicModel_zhuSemisimple
\end{leancode}
\noindent The two inputs for $V_{2m}^+$ are, in \leanpath{VOA/\allowbreak Classification/\allowbreak COne/\allowbreak PlusC2Final.lean}
and \leanpath{VOA/\allowbreak Classification/\allowbreak COne/\allowbreak PlusTwistCharacters.lean}:
\begin{leancode}
theorem V2mPlus_c2Cofinite (m : ℕ) (hm : 1 ≤ m) : C2Cofinite (V2mPlus m hm) :=
  v2mPlusC2CofiniteFact_holds m hm

theorem V2mPlus_zhuSemisimple : ZhuV2mPlusSemisimple :=
  zhuSemisimple_of_fourExtraCharacters
    (fourExtraCharacters_of_twistedOmegaECharacters V2mPlus_twistedOmegaECharacters)

theorem V2mPlus_zhu_finrank (m : ℕ) (hm : 1 ≤ m) (hm2 : 2 ≤ m) :
    Module.finrank ℂ (VOA.A (V2mPlus m hm)) = m + 7 :=
\end{leancode}
\noindent where \lean{ZhuV2mPlusSemisimple} is defined in \leanpath{VOA/\allowbreak Classification/\allowbreak COne/\allowbreak PlusSplitResidual.lean}:
\begin{leancode}
def ZhuV2mPlusSemisimple : Prop :=
  ∀ (m : ℕ) (hm : 1 ≤ m), 2 ≤ m →
    letI := instRingA (V2mPlus m hm)
    IsSemisimpleRing (VOA.A (V2mPlus m hm))
\end{leancode}
\noindent and the two inputs for the tetrahedral model are, in \leanpath{VOA/\allowbreak Classification/\allowbreak COne/\allowbreak OrbA4Assembly.lean}
and \leanpath{VOA/\allowbreak Classification/\allowbreak COne/\allowbreak OrbTopZhu.lean}:
\begin{leancode}
theorem cubicModel_c2Cofinite_holds : C2Cofinite cubicModel := by

theorem cubicModel_zhuSemisimple : IsSemisimpleRing (VOA.A cubicModel) :=
  OrbRat2.isSemisimpleRing_of_charTable OrbA4.cubicModel_c2Cofinite_holds
    OrbSgl.cubicZhuDimensionBound_holds OrbRat6.cubicZhuCharacterTable_holds
\end{leancode}

\begin{dictionary}
\textbf{Reading the statements.} Each is Theorem~\ref{l15:thm:canonical} applied to a specific $V$; the proof terms (shown,
since they are informative) are literally \lean{zhuModularity_}\allowbreak\lean{canonical_of_}\allowbreak\lean{C2Cofinite_semisimple}
applied to a proof of $C_2$-cofiniteness and a proof of semisimplicity of $A(V)$. The \lean{letI} lines are the same
instance constructions as in the headline. \lean{VOA.ModularData}, \lean{VOA.ZhuModularityAnalytic} are the objects of
Definition~\ref{l15:def:modulardata} (the prefix \lean{VOA.} only because the file lives in a sub-namespace).

\textbf{Hypotheses.} \lean{zhuModularity_cubicModel} has none. \lean{zhuModularity_V2mPlus} has only the parameters $m$ and
the bounds \lean{hm : 1 ≤ m} (needed to form $V_{2m}^+$ at all, since the lattice $\sqrt{2m}\Z$ must be positive definite)
and \lean{hm2 : 2 ≤ m} (needed for the semisimplicity input; $V_2^+\cong V_8$ is a lattice VOA, covered up to isomorphism by
Theorem~\ref{l15:thm:instancesV2m}, although the modularity statement is not transported along that isomorphism). These are
ranges of a parameter, not assumptions. As \lean{hm2} implies \lean{hm}, the latter is redundant; the library keeps it only
because \lean{V2mPlus m hm} takes it as an argument.

\textbf{The inputs.} \lean{V2mPlus_c2Cofinite}: $C_2$-cofiniteness of $V_{2m}^+$ for all $m\ge1$ (Yamskulna; Abe--Buhl--Dong;
proved in Chapter~\ref{ch:L22} by an explicit strong-generation argument). \lean{V2mPlus_zhuSemisimple}: semisimplicity of
$A(V_{2m}^+)$ for $m\ge2$ (Abe, Theorem~7.1(1)), proved from the characters of $A(V_{2m}^+)$ on the tops of the four
$\theta$-twisted modules, which the library constructs from actual lattice modules
(\lean{V2mPlus_twistedOmegaECharacters}) and which complete the list needed for semisimplicity; the companion \lean{V2mPlus_zhu_finrank} computes $\dim A(V_{2m}^+)=m+7$. For the
tetrahedral model, \lean{cubicModel_zhuSemisimple} follows from a dimension bound $\dim A(V)\le21$ and a table of $21$
distinct characters of $A(V)$, which make the evaluation map $A(V)\to\C^{21}$ an algebra isomorphism; so $A(V)\cong\C^{21}$
is commutative and semisimple.

\textbf{What is \emph{not} used.} Neither instance passes through rationality: the semisimplicity of $A(V)$ is proved directly.
So these instances illustrate Remark~\ref{l15:rem:hyp}: the hypothesis the library needs can be checked without first
proving that $V$ is rational.

\textbf{(iii) The lattice VOAs $V_{2m}$ go through rationality instead.} $V_{2m}=V_{\sqrt{2m}\Z}$ is strongly rational in
the library (\lean{V2m_stronglyRational}, Chapter~\ref{ch:L22}). So Zhu's theorem for it follows from the headline, with
semisimplicity of $A(V_{2m})$ supplied by Theorem~\ref{l15:thm:zhuss}; this is Theorem~\ref{l15:thm:instancesV2m} below.
(The first release revision of this chapter said here that $V_{2m}$ was not covered, because ``rational $\Rightarrow A(V)$
semisimple'' was unproved. That was wrong; see the correction in the findings.) $\Vnat$ and the lattice VOAs of higher rank
are still not covered: although every even positive-definite lattice VOA is proved $C_2$-cofinite
(\lean{latticeVOAGen_c2Cofinite}, Chapter~\ref{ch:L22}), the library proves \lean{AdmissiblyRational} (or semisimplicity of
$A(V)$) only in rank one; for $\Vnat$ neither input is proved.

\textbf{Could the instances be empty of content?} No: unlike $\C\lfifteenvac$, $V_{2m}^+$ has $m+7$ inequivalent irreducible
modules, with conformal weights that include $r^2/4m$ ($1\le r\le m$) and the twisted weights $1/16$, $9/16$, so the vector
of characters is a genuinely multi-component vector-valued modular function, and $\rho(S)$ (which on paper mixes all the
sectors) is not a permutation of a trivial representation.
\end{dictionary}

\begin{proof}[Proof on paper]
Immediate from Theorem~\ref{l15:thm:canonical} and the two inputs for each $V$. The inputs themselves are theorems of
Chapter~\ref{ch:L22}.
\end{proof}

\begin{physics}
For the orbifold boson at radius $\sqrt{2m}$ the theorem says that the $m+7$ characters --- the two untwisted characters
$(\chi_{V_{2m}}\pm\chi')/2$ of the vacuum and its partner, the $m-1$ charged sectors, the two sectors at the half-lattice
point, and the four twisted characters built from $\eta(\tau)/\eta(\tau/2)$-type expressions --- span a representation of
$\lfifteenSL$, as they must for the orbifold partition function on the torus to be modular invariant. The theorem gives the
representation's existence and its $T$-matrix; it does not compute the $S$-matrix (Dijkgraaf--Vafa--Verlinde--Verlinde
did, by hand). The tetrahedral model is the analogous statement for an isolated exceptional $c=1$ theory with
$21$ sectors.
\end{physics}

\begin{status}
\lean{VOA.Classification.COne.zhuModularity_V2mPlus}, \lean{VOA.Classification.COne.zhuModularity_cubicModel}: both
\texttt{[propext, Classical.choice, Quot.sound]}, run at an earlier commit (not re-run for this edition) (the file had no olean; its text was elaborated
in the scratch run together with \leanpath{ZhuModularityCanonical.lean}). Inputs \lean{V2mPlus_c2Cofinite},
\lean{V2mPlus_zhuSemisimple}, \lean{V2mPlus_zhu_finrank}, \lean{OrbA4.cubicModel_c2Cofinite_holds},
\lean{OrbTop.cubicModel_zhuSemisimple}: the standard three, same run. \textbf{Unconditional} (no hypotheses beyond the range
of $m$; the full hypothesis closure is the two proved inputs), using the theorem \lean{VOA.pbw}. Content: one-point
functions only; convergence by the separate growth theorem; $\rho(S)$ only shown to exist. ($V_{2m}$ is covered by the next
theorem; $\Vnat$ and the higher-rank lattice VOAs are not covered: the latter are $C_2$-cofinite
(\lean{latticeVOAGen_c2Cofinite}) but not proved rational.)
\end{status}

\begin{theorem}[Zhu's theorem for the lattice VOAs $V_{\sqrt{2m}\Z}$, and for every nice VOA]\label{l15:thm:instancesV2m}
(i) For every $m\ge1$, the conclusion of Theorem~\ref{l15:thm:canonical} holds for the rank-one lattice VOA
$V=V_{2m}=V_{\sqrt{2m}\Z}$, with no further hypotheses. (ii) It holds for every \emph{nice} vertex operator algebra $V$, i.e.\
every $V$ that is strongly rational and pseudo-unitary (Gannon--Rayhaun, Definition~2.6).
\end{theorem}
\noindent File \leanpath{VOA/\allowbreak Classification/\allowbreak COne/\allowbreak ZhuModularityV2mPlus.lean},
namespace \lean{VOA.Classification.COne}.
\begin{leancode}
theorem zhuModularity_V2m (m : ℕ) (hm : 1 ≤ m) :
    letI : Finite (Φ (V2m m hm)) :=
      VOA.Huang.ZhuFiniteIrr.finite_Φ_of_C2Cofinite_pbw (V2m_stronglyRational m hm).c2Cofinite
    letI : Fintype (Φ (V2m m hm)) := Fintype.ofFinite _
    letI : DecidableEq (Φ (V2m m hm)) := Classical.decEq _
    ∃ D : VOA.ModularData (V2m m hm) (Φ (V2m m hm)), D.Mfam = canonicalIrrModule ∧
      D.h = (fun c => conformalWeightOf (canonicalIrrModule c)) ∧ VOA.ZhuModularityAnalytic D :=
  zhuModularity_canonical_of_stronglyRational (V2m m hm) (V2m_stronglyRational m hm)

theorem zhuModularity_of_nice {𝒱 : Type*} [AddCommGroup 𝒱] [Module ℂ 𝒱]
    (V : VertexOperatorAlgebra 𝒱) (hV : Nice V) :
    letI : Finite (Φ V) := VOA.Huang.ZhuFiniteIrr.finite_Φ_of_C2Cofinite_pbw hV.1.c2Cofinite
    letI : Fintype (Φ V) := Fintype.ofFinite _
    letI : DecidableEq (Φ V) := Classical.decEq _
    ∃ D : VOA.ModularData V (Φ V), D.Mfam = canonicalIrrModule ∧
      D.h = (fun c => conformalWeightOf (canonicalIrrModule c)) ∧ VOA.ZhuModularityAnalytic D :=
  zhuModularity_canonical_of_stronglyRational V hV.1
\end{leancode}
\noindent The objects in these statements are, in \leanpath{VOA/\allowbreak Classification/\allowbreak COne/\allowbreak Statement.lean}
and \leanpath{VOA/\allowbreak Classification/\allowbreak COne/\allowbreak LatticeStronglyRational.lean} (Chapter~\ref{ch:L22}):
\begin{leancode}
def Nice (V : VertexOperatorAlgebra 𝒱) : Prop :=
  StronglyRational V ∧ PseudoUnitary V

noncomputable def V2m (m : ℕ) (hm : 1 ≤ m) : VertexOperatorAlgebra (LatModuleGen (latP m)) :=
  latticeVOAGen (latP m) (latP_even m) (latP_posDef m hm)

theorem V2m_stronglyRational (m : ℕ) (hm : 1 ≤ m) : StronglyRational (V2m m hm) :=
  { V2m_stronglyFinite m hm with rational := V2m_admissiblyRational hm }
\end{leancode}

\begin{dictionary}
\textbf{The objects.} \lean{V2m m hm} is the lattice VOA (\lean{latticeVOAGen}, Chapter~\ref{ch:L05}) of the rank-one lattice
\lean{latP m} $=\sqrt{2m}\,\Z$, which is even (\lean{latP_even}) and, for $m\ge1$, positive definite (\lean{latP_posDef m hm});
its central charge is $1$. Chapter~\ref{ch:L22}, Definition~\ref{l22:def:V2m}, decodes it. \lean{Nice V} is the conjunction
(\lean{∧}) of \lean{StronglyRational V} and \lean{PseudoUnitary V}: every irreducible module not isomorphic to $V$ has all its
conformal weights of positive real part (Chapter~\ref{ch:L22}, Definitions~\ref{l22:def:pu} and~\ref{l22:def:nice};
Chapter~\ref{ch:L21}). \lean{hV.1} is the first component, the strong rationality. \lean{V2m_stronglyRational} combines strong
finiteness (simple, self-dual, $C_2$-cofinite; \lean{V2m_stronglyFinite}) with \lean{V2m_admissiblyRational}, the library's
proof that every admissible $V_{2m}$-module is completely reducible (Dong's theorem for lattice VOAs, in the library's reading of
rationality).

\textbf{Reading the statements.} The proof terms are shown because they are the whole proof: each is the headline form
\lean{zhuModularity_}\allowbreak\lean{canonical_of_}\allowbreak\lean{stronglyRational} (Theorem~\ref{l15:thm:rational}) applied to
a proof of strong rationality. The \lean{letI} lines are the same instance constructions as in the headline, with the
$C_2$-cofiniteness taken from the strong-rationality proof. The conclusion is pinned to the canonical family and its weights.

\textbf{Hypotheses.} \lean{zhuModularity_V2m} has only the parameter $m$ and the bound \lean{hm : 1 ≤ m}, which is needed to form
$V_{2m}$ at all (positive definiteness). It is \textbf{unconditional}. \lean{zhuModularity_of_nice} has the hypothesis
\lean{hV : Nice V}, of which only strong rationality is used; pseudo-unitarity is not. So it is the special case of
Theorem~\ref{l15:thm:rational} for nice $V$, stated because the classification and Huang chapters work with nice VOAs. Every
VOA that the library proves nice is an unconditional instance, for example $V_{2m}$ (\lean{V2m_nice}) and $V_{2m}^+$
(\lean{V2mPlus_nice}, every $m\ge1$, which also covers $m=1$, excluded from \lean{zhuModularity_V2mPlus}).

\textbf{What is used, and what is not.} Unlike the orbifold instances of Theorem~\ref{l15:thm:instances}, these do not prove
semisimplicity of $A(V)$ by a computation with the Zhu algebra. They go through rationality and
\lean{isSemisimpleRing_}\allowbreak\lean{A_of_}\allowbreak\lean{admissiblyRational}. So the two kinds of instance check the two layers of the theorem: the engine
directly, and the headline through Zhu's semisimplicity theorem.

\textbf{Could the instances be empty of content?} No. $V_{\sqrt{2m}\Z}$ has $2m$ inequivalent irreducible modules
$V_{\lambda+\sqrt{2m}\Z}$, $\lambda=r/\sqrt{2m}$, $r\in\Z/2m$, with conformal weights $\min(r,2m-r)^2/4m$ for $0\le r<2m$
(on paper; Dong's classification, not a Lean theorem here). Their characters
$\theta_{r}(\tau)/\eta(\tau)$, $\theta_r=\sum_{n\in\Z}q^{(2mn+r)^2/4m}$, form a genuinely $2m$-dimensional vector-valued modular
function, and $\rho(S)$ is, on paper and up to a sign convention, the discrete Fourier transform $e^{-2\pi i rs/2m}/\sqrt{2m}$
on $\Z/2m$.
\end{dictionary}

\begin{proof}[Proof on paper]
(i) $V_{2m}$ is strongly rational (\lean{V2m_stronglyRational}); apply Theorem~\ref{l15:thm:rational}. (ii) A nice VOA is
strongly rational by definition; apply Theorem~\ref{l15:thm:rational}.
\end{proof}

\begin{physics}
$V_{\sqrt{2m}\Z}$ is the chiral algebra of the compact boson at the rational radius where the chiral momentum lattice is
$\sqrt{2m}\Z$; $m=1$ is the self-dual $SU(2)_1$ point. Its sectors are labelled by the momentum $r\in\Z/2m$, and the theorem
says that the $2m$ characters $\theta_r/\eta$ span a representation of $\lfifteenSL$. This is the rational point of the
$c=1$ moduli space, where the chiral algebra is extended by the vertex operators $e^{\pm i\sqrt{2m}\phi}$ and the torus partition
function becomes a finite sum of products of these characters. Pseudo-unitarity, which (ii) assumes but does not use, says
that every non-vacuum sector has positive energy, as in a unitary CFT.
\end{physics}

\begin{status}
\lean{VOA.Classification.COne.}\allowbreak\lean{zhuModularity_V2m}, \lean{VOA.Classification.COne.}\allowbreak\lean{zhuModularity_of_nice}: \texttt{[propext,
Classical.choice, Quot.sound]}, run at an earlier commit (not re-run for this edition) in the same scratch run as Theorem~\ref{l15:thm:rational} (the olean of
\leanpath{ZhuModularityV2mPlus.lean} also predates the theorems, so the file's text was elaborated in the scratch file).
Inputs \lean{V2m_stronglyRational}, \lean{V2m_nice}, \lean{V2mPlus_nice}: the standard three, same run.
\lean{zhuModularity_V2m}: \textbf{unconditional} (no hypotheses beyond the range of $m$), using the theorem \lean{VOA.pbw}.
\lean{zhuModularity_of_nice}: \texttt{\#print axioms}-clean, \emph{conditional on} \lean{Nice V}. One-point functions only.
Per \texttt{Notes/RELEASE\_PLAN.md} the umbrella build containing them was green (12,271 jobs); not re-run here.
\end{status}

\section[Twisted modules: a statement only]{Twisted modules and the orbifold conjecture: a statement only}\label{l15:sec:twisted}

Zhu's theorem has an orbifold generalisation (Dong--Li--Mason 2000): for a finite group $G$ of automorphisms of a rational
$C_2$-cofinite $V$, the twisted trace functions $\lfifteenZ(M,h;\tau)=\tr_M\varphi(h)q^{L_0-c/24}$ over $g$-twisted modules $M$
($g,h\in G$ commuting) span an $\lfifteenSL$-representation, and each is invariant under a congruence subgroup. The library
states (part of) this as a proposition and never asserts it.

\begin{definition}[Twisted trace; Conjecture 5.3.10 as a proposition]\label{l15:def:conj}
For a $g$-twisted module $T$ ($g$ of order $N$) with $L_0$-weights in $h_0+\frac1N\N$ and an operator $\varphi$ commuting with
the twisted action, the twisted trace is the formal series $\sum_e\tr(\varphi|_{T_{e+c/24}})q^e$. The proposition asserts that,
evaluated along $h_0+\frac1N\N$, it is invariant under $\Gamma(N')$ for some $N'>0$, for all such $g,T,h_0,\varphi$.
\end{definition}
\noindent File \leanpath{VOA/\allowbreak Modularity/\allowbreak Twisted.lean}.
\begin{leancode}
noncomputable def twistedTrace {g : Aut W} {N : ℕ} {M : Type u} [AddCommGroup M] [Module ℂ M]
    (T : TwistedModule g N M) (φ : Module.End ℂ M) : FormalQSeries :=
  fun e =>
    haveI := T.finrank (e + W.centralCharge / 24)
    LinearMap.trace ℂ _ (T.weightBlock φ (e + W.centralCharge / 24))

noncomputable def qEvalAtFrac (χ : FormalQSeries) (h₀ : ℂ) (N : ℕ) (τ : UpperHalfPlane) : ℂ :=
  Complex.exp (2 * (Real.pi : ℂ) * Complex.I * (τ : ℂ) * h₀) *
    ∑' n : ℕ, χ (h₀ + (n : ℂ) / N) * Complex.exp (2 * (Real.pi : ℂ) * Complex.I * (τ : ℂ) * ((n : ℂ) / N))

def TwistedTraceModular {g : Aut W} {N : ℕ} {M : Type u} [AddCommGroup M] [Module ℂ M]
    (T : TwistedModule g N M) (φ : Module.End ℂ M) (h₀ : ℂ) : Prop :=
  ∃ N' : ℕ, 0 < N' ∧ ∀ A ∈ CongruenceSubgroup.Gamma N', ∀ τ : UpperHalfPlane,
    qEvalAtFrac (twistedTrace T φ) h₀ N (A • τ) = qEvalAtFrac (twistedTrace T φ) h₀ N τ

def Conjecture_5_3_10 (W : VertexOperatorAlgebra 𝒱) : Prop :=
  ∀ (g : Aut W) (N : ℕ), 0 < N → g ^ N = 1 →
    ∀ (M : Type u) [AddCommGroup M] [Module ℂ M] (T : TwistedModule g N M) (h₀ : ℂ),
      (∀ α : ℂ, T.weight α ≠ ⊥ → ∃ n : ℕ, α = h₀ + (n : ℂ) / N) →
      ∀ φ : Module.End ℂ M, (∀ (u : 𝒱) (k : ℤ), φ * T.YM u k = T.YM u k * φ) →
        TwistedTraceModular T φ h₀
\end{leancode}

\begin{dictionary}
\lean{TwistedModule g N M} is the twisted-module structure of \leanpath{Twisted.lean} (fractional modes \lean{YM u k}
$=u_{(k/N)}$, with the repaired Jacobi identity restricted to the support coset; Chapter~\ref{ch:06} explains the
repair). \lean{CongruenceSubgroup.Gamma N'} is Mathlib's $\Gamma(N')$. \lean{T.weight α ≠ ⊥ → …} says the $L_0$-weights lie in
$h_0+\frac1N\N$. The proposition takes no hypothesis on $V$ (no $C_2$, no rationality); it is meant to be assumed for, or
proved for, suitable $V$. It is used nowhere.
\end{dictionary}

\begin{warning}
\textbf{\lean{Conjecture_5_3_10} reads the twisted trace along the wrong exponents (serious).} The coefficient of
\lean{twistedTrace T φ} at $e$ is the trace on the weight space $T_{e+c/24}$ (exponent convention of
Definition~\ref{l15:def:torusonepoint}). The hypothesis puts the weights in $h_0+\frac1N\N$, so the nonzero coefficients sit at
$e=h_0-\frac c{24}+\frac nN$. But \lean{qEvalAtFrac (twistedTrace T φ) h₀ N} reads the coefficients at $e=h_0+\frac nN$, i.e.\ the
traces on $T_{h_0+n/N+c/24}$. (Compare \lean{traceFun}, which correctly evaluates along $h_i-c/24$.) Consequences:
\begin{itemize}
\item If $\frac c{24}\notin\frac1N\Z$ (e.g.\ $c=\frac12$, $N=1$), every coefficient read is $0$, the evaluated function is $0$,
  and \lean{TwistedTraceModular} holds trivially: the ``conjecture'' is vacuous for every twisted module of every such $V$.
  Checked in Lean against the library (the check is not part of this release): \lean{VOA.L15Cert.twistedTraceModular_}\allowbreak\lean{of_weight_bot}
  and \lean{twistedTraceModular_}\allowbreak\lean{of_c_off_lattice} (standard three axioms).
\item If $\frac c{24}\in\frac1N\Z$, the function read is a truncated, shifted trace and is in general \emph{not} modular. For
  example, for $V=\Vnat$, $g=1$, $N=1$, $T=\Vnat$, $\varphi=1$, $h_0=0$ ($c=24$) one reads
  $\sum_n\dim\Vnat_{n+1}q^n=J(\tau)-q^{-1}$, which is invariant under no $\Gamma(N')$ (on paper; the genuine statement, with
  $h_0-c/24$, gives $J$, which is invariant).
\end{itemize}
So the proposition is not Dong--Li--Mason's conjecture/theorem. Fix: evaluate along $h_0-\frac c{24}$, i.e.\ use
\lean{qEvalAtFrac (twistedTrace T φ) (h₀ - W.centralCharge / 24) N}. Two further, milder deviations: the proposition does not
ask $\varphi$ to come from $h\in C_G(g)$ (any commuting operator is allowed; harmless for irreducible $T$ by Schur), and it
omits the $\lfifteenSL$-covariance of the whole family (5.3.25), as its docstring says.
\end{warning}

\begin{physics}
Twisted sectors are the sectors of an orbifold theory: strings closed up to the symmetry $g$. The twisted trace with $h$
inserted is the torus partition function with boundary conditions twisted by $g$ along one cycle and $h$ along the other,
and modular transformations permute the pairs $(g,h)$. This is how $\Vnat$ arises from the Leech lattice theory and how
Monstrous Moonshine's McKay--Thompson series become modular functions. The library contains this only as a statement.
\end{physics}

\begin{status}
\lean{VOA.Conjecture_5_3_10}: a definition (\texttt{\#print axioms}: standard three). Nothing proves or uses it. The
theorems of the separate Lean check about it (not part of this release): standard three axioms.
\end{status}

\section{What is not proved}\label{l15:sec:notproved}

For the reader comparing with textbook accounts of Zhu's theorem (Gannon \S5.3; Dong--Li--Mason; Huang), the following
standard companions are \emph{not proved about the representation $\rho$ of this chapter}. (Chapter~\ref{ch:L21} develops, for
the class of VOAs it calls \emph{nice} --- strongly rational and pseudo-unitary --- Huang's modular tensor category, a
Verlinde theorem and congruence statements; their hypotheses and their relation to this chapter's $\rho$ are discussed
there, not here.)
\begin{itemize}
\item properties of $\rho(S)$: symmetry, unitarity, $S^2=C$ (charge conjugation), $(ST)^3=S^2$ beyond what $\rho$ being a
  homomorphism implies; the \emph{Verlinde formula} (Huang 2005), and indeed any nonzero fusion coefficient (the library's
  \lean{fusionCoeff} is identically $0$, a placeholder rather than the fusion rules; Chapter~\ref{ch:L12});
\item the congruence property ($\ker\rho\supseteq\Gamma(N)$; Dong--Lin--Ng 2015) for general $V$: for nice $V$ (strongly
  rational and pseudo-unitary) it is proved for Zhu's representation in the form of a Zhu datum
  (\lean{Huang.}\allowbreak\lean{CongruenceNamed.}\allowbreak\lean{zhu_kernelContainsGamma_at_level}, Chapter~\ref{ch:L21}; it uses
  only the kernel half of Ng--Schauenburg); and rationality of $c$ and $h_i$;
\item a modular linear differential equation with modular-form coefficients (Section~\ref{l15:sec:mlde});
\item the theorem for torus $n$-point functions, and Miyamoto's extension to non-rational $C_2$-cofinite $V$ via pseudo-traces
  (the logarithmic files are Chapter~\ref{ch:L18});
\item the orbifold theorem (Section~\ref{l15:sec:twisted});
\item the converse half of the comparison of the library's rationality \lean{AdmissiblyRational} with Dong--Li--Mason
  rationality. (The direction DLM $\Rightarrow$ \lean{AdmissiblyRational}, needed to apply the headline to a $V$ known
  only to be rational in the paper's sense, is now proved for $C_2$-cofinite $V$, \lean{admissiblyRational_of_dlmRational},
  Theorem~\ref{l12:thm:dlmrat}, which is all the headline needs.) For the converse, Lean proves that
  \lean{AdmissiblyRational} implies the level-graded reading (\lean{admissiblyRationalGraded_of_admissiblyRational}) and,
  under $C_2$, DLM's decomposition into irreducible ordinary modules
  (\lean{AdmissiblyRational.}\allowbreak\lean{exists_}\allowbreak\lean{iSupIndep_}\allowbreak\lean{irreducible_}\allowbreak\lean{ordinary}),
  but not DLM's definition itself (Remark~\ref{l15:rem:hyp});
\item instances for $\Vnat$ and for lattice VOAs of rank $>1$: the latter are proved $C_2$-cofinite
  (\lean{latticeVOAGen_c2Cofinite}) but not rational, and for $\Vnat$ neither is proved (Section~\ref{l15:sec:instances});
\item the equality $\dim\lfifteenB=\#\Phi(V)$ (only $\le$ is a Lean theorem; Theorem~\ref{l15:thm:533}).
\end{itemize}
(Four items of earlier versions of this list are now proved, or were proved all along: nontrivial $V$ satisfying the
hypotheses, Section~\ref{l15:sec:instances}; an axiom-free \lean{Pr}, supplied by the theorem \lean{VOA.pbw}; the
implication ``rational $\Rightarrow A(V)$ semisimple'', \lean{isSemisimpleRing_}\allowbreak\lean{A_of_}\allowbreak\lean{admissiblyRational}
(Theorem~\ref{l15:thm:zhuss}), which the first release revision wrongly listed here as unproved; and with it Zhu's theorem for
the lattice VOAs $V_{\sqrt{2m}\Z}$, $m\ge1$ (Theorem~\ref{l15:thm:instancesV2m}).)
The corollary $\chi_{\Vnat}=J$ is treated in Chapter~\ref{ch:L14} and in the Moonshine part of the book (Chapter~\ref{ch:05});
it is not derived from the headline of this chapter.

\section{The remaining files, by role}\label{l15:sec:files}

The $124$ files (the $122$ of the snapshot and the two release files) fall into the following groups. Files not named are helpers of the named ones.
\begingroup\raggedright
\begin{itemize}
\item \textbf{Base objects:} \leanpath{OnePointFunctions}, \leanpath{GradedTrace}, \leanpath{ZhuTheorem} (modular data, $T$-law),
  \leanpath{ChiralBlock} (multilinear $n$-point block interface; not used by the headline), \leanpath{Twisted}.
\item \textbf{Square bracket:} \leanpath{SquareBracketWitt}, \leanpath{SquareBracketVirasoro}, \leanpath{SquareBracketGrading},
  \leanpath{SquareBracketSubst}, \leanpath{SquareBracketC2}, \leanpath{SqWeightSpanning}, \leanpath{SqWeightNat},
  \leanpath{SqModeLinear}, \leanpath{SqModeWeight}.
\item \textbf{Recursion and block conditions:} \leanpath{ZhuRecursion}, \leanpath{ZhuRecursionProof}, \leanpath{ZhuRecursionClose},
  \leanpath{ZhuCoefficientIdentity}, \leanpath{ZhuAssembly}, \leanpath{ZhuNormalization}, \leanpath{ZhuBlockConditions},
  \leanpath{JacobiZhu}, \leanpath{BlockZeroMode}.
\item \textbf{Convergence:} \leanpath{BlockSystem}, \leanpath{BlockGrowth}, \leanpath{OnePointDescent}, \leanpath{OnePointGrowth},
  \leanpath{OnePointODE}, \leanpath{ZhuODEAssembly}, \leanpath{FormalBlockFinite}, \leanpath{UniCollapse} (dead-end certificate),
  \leanpath{TraceGrowthReduction}, \leanpath{C2Corollaries}; the Buhl/Li ordered-spanning route \leanpath{LiFiltration},
  \leanpath{OrderedSpan}, \leanpath{OrderedMonomialBridge}, \leanpath{BuhlSpanning*} (an earlier route to the $v=\lfifteenvac$ growth,
  superseded).
\item \textbf{Blocks and finiteness:} \leanpath{AnalyticBlock}, \leanpath{BlockSpace}, \leanpath{BlockRank}, \leanpath{BlockDeterminacy},
  \leanpath{BlockFiniteness}, \leanpath{BlockFiniteDim}, \leanpath{BlockRelSpanning}, \leanpath{BlockRelSpanE}, \leanpath{BlockRelSpanG},
  \leanpath{BlockRelSpanTransport}, \leanpath{GradedBlock}, \leanpath{SqGradedBlock}, \leanpath{SqGradedBlockSlash},
  \leanpath{FormalBlockFinite}, \leanpath{GNBlock}, \leanpath{BlockFiniteEtingof} (Etingof-route variant).
\item \textbf{Count:} \leanpath{CuspLimit}, \leanpath{CuspPoint}, \leanpath{ZhuCount}, \leanpath{ZhuCountC2}, \leanpath{ZhuBlockCount},
  \leanpath{ZhuAlgebraFinite}, \leanpath{SemisimpleTrace}, \leanpath{SymmetricFunctionals}, \leanpath{TraceSpace},
  \leanpath{WedderburnAMod}, \leanpath{WedderburnSimples}, \leanpath{SimpleAModFinite}, \leanpath{AModBridge}.
\item \textbf{Existence of modules:} \leanpath{StrictSpanning}, the \leanpath{GN*} files (Gaberdiel--Neitzke Prop.~8 in five
  ``links''), \leanpath{RealizationCofinite}, \leanpath{RealizationAssembly}, \leanpath{OmegaCentral}, \leanpath{HL0Scoped},
  \leanpath{AnFiniteDim}, \leanpath{Anm0FiniteDim}, \leanpath{C2Extraction}, \leanpath{C2GradedVanishing}, \leanpath{LowestWeightExistence}.
\item \textbf{Independence:} \leanpath{TraceIndependence}, \leanpath{FibreIndependence}, \leanpath{CharIndependence},
  \leanpath{CharIndependenceAMod}, \leanpath{BottomCharacter}, \leanpath{TraceNonvanish}.
\item \textbf{Assembly:} \leanpath{BlockSpanning}, \leanpath{SlashAmbientObstruction}, \leanpath{SqAmbientRank},
  \leanpath{SqAmbientAssembly}, \leanpath{SqSlashBridge}, \leanpath{MinimalClosure}, \leanpath{FamilyIndexing},
  \leanpath{ConformalWeightChoice}, \leanpath{ZhuModularityFibre}, \leanpath{ZhuModularityODE}, \leanpath{ZhuModularityC2},
  \leanpath{ZhuModularitySurface}, \leanpath{ZhuModularityGN}, \leanpath{ZhuModularityPinned}, \leanpath{GNProp8Closure},
  \leanpath{GNProp8PBW}, \leanpath{StronglyRationalFamily}, \leanpath{ZhuStronglyRational}, \leanpath{ZhuModularityCanonical}
  (the engine, the rational and strongly rational headline forms, the closed trivial instance and the combined Zhu
  correspondence of Chapter~\ref{ch:L13}). Zhu's semisimplicity theorem, which the headline uses, is in
  \leanpath{VOA/\allowbreak Huang/\allowbreak ZhuSemisimple.lean} (Chapter~\ref{ch:L21}).
  The instances \leanpath{ZhuModularityV2mPlus} (for $V_{2m}^+$, the tetrahedral model, $V_{2m}$ and nice VOAs; assigned to this chapter) live in
  \leanpath{VOA/\allowbreak Classification/\allowbreak COne/}, next to their inputs (Chapter~\ref{ch:L22}).
\item \textbf{Examples:} \leanpath{ZhuInstantiation}, \leanpath{TrivialRationalityLattice}, \leanpath{TrivialAdmissiblyRational},
  \leanpath{TrivialStronglyRational}.
\end{itemize}
\endgroup

\section{Referee's checklist}\label{l15:sec:checklist}

\begingroup\small
\begin{longtable}{@{}p{4.3cm}p{4.2cm}p{3.6cm}p{2.3cm}@{}}
\toprule
\textbf{Declaration} & \textbf{Intended meaning} & \textbf{Hypotheses} & \textbf{Check}\\
\midrule\endhead
\lean{zhuModularity_}\allowbreak\lean{canonical_of_}\allowbreak\lean{C2Cofinite_}\allowbreak\lean{admissiblyRational} & \textbf{headline}: Zhu's
modularity theorem for one-point functions of rational $C_2$-cofinite $V$, all irreducibles (canonical family) & $C_2$;
\lean{AdmissiblyRational} (implied by DLM-rationality under $C_2$; converse open); conditional; via theorem \lean{VOA.pbw} & Thm.~\ref{l15:thm:rational}\\
\lean{zhuModularity_}\allowbreak\lean{canonical_of_}\allowbreak\lean{stronglyRational} & \textbf{headline}, strongly rational
form & \lean{StronglyRational}; conditional; via theorem \lean{VOA.pbw} & Thm.~\ref{l15:thm:rational}\\
\lean{isSemisimpleRing_}\allowbreak\lean{A_of_}\allowbreak\lean{admissiblyRational}, \lean{zhuAdmissibleSemisimple_of_Pr} & rational
$\Rightarrow A(V)$ semisimple (Zhu 2.2.3; DLM 4.10(a)) & \lean{Pr} (supplied by \lean{VOA.pbw}), \lean{AdmissiblyRational};
proved & Thm.~\ref{l15:thm:zhuss}\\
\lean{zhuModularity_}\allowbreak\lean{canonical_of_}\allowbreak\lean{C2Cofinite_semisimple} & engine: the same conclusion from
semisimplicity of $A(V)$ & $C_2$; $A(V)$ semisimple; conditional; via theorem \lean{VOA.pbw} & Thm.~\ref{l15:thm:canonical}\\
\lean{finite_Φ_of_C2Cofinite_pbw} & $\Phi(V)$ finite & $C_2$ & Thm.~\ref{l15:thm:finitephi}\\
\lean{zhuModularity_V2mPlus}, \lean{zhuModularity_cubicModel} & first nontrivial instances ($V_{2m}^+$, $m\ge2$; tetrahedral) &
none; \textbf{unconditional}; via theorem \lean{VOA.pbw} & Thm.~\ref{l15:thm:instances}\\
\lean{zhuModularity_V2m} & Zhu's theorem for every lattice VOA $V_{\sqrt{2m}\Z}$, $m\ge1$ & none; \textbf{unconditional}; via
theorem \lean{VOA.pbw} & Thm.~\ref{l15:thm:instancesV2m}\\
\lean{zhuModularity_of_nice} & Zhu's theorem for every nice VOA & \lean{Nice V} (only strong rationality used); conditional &
Thm.~\ref{l15:thm:instancesV2m}\\
\lean{trivialVOA_}\allowbreak\lean{zhuModularity_closed} & instance at $\C\lfifteenvac$ & none; unconditional & Thm.~\ref{l15:thm:trivial}\\
\lean{zhuModularity_}\allowbreak\lean{pinned_of_}\allowbreak\lean{C2Cofinite} & Zhu's modularity theorem, pinned to the given family & $C_2$; $A(V)$ ss; $\Phi$ finite;
\lean{Pr}; \lean{hcard}, \lean{hcw}, \lean{hsep}; conditional & Thm.~\ref{l15:thm:headline}\\
\lean{zhuModularity_}\allowbreak\lean{pinned_of_}\allowbreak\lean{C2Cofinite_voaPBW} & same, \lean{Pr} from the theorem \lean{VOA.pbw} & as above minus \lean{Pr}; via theorem \lean{VOA.pbw} &
Thm.~\ref{l15:thm:headline}\\
\lean{zhuModularity_}\allowbreak\lean{stronglyRational_voaPBW} & same for strongly rational $V$, canonical family & \lean{StronglyRational}, DLM finiteness,
Zhu semisimplicity (both carried, both now theorems for strongly rational $V$); via theorem \lean{VOA.pbw}; superseded & Thm.~\ref{l15:thm:sr}\\
\lean{trivialVOA_zhuModularity} & non-vacuity at $\C\lfifteenvac$ & $\#\Phi=1$ (a theorem); via theorem \lean{VOA.pbw} & Thm.~\ref{l15:thm:trivial}\\
\lean{onePointSubexp_of_c2} & subexponential growth of all one-point coefficients & $C_2$, conformal weight & Thm.~\ref{l15:thm:growth}\\
\lean{qSummableAt_}\allowbreak\lean{of_subexponential} & growth $\Rightarrow$ convergence on $\lfifteenH$ & subexponential & Thm.~\ref{l15:thm:growth}\\
\lean{zhuRecursionPinned_holds} & Zhu's recursion & none & Thm.~\ref{l15:thm:recursion}\\
\lean{torusOnePointHom_}\allowbreak\lean{isFormalBlock} & traces satisfy the block conditions & none & Thm.~\ref{l15:thm:formalblock}\\
\lean{torusOnePoint_}\allowbreak\lean{isAnalyticBlock} & traces are analytic blocks & support, growth & Thm.~\ref{l15:thm:traceblock}\\
\lean{exists_finrank_}\allowbreak\lean{analyticBlocks_le} & $\dim\lfifteenB<\infty$ (Zhu 5.3.2) & $C_2$ & Thm.~\ref{l15:thm:532}\\
\lean{finiteDimensional_}\allowbreak\lean{A_of_C2Cofinite} & $\dim A(V)<\infty$ & $C_2$ & Thm.~\ref{l15:thm:Afin}\\
\lean{finiteDimensional_}\allowbreak\lean{and_finrank_}\allowbreak\lean{analyticBlocks_}\allowbreak\lean{le_card_Φ} & $\dim\lfifteenB<\infty$ and $\dim\lfifteenB\le\#\Phi$ (Zhu 5.3.3, \textbf{upper bound only}) & $C_2$, ss, $\Phi$ finite,
\lean{hreal}, \lean{Pr} & Thm.~\ref{l15:thm:533}\\
\lean{gnProp8_GNReach}, \lean{cofinite_}\allowbreak\lean{of_strictSpanning} & strict spanning in VA modes (cf.\ GN Prop.~8; convention unverified), GN Thm.~11 & $C_2$ & Thm.~\ref{l15:thm:gn}\\
\lean{simpleAModsRealized_}\allowbreak\lean{of_C2Cofinite} & every simple $A(V)$-module is the bottom of an irreducible \emph{ordinary} module & $C_2$, \lean{Pr} (explicit) & Thm.~\ref{l15:thm:realized}\\
\lean{traceBlock_}\allowbreak\lean{S_sqRestrict_}\allowbreak\lean{mem_analyticBlocks} & $S$-slash of a trace is a block & $C_2$, conformal weight & Thm.~\ref{l15:thm:slashblock}\\
\lean{fibreIndependent_}\allowbreak\lean{of_bottomSeparation} & traces independent & growth, simple separated bottoms & Thm.~\ref{l15:thm:indep}\\
\lean{zhu_modularity_}\allowbreak\lean{of_spanning} & $\rho$ from spanning + independence & the two data & Thm.~\ref{l15:thm:spanning}\\
\lean{slashSpannedAtS_}\allowbreak\lean{of_sqRank_le} & $S$-spanning from rank bound & rank, independence & Thm.~\ref{l15:thm:rank}\\
\lean{char_T_modularity}, \lean{zhu_modularity_}\allowbreak\lean{T_onePoint} & law at $T$ & conformal weight & Thm.~\ref{l15:thm:T}\\
\lean{Lsq_virasoro_apply'}, \lean{sqWeightSpace_isInternal} & torus Virasoro, $L[0]$-grading & none & Thm.~\ref{l15:thm:sqvir}\\
\lean{not_rank_}\allowbreak\lean{analyticSlashAmbient_le} & legacy rank hypothesis is contradictory & displayed & Sec.~\ref{l15:sec:slash}\\
\lean{Conjecture_5_3_10} & orbifold modularity (statement) & --- ; \textbf{mis-normalised} & Sec.~\ref{l15:sec:twisted}\\
\bottomrule
\end{longtable}
\endgroup

\noindent All entries were checked with \texttt{\#print axioms} under \texttt{lake env lean -DmaxSynthPendingDepth=3}: the
standard three axioms for every entry (at the snapshot, \lean{VOA.pbw} was added for the three entries marked ``via theorem
\lean{VOA.pbw}'' and for the two wrappers of the separate Lean check below; re-run at a later commit, after \lean{VOA.pbw} became a theorem (not re-run for this edition), for the changed entries, the \lean{\_voaPBW} items and those of the check: standard three; the
unchanged entries were not re-run). The separate Lean check of this chapter against the library (not part of this release) compiled with the same
flag; its theorems \lean{vacuumGradeZero_always}, \lean{hsimp_of_semisimple}, \lean{zhuModularity_}\allowbreak\lean{pinned_lean},
\lean{headline_}\allowbreak\lean{traces_are_genuine}, \lean{twistedTraceModular_}\allowbreak\lean{of_weight_bot}, \lean{twistedTraceModular_}\allowbreak\lean{of_c_off_lattice}
use the standard three axioms; \lean{trivialVOA_}\allowbreak\lean{zhuModularity_uncond} and \lean{stronglyRational_}\allowbreak\lean{wrapper_fires} added
\lean{VOA.pbw} at the snapshot and report the standard three in that later run.

\paragraph{Release audit 2026-10-27: RELEASE\_AUDIT\_ZHU.md.} The independent release audit
(\texttt{Notes/\allowbreak Textbook/\allowbreak referee/\allowbreak RELEASE\_AUDIT\_ZHU.md}) re-ran \texttt{\#print axioms} on the
chapter's headlines (all standard) and required the corrections applied in this revision: the hypothesis is semisimplicity
of $A(V)$, not rationality (Remark~\ref{l15:rem:hyp}); one-point functions only; convergence is a separate theorem; the block
count is an upper bound; \lean{[Finite (Φ W)]} is a theorem; the mode convention of strict spanning is unverified against
Gaberdiel--Neitzke; \lean{simpleAModsRealized_}\allowbreak\lean{of_C2Cofinite} takes \lean{Pr} and $C_2$. The theorems it
recommended were then added to the library and are presented here. For this revision a scratch file was run at
an earlier commit (not re-run for this edition) with \texttt{lake env lean -DmaxSynthPendingDepth=3} from the repository root, against the oleans built
2026-10-25 for the imported modules; the two new files had no oleans, so their text was included verbatim in the scratch
file and elaborated there. Output:
\begin{plaincode}
'VOA.zhuModularity_canonical_of_C2Cofinite_semisimple' depends on axioms: [propext, Classical.choice, Quot.sound]
'VOA.trivialVOA_zhuModularity_closed' depends on axioms: [propext, Classical.choice, Quot.sound]
'VOA.zhuCorrespondence_of_C2Cofinite' depends on axioms: [propext, Classical.choice, Quot.sound]
'VOA.Classification.COne.zhuModularity_V2mPlus' depends on axioms: [propext, Classical.choice, Quot.sound]
'VOA.Classification.COne.zhuModularity_cubicModel' depends on axioms: [propext, Classical.choice, Quot.sound]
'VOA.Huang.ZhuFiniteIrr.finite_Φ_of_C2Cofinite_pbw' depends on axioms: [propext, Classical.choice, Quot.sound]
'VOA.Classification.COne.V2mPlus_zhuSemisimple' depends on axioms: [propext, Classical.choice, Quot.sound]
'VOA.Classification.COne.V2mPlus_c2Cofinite' depends on axioms: [propext, Classical.choice, Quot.sound]
'VOA.Classification.COne.V2mPlus_zhu_finrank' depends on axioms: [propext, Classical.choice, Quot.sound]
'VOA.Classification.COne.OrbTop.cubicModel_zhuSemisimple' depends on axioms: [propext, Classical.choice, Quot.sound]
'VOA.Classification.COne.OrbA4.cubicModel_c2Cofinite_holds' depends on axioms: [propext, Classical.choice, Quot.sound]
'VOA.finiteDimensional_and_finrank_analyticBlocks_le_card_Φ' depends on axioms: [propext, Classical.choice, Quot.sound]
'VOA.finrank_analyticBlocks_le_card_Φ' depends on axioms: [propext, Classical.choice, Quot.sound]
'VOA.onePointSubexp_of_c2' depends on axioms: [propext, Classical.choice, Quot.sound]
'VOA.qSummableAt_of_subexponential' depends on axioms: [propext, Classical.choice, Quot.sound]
'VOA.simpleAModsRealized_of_C2Cofinite' depends on axioms: [propext, Classical.choice, Quot.sound]
'VOA.nonempty_strictSpanning_of_C2Cofinite' depends on axioms: [propext, Classical.choice, Quot.sound]
'VOA.zhuBijective_of_voaPBW' depends on axioms: [propext, Classical.choice, Quot.sound]
\end{plaincode}
No \texttt{lake build} was run for this revision. The release plan records that the umbrellas were green when these
theorems were merged (2026-10-28 entry); the release audit's own caveat stands: nothing counts as released until an aggregate
\texttt{lake build} containing the modularity umbrella is green at the release commit.

\paragraph{Second release revision 2026-10-28: RELEASE\_AUDIT\_NEW\_DECLS.md, finding F2.} The audit of the new declarations
(\texttt{Notes/\allowbreak Textbook/\allowbreak referee/\allowbreak RELEASE\_AUDIT\_NEW\_DECLS.md}) found that the module
docstrings of \leanpath{ZhuModularityCanonical.lean} and \leanpath{ZhuModularityV2mPlus.lean} understated the library: the
implication ``rational $\Rightarrow A(V)$ semisimple'' was already proved
(\lean{VOA.Huang.ZhuSemisimple.}\allowbreak\lean{isSemisimpleRing_}\allowbreak\lean{A_of_}\allowbreak\lean{admissiblyRational}), and the $V_{2m}$ instance followed in two lines. The
first release revision of this chapter repeated both errors. The library then added
\lean{zhuModularity_}\allowbreak\lean{canonical_of_}\allowbreak\lean{C2Cofinite_}\allowbreak\lean{admissiblyRational}, \lean{zhuModularity_}\allowbreak\lean{canonical_of_}\allowbreak\lean{stronglyRational},
\lean{zhuModularity_V2m} and \lean{zhuModularity_of_nice}, and corrected the two docstrings. This revision presents them
(Sections~\ref{l15:sec:rational} and~\ref{l15:sec:instances}) and corrects the chapter. For it, a scratch file was run at
an earlier commit (not re-run for this edition) with \texttt{lake env lean -DmaxSynthPendingDepth=3} from the repository root. The oleans of
\leanpath{ZhuModularityCanonical.lean} and \leanpath{ZhuModularityV2mPlus.lean} (2026-10-28 00:02) predate their current sources
(04:27), so the scratch file imported the two files' imports and contained their text verbatim; all other modules were taken
from existing oleans and nothing was built. Output (the run exited with status 0):
\begin{plaincode}
'VOA.zhuModularity_canonical_of_C2Cofinite_admissiblyRational' depends on axioms: [propext,
 Classical.choice,
 Quot.sound]
'VOA.zhuModularity_canonical_of_stronglyRational' depends on axioms: [propext, Classical.choice, Quot.sound]
'VOA.zhuModularity_canonical_of_C2Cofinite_semisimple' depends on axioms: [propext, Classical.choice, Quot.sound]
'VOA.Classification.COne.zhuModularity_V2m' depends on axioms: [propext, Classical.choice, Quot.sound]
'VOA.Classification.COne.zhuModularity_of_nice' depends on axioms: [propext, Classical.choice, Quot.sound]
'VOA.Classification.COne.zhuModularity_V2mPlus' depends on axioms: [propext, Classical.choice, Quot.sound]
'VOA.Classification.COne.zhuModularity_cubicModel' depends on axioms: [propext, Classical.choice, Quot.sound]
'VOA.Huang.ZhuSemisimple.isSemisimpleRing_A_of_admissiblyRational' depends on axioms: [propext,
 Classical.choice,
 Quot.sound]
'VOA.Huang.ZhuSemisimple.zhuAdmissibleSemisimple_of_Pr' depends on axioms: [propext, Classical.choice, Quot.sound]
'VOA.Classification.COne.V2m_stronglyRational' depends on axioms: [propext, Classical.choice, Quot.sound]
'VOA.Classification.COne.V2m_nice' depends on axioms: [propext, Classical.choice, Quot.sound]
'VOA.Classification.COne.V2mPlus_nice' depends on axioms: [propext, Classical.choice, Quot.sound]
'VOA.VHat.DLM3.voaPBWBottomProjection' depends on axioms: [propext, Classical.choice, Quot.sound]
\end{plaincode}
The same file also checked the \lean{example} of Theorem~\ref{l15:thm:zhuss} (\lean{ZhuAdmissibleSemisimple W} for every $V$).
No \texttt{lake build} was run for this revision. \texttt{Notes/RELEASE\_PLAN.md} records that the umbrella builds of
\lean{VOA.Modularity} and \lean{VOA.Classification} containing the new theorems were green (12,271 jobs) with standard axioms.
The audit's process finding F1 stands: nothing counts as released until the aggregate \texttt{lake build} is green at the
release commit.

\section{Notes and references}\label{l15:sec:notes}

The theorem is due to Yongchang Zhu, \emph{Modular invariance of characters of vertex operator algebras}, J.\ Amer.\ Math.\
Soc.\ \textbf{9} (1996), 237--302; the square-bracket structure, the recursion (Prop.~4.3.2), the block space and
Theorems~5.3.1--5.3.3 are from \S\S4--5 there. The orbifold and twisted generalisation, the systematic treatment of
convergence and the symmetric-functional lemma (Lemma~10.7) are in C.~Dong, H.~Li, G.~Mason, \emph{Modular-invariance of
trace functions in orbifold theory and generalized Moonshine}, Comm.\ Math.\ Phys.\ \textbf{214} (2000), 1--56
(arXiv:q-alg/9703016); finiteness of irreducible admissible modules is C.~Dong, H.~Li, G.~Mason, \emph{Twisted
representations of vertex operator algebras}, Math.\ Ann.\ \textbf{310} (1998), 571--600. The extension to non-rational
$C_2$-cofinite algebras by pseudo-traces is M.~Miyamoto, \emph{Modular invariance of vertex operator algebras satisfying
$C_2$-cofiniteness}, Duke Math.\ J.\ \textbf{122} (2004), 51--91 (arXiv:math/0209101). The spanning theorem used for
existence is M.~R.~Gaberdiel and A.~Neitzke, \emph{Rationality, quasirationality and finite W-algebras}, Comm.\ Math.\
Phys.\ \textbf{238} (2003), 305--331 (arXiv:hep-th/0009235), Prop.~8 and Thm.~11; related spanning sets are in G.~Buhl,
\emph{A spanning set for VOA modules}, J.\ Algebra \textbf{254} (2002), 125--151, and H.~Li, \emph{Abelianizing vertex
algebras}, Comm.\ Math.\ Phys.\ \textbf{259} (2005), 391--411. The $C_2$-implies-finite-$A(V)$ and related results are in T.~Abe,
G.~Buhl, C.~Dong, \emph{Rationality, regularity, and $C_2$-cofiniteness}, Trans.\ Amer.\ Math.\ Soc.\ \textbf{356} (2004),
3391--3402 (arXiv:math/0204021). The bimodules $A_{n,m}(V)$ and the characterisation of rationality by the higher Zhu
algebras are C.~Dong and C.~Jiang, \emph{Bimodules associated to vertex operator algebras}, Math.\ Z.\ \textbf{259} (2008),
799--826. The Verlinde formula for rational $C_2$-cofinite VOAs is Y.-Z.~Huang, \emph{Vertex operator algebras and the
Verlinde conjecture}, Commun.\ Contemp.\ Math.\ \textbf{10} (2008), 103--154 (arXiv:math/0406291, 2004--2005), and the
congruence property is C.~Dong, X.~Lin, S.-H.~Ng, \emph{Congruence property in conformal field theory}, Algebra Number Theory
\textbf{9} (2015), 2121--2166. The textbook account followed by the library's docstrings is T.~Gannon, \emph{Moonshine
beyond the Monster}, Cambridge University Press (2006), \S5.3 (Theorem~5.3.8 there is Zhu's theorem). Modular linear
differential equations in physics: S.~D.~Mathur, S.~Mukhi, A.~Sen, \emph{On the classification of rational conformal field
theories}, Phys.\ Lett.\ B \textbf{213} (1988), 303--308; the Verlinde formula: E.~Verlinde, \emph{Fusion rules and modular
transformations in 2D conformal field theory}, Nucl.\ Phys.\ B \textbf{300} (1988), 360--376. The orbifold boson and its modular data:
R.~Dijkgraaf, C.~Vafa, E.~Verlinde, H.~Verlinde, \emph{The operator algebra of orbifold models}, Comm.\ Math.\ Phys.\
\textbf{123} (1989), 485--526; the Zhu algebra of $V_L^+$ and its semisimplicity: T.~Abe, \emph{Rationality of the vertex
operator algebra $V_L^+$ for a positive definite even lattice $L$}, Math.\ Z.\ \textbf{249} (2005), 455--484 (Theorem~7.1);
$C_2$-cofiniteness of $V_L^+$: G.~Yamskulna, \emph{$C_2$-cofiniteness of the vertex operator algebra $V_L^+$ when $L$ is a
rank one lattice}, Comm.\ Algebra \textbf{32} (2004), 927--954. The quasi-lisse extension
(Chapter~\ref{ch:L16}) is T.~Arakawa and K.~Kawasetsu, \emph{Quasi-lisse vertex algebras and modular linear differential
equations}, in \emph{Lie groups, geometry, and representation theory}, Progr.\ Math.\ \textbf{326} (2018), 41--57
(arXiv:1610.05865).

\endgroup
